\documentclass[11pt]{amsart}
\usepackage{amsmath,amssymb,amsthm,mathtools,mathrsfs}
\usepackage[margin=1.08in]{geometry}
\usepackage{booktabs}
\usepackage{array}
\usepackage{longtable}
\usepackage{tabularx}
\usepackage{hyperref}
\hypersetup{hidelinks}
\usepackage[final,expansion=false]{microtype}
\setlength{\parskip}{4pt plus 1pt}
\emergencystretch=2.5em

\newtheorem{theorem}{Theorem}[section]
\newtheorem{proposition}[theorem]{Proposition}
\newtheorem{lemma}[theorem]{Lemma}
\newtheorem{corollary}[theorem]{Corollary}
\theoremstyle{definition}
\newtheorem{definition}[theorem]{Definition}
\newtheorem{convention}[theorem]{Convention}
\theoremstyle{remark}
\newtheorem{remark}[theorem]{Remark}

\newcommand{\C}{\mathbb C}
\newcommand{\R}{\mathbb R}
\newcommand{\N}{\mathbb N}
\newcommand{\Z}{\mathbb Z}
\newcommand{\Q}{\mathbb Q}
\newcommand{\A}{\mathbb A}
\newcommand{\dd}{\,d}
\newcommand{\dstar}{\,d^{*}}
\newcommand{\Tr}{\operatorname{Tr}}
\newcommand{\Rel}{\operatorname{Re}}
\newcommand{\Pic}{\operatorname{Pic}}
\newcommand{\Det}{\operatorname{Det}}
\newcommand{\Picb}{\operatorname{Pic}_{b}}
\newcommand{\Hom}{\operatorname{Hom}}
\newcommand{\Cone}{\operatorname{Cone}}
\newcommand{\supp}{\operatorname{supp}}
\newcommand{\vol}{\operatorname{vol}}
\newcommand{\sgn}{\operatorname{sgn}}
\newcommand{\id}{\operatorname{id}}
\newcommand{\HH}{\operatorname{HH}}
\newcommand{\Sh}{\operatorname{Sh}}
\newcommand{\Perf}{\operatorname{Perf}}
\newcommand{\Aut}{\operatorname{Aut}}
\newcommand{\Res}{\operatorname{Res}}
\newcommand{\Inf}{\operatorname{Inf}}
\newcommand{\Rest}{\operatorname{Res}}
\newcommand{\Eff}{\operatorname{Eff}}
\newcommand{\Th}{\widehat{T}}
\newcommand{\Wh}{\widehat{W}}
\newcommand{\Ch}{\widehat{C}_S}
\newcommand{\DA}{D_A}
\newcommand{\Atr}{\mathcal A_{\mathrm{tr}}}
\newcommand{\Corrtr}{\mathrm{Corr}^{\mathrm{tr}}_{A}}
\newcommand{\Nrad}{\mathcal N}
\newcommand{\BT}{B_T}
\newcommand{\NAd}{N_A^{\mathrm{dist}}}
\newcommand{\Log}{\operatorname{Log}}
\newcommand{\cond}{\operatorname{cond}}
\newcolumntype{Y}{>{\raggedright\arraybackslash}X}

\begin{document}

\title[A correspondence algebra on the idele class group]
{A Correspondence Algebra with a Canonical Trace\\
on the Idele Class Group,\\
its Tate Symbol and its Numerical Form}

\author{David Alejandro Trejo Pizzo}
\thanks{Independent Researcher. \texttt{dtrejopizzo@gmail.com}}

\date{\today}

\begin{abstract}
Let $S$ be a finite set of places of $\Q$ containing the real place and at
least one prime, and let $C_S$ be the semilocal idele class group.  We study
the graph bimodules of idele classes over the crossed product of a function
algebra on the semilocal adelic space by the positive $S$-units, with their
composition and involution, and show that the Hochschild complex of such a
bimodule does not detect the class: for a class whose real coordinate is not
an $S$-unit, each lattice component of the complex is homotopy equivalent to
an exterior algebra on $|S_f|$ generators, the local factor at a finite place
being one class in degree zero.  On the convolution algebra $C_c^{\infty}(C_S)$ we
define a trace as a sum of local fixed-point distributions and compute its
Fourier symbol, a real-analytic function whose continuation on the principal
component has simple poles at $\pm i/2$; the residue terms there are the two
Tate characters.  The Hermitian form attached to
the trace is nondegenerate, and in the normalization with zero constant term
it has infinite index of both signs.  The two Tate characters give the
bidegree; their common kernel is the primitive subspace.  The two ruling
classes are not test functions.  They are limits, in a completion, of test
functions along which the scalar form tends to zero.  In the limit over $S$
the value of the form on a pair of slowly varying radial approximations of
the two rulings tends to $1$, the value of the hyperbolic form, by Mertens'
theorem.
\end{abstract}

\maketitle

\tableofcontents

\section{Introduction}\label{sec:intro}

On a smooth projective curve $C$ over a finite field, the correspondences on
$C\times C$ form a ring with a transpose and an intersection pairing
\cite{Weil1948}.  The two rulings $C\times\{pt\}$ and $\{pt\}\times C$ span a
hyperbolic plane in the numerical group, and every correspondence is the sum
of a ruling part, recorded by its bidegree, and a primitive part.

In this paper we construct the same operations on an adelic space.  The
space is the quotient of the semilocal adelic space $A_S$ by the positive
$S$-units, the space on which Connes based his trace formula
\cite{Connes99}; see also \cite{MeyerIdele}.  The correspondences are the
graph bimodules of idele classes.  The intersection pairing is an analytic
trace, a sum of local distributions with density $|x|_v^{1/2}/|1-x|_v$, the
half-density form of the fixed-point contribution of a dilation.  Outside
Appendix~\ref{app:mertens} the only arithmetic input is the product formula.  No explicit formula and no
positivity statement is used.

The theory differs from the theory for curves in three respects.  First,
$C_S$ is a continuum, so there are infinitely many pairwise non-isomorphic
generators, and the pairing is a Hermitian form on a nuclear space, not a
bilinear form on a finitely generated group.  Second, the Hochschild complex
of a graph bimodule does not see the dilation.  At a finite place the fixed
point of a dilation is the origin of $\Q_p$, and the twisted Hochschild
complex of the local function algebra is homotopy equivalent to one class in
degree zero for every dilation other than the identity; at the real place it
is contractible for every dilation other than the identity
(Theorem~\ref{thm:shadow}, Corollary~\ref{cor:blind}).  Assembling the places
uses a twisted form of Wodzicki's excision theorem \cite{Wodzicki}, which we
prove (Theorem~\ref{thm:excision-input}).  The dependence on the dilation is
carried by the trace.  Third, the two ruling directions are not represented
by elements of the test algebra.  They are eigenvectors of the whole
convolution algebra whose eigenvalues are the values at $\mp i/2$ of the
Fourier transform of the acting element, off the unitary axis, and no nonzero
finite-dimensional subspace of the test algebra is stable under convolution
(Theorem~\ref{thm:nosplit}).  The rulings are limits of test functions whose
mass moves to infinity (Theorem~\ref{thm:rulings}).

\subsection*{Main results}

Fix a finite set $S$ of places of $\Q$ with $\infty\in S$ and $S_f\ne
\emptyset$.  Write $A_S=\R\times\prod_{p\in S_f}\Q_p$, and let $\Gamma_S$ be
the group of positive rationals whose prime factorization uses only the
primes of $S_f$.

\begin{theorem}[Lemma~\ref{thm:split} and
Corollary~\ref{cor:trivial}]\label{thm:introA}
The map $\varphi(x)=\prod_{p\in S_f}p^{v_p(x_p)}$ is a continuous
homomorphism $A_S^{\times}\to\Gamma_S$ restricting to the identity on
$\Gamma_S$.  Hence, as topological groups,
\[
A_S^{\times}=\Gamma_S\times\ker\varphi,\qquad
C_S:=\Gamma_S\backslash A_S^{\times}\cong\ker\varphi=
\R^{\times}\times\prod_{p\in S_f}\Z_p^{\times},
\]
and $|x|_S=|x_{\infty}|_{\infty}$ on $\ker\varphi$.  The class map therefore
has a section $\sigma$ which is a homomorphism.  For $F$ in
$D_A:=C_c^{\infty}(C_S)$ let $E_F$ be the smeared object
$u\mapsto F(u)\cdot1\in E_{\sigma(u)}$.  Then
$E_F\star E_G=E_{F\star G}$ and $E_F^{\dagger}=E_{F^{\dagger}}$, where
$\star$ on the right is convolution and
$F^{\dagger}(u)=\overline{F(u^{-1})}$, and $F\mapsto E_F$ is a bijection of
$D_A$ onto the smeared objects.
\end{theorem}

For a place $v$ of $\Q$ and $x\in\Q_v^{\times}$, $x\ne1$, put
\[
\tau_v(x):=\frac{|x|_v^{1/2}}{|1-x|_v}.
\]
At a finite place $\tau_p(p^ku)=p^{-|k|/2}$ for $k\ne0$, and
$\tau_p(u)=|1-u|_p^{-1}$ for a unit $u\ne1$.  For $u\in1+q_p\Z_p$ ($q_p=p$
for odd $p$, $q_2=4$) the second value is $N(u)/(1-p^{-1})$, where $N(u)$ is
the index in $\Z_p^{\times}$ of the closed subgroup generated by $u$, and
also the dimension of the quotient of the functions on a valuation shell by
the closure of the image of $1-\sigma$, $\sigma$ the dilation by $u$
(Lemma~\ref{lem:tau}, Proposition~\ref{prop:haarrank}).  The trace is
\[
T(F)=\sum_{v\in S}W_v(F)+c\,F(1),
\]
where $W_p$ is the integral of $F$ against $\tau_p$ over the shells of
nonzero valuation plus a regularized integral over the units, $W_{\infty}$
is the archimedean distribution with its regularization at the identity, and
$c$ is a real constant (Definition~\ref{def:trace}).  On test functions
vanishing near the identity class, $T$ is a finite sum of absolutely
convergent integrals and does not depend on $c$
(Theorem~\ref{thm:absconv}).  A continuous functional with the same values
there differs from $T$ by a finite combination of derivatives of $F$ at the
identity class (Theorem~\ref{thm:unique}).

The dual group $\widehat C_S$ is a countable disjoint union of copies of
$\R$, indexed by a parity $\varepsilon\in\{0,1\}$ and a system $(\kappa_q)$
of characters of the groups $\Z_q^{\times}$, $q\in S_f$.  We write
$\cond(\kappa)$ for the conductor exponent of $\kappa$.

\begin{theorem}[Theorems~\ref{thm:symbol} and \ref{thm:ML}]\label{thm:introB}
For every $H\in D_A$,
\[
T(H)=\int_{\widehat C_S}\widehat H(\chi)\bigl(\widehat T(\chi)+c\bigr)\,d\chi,
\qquad
\widehat T=\widehat W_{\infty}+\sum_{p\in S_f}\widehat W_p ,
\]
the integral being absolutely convergent.  The function $\widehat T$ is
real-valued, real-analytic on each component, bounded above on
$\widehat C_S$, and tends to $-\infty$ as $|\tau|\to\infty$ on each
component.  At a finite place
\[
\widehat W_p(\chi)=
\begin{cases}
2(\log p)\,\Rel\bigl(\zeta_p/(1-\zeta_p)\bigr), & \kappa_p=1,\\[2pt]
-\cond(\kappa_p)\log p, & \kappa_p\neq1,
\end{cases}
\qquad \zeta_p=p^{-1/2}\overline{\eta_p(\chi)},
\]
with $\eta_p(\chi)=p^{-i\tau}\prod_{q\ne p}\kappa_q(p)^{-1}$, and on the even
components $\widehat W_{\infty}=\log\pi-\Rel\,\psi(\tfrac14+\tfrac{i\tau}2)$.
For $\kappa_p=1$ the function $\widehat W_p$ extends to a meromorphic
$2\pi/\log p$-periodic function of $\tau$ whose poles are simple, lie on the
lines $\Im\tau=\pm\tfrac12$, and have residues $\mp i$.  The even
archimedean symbol is meromorphic with simple poles at
$\pm i(2n+\tfrac12)$, $n\ge0$, and the same residues.  On the principal
component every place of $S$ contributes one pole at each of $\pm i/2$, so
the residue of $\widehat T$ there is $\mp i(|S_f|+1)$.
\end{theorem}

The next theorem uses only that $\widehat T+c$ is real-analytic, not
identically zero on any component, and of at most logarithmic growth.

\begin{theorem}[Theorem~\ref{thm:radical}]\label{thm:introC}
Let $I_A(F,G):=T(F\star G^{\dagger})$; we call $I_A$ the scalar form.  If $I_A(F,G)=0$ for all $G\in D_A$,
then $F=0$, for every value of $c$.  Thus the Hermitian form $I_A$ is
nondegenerate on the test algebra.
\end{theorem}

By the product formula the modulus $|\cdot|_S$ is well defined on $C_S$, so
$\lambda_{\mp}(F)=\int_{C_S}F(u)|u|_S^{\mp1/2}\dstar u$ are characters of the
convolution algebra.  They are the values $\widehat F(\mp i/2)$ of the
entire continuation of $\widehat F$ on the principal component, and they are
the residue terms produced by the poles of the archimedean symbol at
$\pm i/2$ when the contour is moved off the real axis
(Theorem~\ref{thm:ML}(e); for the finite places see part (d)).  Put
$M=(\lambda_-,\lambda_+)$, $B_T=\C^{2}$ and
$\beta_T(z,w)=z_-\overline{w_+}+z_+\overline{w_-}$.  There is an exact
sequence of $D_A$-modules
\[
0\longrightarrow D_A^{\mathrm{prim}}\longrightarrow D_A
\stackrel{M}{\longrightarrow}B_T\longrightarrow 0,
\]
in which the involution acts on $B_T$ by the conjugate-linear exchange of the
two coordinates, and $I_A$ is nondegenerate on the kernel
(Theorem~\ref{thm:tateseq}).  We call $M(F)$ the bidegree of $F$.  The
trace is not a function of $M$: for $u,u'\in1+p\Z_p$ with
$v_p(1-u)=j\ne j'=v_p(1-u')$ there are test functions $F,F'$, supported over
the cosets $u(1+p^m\Z_p)$ and $u'(1+p^m\Z_p)$, with $M(F)=M(F')$ and
$T(F)-T(F')=\nu(\log p)(p^{j}-p^{j'})$, $\nu>0$
(Theorem~\ref{thm:separate}).  The sequence has no $D_A$-module section
(Theorem~\ref{thm:nosplit}).

Let $D_A'$ be the space of continuous conjugate-linear functionals on $D_A$.
The map $J(F)=(I_A(F,\cdot),M(F))$ embeds $D_A$ in $D_A'\oplus B_T$ with
dense image (Theorem~\ref{thm:completion}).  In this space the rulings
appear.  For the sequences
$F_n(w)=n^{-1}|w|_S^{-1/2}\vartheta(n^{-1}\log|w|_S)$ and
$G_n=F_n^{\dagger}$, with $\vartheta\ge0$ a bump function on $(0,\infty)$ of
integral one, $J(F_n)\to(0,e_+)$ and $J(G_n)\to(0,e_-)$, where
$e_-=(1,0)$ and $e_+=(0,1)$, and the limits are
eigenvectors of $D_A$ with eigenvalues $\lambda_{\pm}$
(Theorem~\ref{thm:rulings}).  The plane $0\oplus B_T$ is the set of elements
whose scalar pairing with every test function vanishes
(Theorem~\ref{thm:rulnull}).  More generally $I_A(F_\alpha,G_\beta)\to0$,
uniformly in $\beta$, whenever $M(F_\alpha)\to e_+$,
$\lVert F_\alpha\rVert_{L^1}\to0$, and $(G_\beta)$ approximates $e_-$ in the
same sense with the seminorm $N$ of \eqref{eq:Nsemi} bounded on it, while
$\beta_T(MF_\alpha,MG_\beta)\to1$ (Theorem~\ref{thm:rulnullclass}).

\begin{theorem}[Theorem~\ref{thm:signature}]\label{thm:introD}
For $F\in D_A$,
$I_A(F,F)=\int_{\widehat C_S}|\widehat F(\chi)|^{2}(\widehat T(\chi)+c)\,
d\chi$ is real.  For every $c$ the form has infinite negative index.  With
$m=\sup\widehat T$, it has infinite positive index if $c>-m$ and is negative
definite if $c\le-m$; and $m>\log\pi+\gamma+3\log2+\pi/2>0$.  In particular
for $c=0$ the form has infinite index of both signs.  The same statements
hold on $D_A^{\mathrm{prim}}$.
\end{theorem}

The remaining sections treat the determinant and the change of $S$.  The
norm $\lVert F\rVert_{\mathrm{tr}}=\Lambda(F)+N(F)$, where $\Lambda$ weights
$F$ by $|w|_S^{1/2}+|w|_S^{-1/2}$ and $N$ weights $\widehat F$ by
$|\widehat T|+|c|$, defines a commutative Banach $\dagger$-algebra $\Atr$ on
which $T$ and both Tate characters are continuous.  On it the series
$\log\Delta_F(z)=\sum_{r\ge1}(-1)^{r+1}z^{r}T(F^{\star r})/r$ converges for
$|z|<1/\lVert F\rVert_{L^1}$ and equals
$\int\Log(1+z\widehat F)(\widehat T+c)\,d\chi$, which gives the
multiplicativity of $\Delta$ under $(F,G)\mapsto F+G+zF\star G$ in the form
stated in Theorem~\ref{thm:banach}(f).  The Riemann--Roch pairing is the
pair
\[
\lambda_{RR}(F,G)=\bigl(e^{I_A(F,G)},\,e^{\beta_T(MF,MG)}\bigr),
\]
whose logarithm $(I_A,\beta_T\circ(M\times M))$ has the shape of Weil's
decomposition $D\cdot D'=P\cdot P'+d_1d_2'+d_2d_1'$ \cite{Weil1948}.  The
analogy is formal: the first component is the full $I_A(F,G)$, not a pairing
of primitive parts (Theorem~\ref{thm:RR}).  Along such approximations
of the rulings the first component tends to $1$ and the second to $e$
(Theorem~\ref{thm:scalartrivial}).

The vanishing of the scalar form on the rulings belongs to a fixed finite
$S$.  The sets $S$ form a directed system; inflation along it changes the
trace of a given test function only at finitely many new primes, so the
trace stabilizes, and the stable value is invariant under inflation
(Theorems~\ref{thm:chartcomp} and \ref{thm:colimtrace}).  For the limit form
$I_{\infty}$ one has $I_{\infty}(F_n,G_m)\to1$ and $I_{\infty}(F_n,F_n)\to0$
along the sequences above, by Mertens' theorem \cite{Mertens}; the value $1$
is the logarithmic density of prime powers.  This limit depends on the
approximation: it holds for radial approximations whose profiles vary
slowly, and there are approximations of the rulings, of the kind considered
above, along which
$I_{\infty}$ tends to $0$ (Theorem~\ref{thm:mertens}).  Mertens' theorem is
used only in Appendix~\ref{app:mertens}.

Section~\ref{sec:carrier} introduces the space and the test algebra.
Sections~\ref{sec:cat} and \ref{sec:dagger} treat the graph bimodules, their
composition, excision, and the involution, and Section~\ref{sec:shadow}
computes their Hochschild complexes.  Sections~\ref{sec:trace} and
\ref{sec:symbol} define the trace and compute its symbol.
Section~\ref{sec:numerical} studies the pairing, the Tate sequence and the
completion.  Sections~\ref{sec:presented}--\ref{sec:charts} treat
presentations, the Banach algebra and the change of $S$.
Sections~\ref{sec:eff}--\ref{sec:classical} treat the cone of non-negative
test functions, the Riemann--Roch pairing and the comparison with
correspondences on a curve, and Section~\ref{sec:signature} computes the
signature.

\section{The semilocal space and the test algebra}\label{sec:carrier}

Throughout, $S$ is a finite set of places of $\Q$ with $\infty\in S$ and
$S_f:=S\setminus\{\infty\}\ne\emptyset$.  We write
\[
A_S:=\R\times\prod_{p\in S_f}\Q_p,\qquad
\Gamma_S:=\Z_S^{\times}\cap\Q_{>0}
=\Bigl\{\textstyle\prod_{p\in S_f}p^{e_p}: e_p\in\Z\Bigr\},
\]
$\Gamma_S$ being embedded diagonally in $A_S^{\times}$.  It is a free abelian
group of rank $|S_f|$, and it acts on $A_S$ by multiplication.  The geometric
object behind the constructions of this paper is the quotient groupoid and its
square,
\begin{equation}\label{eq:carrier}
X_S:=[A_S/\Gamma_S],\qquad Y_S:=X_S\times X_S .
\end{equation}
They serve as a guide only: every statement below is about the algebras and
bimodules attached to the action of $\Gamma_S$ on $A_S$.

We equip $A_S$ with the additive product Haar measure $\mu$ normalized so that
the unit ball at each finite place has volume one, and we write
$|a|_S=\prod_{v\in S}|a_v|_v$ for the module of $a\in A_S^{\times}$, so that
$d\mu(ax)=|a|_S\,d\mu(x)$.

\begin{lemma}[Product formula]\label{lem:prodformula}
$|\gamma|_S=1$ for every $\gamma\in\Gamma_S$.  Consequently multiplication by
$\gamma$ preserves $\mu$.
\end{lemma}

\begin{proof}
A positive $S$-unit $\gamma$ has $|\gamma|_v=1$ at every place $v\notin S$, so
the global product formula $\prod_{v}|\gamma|_v=1$ reduces to
$\prod_{v\in S}|\gamma|_v=1$.
\end{proof}

\subsection{Coefficients and the source algebra}

\subsubsection*{The local algebras}
Let $p$ be a finite place.  For $j\in\Z$ we call $p^{j}\Z_p^{\times}$ the
\emph{shell} of valuation $j$, and for a character $\kappa$ of
$\Z_p^{\times}$ we write $\cond_p(\kappa)\ge0$ for its conductor exponent, the
least integer $e\ge0$ with $\kappa$ trivial on $1+p^{e}\Z_p$, where
$1+p^{0}\Z_p:=\Z_p^{\times}$.  We let $F_p$ be the space of functions $f$ on $\Q_p$ with
$f(0)=0$ which have an expansion
\[
f(p^{j}u)=\sum_{\kappa}f_{j,\kappa}\,\kappa(u)\qquad(j\in\Z,\ u\in\Z_p^{\times})
\]
such that
\begin{equation}\label{eq:seminorms}
v_{l,N}(f)=\sum_{j,\kappa}p^{l|j|}\,(1+\cond_p(\kappa))^{N}\,|f_{j,\kappa}|<\infty ,
\qquad l,N\ge0 .
\end{equation}
With these seminorms $F_p$ is a Fr\'echet space.  Since
$\sup_{u}|f(p^{j}u)|\le p^{-l|j|}v_{l,0}(f)$, the elements of $F_p$ are
continuous functions on $\Q_p$ which decay faster than every power of
$p^{-|j|}$ as $x_p\to0$ and as $x_p\to\infty$.  Under pointwise multiplication
$(fg)_{j,\kappa}=\sum_{\kappa'\kappa''=\kappa}f_{j,\kappa'}g_{j,\kappa''}$, and
$\cond_p(\kappa'\kappa'')\le\max(\cond_p(\kappa'),\cond_p(\kappa''))$, so
\[
v_{l,N}(fg)\le v_{l,N}(f)\,v_{l,N}(g);
\]
thus $F_p$ is a Fr\'echet algebra.  We write $1_{v_p=j}$ for the indicator
function of the shell of valuation $j$; it lies in $F_p$, it is idempotent,
and $v_{l,N}(1_{v_p=j})=p^{l|j|}$.

Put $\chi_p:=\mathbf 1_{\Z_p}$ and
\[
E_p:=\C\cdot\chi_p\oplus F_p,\qquad
v^{E}_{l,N}(\lambda\chi_p+f):=|\lambda|+v_{l,N}(f).
\]
Since $\chi_p^{2}=\chi_p$ and $\chi_pf=\sum_{j\ge0}1_{v_p=j}f\in F_p$ with
$v_{l,N}(\chi_pf)\le v_{l,N}(f)$, the space $E_p$ is a Fr\'echet algebra of
continuous functions on $\Q_p$ under pointwise multiplication, the seminorms
$v^{E}_{l,N}$ are submultiplicative, and $F_p$ is a closed ideal of
codimension one.  The summand $\C\chi_p$ records the value at $x_p=0$.

At the archimedean place, $E_{\R}$ is the algebra, under pointwise
multiplication, of smooth functions $f$ on $\R^{\times}$ such that, on each of
the two half-lines and in the logarithmic coordinate $t=\log|x_{\infty}|$,
\[
p_{l,m}(f):=\max_{m'\le m}\,\sup_{x_{\infty}\in\R^{\times}}e^{l|t|}\,|D^{m'}f|<\infty
\qquad(l,m\ge0,\ D=d/dt).
\]
Thus $E_{\R}$ is the direct sum of two copies, one for each sign of
$x_{\infty}$, of the algebra of smooth functions of $t\in\R$ all of whose
derivatives decay faster than every exponential.

\subsubsection*{The action of the ideles}
For $a\in A_S^{\times}$ and a function $f$ of $x\in A_S$ put
$(\beta_af)(x):=f(a^{-1}x)$.  On $E_{\R}$ this is translation by
$\log|a_{\infty}|$ in $t$, composed with the exchange of the two half-lines
when $a_{\infty}<0$.  On $E_p$, writing $a_p=p^{k}w$ with $w\in\Z_p^{\times}$,
it sends $f\in F_p$ to the element with coefficients
$\kappa(w)^{-1}f_{j-k,\kappa}$, and it sends $\chi_p$ to
$\mathbf 1_{p^{k}\Z_p}$, which differs from $\chi_p$ by the sum, with sign
$-\sgn(k)$, of $|k|$ shell idempotents.  Hence $\beta_a$ is a continuous
automorphism of each local algebra, preserving $F_p$, and
$\beta_a\beta_b=\beta_{ab}$.

The coefficient algebra is the completed projective tensor product
\begin{equation}\label{eq:coeff}
E_P:=E_{\R}\,\widehat\otimes\,\bigotimes_{p\in S_f}E_p ,
\end{equation}
a Fr\'echet algebra on which $A_S^{\times}$ acts by the automorphisms
$\beta_a$; its topology is defined by the projective tensor products
$\pi_{l,m}$ of the seminorms $p_{l,m}$ and $v^{E}_{l,m}$ $(p\in S_f)$.  We
write $\alpha_{\gamma}:=\beta_{\gamma}$ for $\gamma\in\Gamma_S$.  Put
\[
D:=\R^{\times}\times\prod_{p\in S_f}\Q_p .
\]

\begin{lemma}\label{lem:functions}
The map which sends an elementary tensor to the product of its factors extends
to a continuous injective algebra homomorphism $\iota$ from $E_P$ to the
algebra $C_b(D)$ of bounded continuous functions on $D$, and
\[
q(e):=\sup_{x\in D}e^{|t(x)|}\,|\iota(e)(x)|,\qquad t(x)=\log|x_{\infty}|,
\]
is a continuous seminorm on $E_P$.  The extension of $\iota$ to
$E_P\oplus\C1$ which sends $1$ to the constant function $1$ is injective.
\end{lemma}

\begin{proof}
For an elementary tensor $e=f_{\infty}\otimes\bigotimes_p(\lambda_p\chi_p+f_p)$
one has
\[
\sup_D|\iota(e)|\le p_{0,0}(f_{\infty})\prod_pv^{E}_{0,0}(\lambda_p\chi_p+f_p),
\qquad
q(e)\le p_{1,0}(f_{\infty})\prod_pv^{E}_{0,0}(\lambda_p\chi_p+f_p).
\]
Hence
$\iota$ extends continuously to $E_P$ with $\sup_D|\iota(e)|\le\pi_{0,0}(e)$
and $q\le\pi_{1,0}$.  It is multiplicative on the algebraic tensor product and
therefore on $E_P$.

For injectivity, fix $p\in S_f$ and consider on $E_p$ the continuous linear
functionals
\[
\mathrm{ev}_0(\lambda\chi_p+f):=\lambda,\qquad
\phi_{j,\kappa}(g):=\int_{\Z_p^{\times}}g(p^{j}u)\,\overline{\kappa(u)}\,du ,
\]
$du$ being the Haar probability measure; thus
$\phi_{j,\kappa}(\lambda\chi_p+f)=f_{j,\kappa}$ unless $j\ge0$ and $\kappa=1$,
in which case it is $f_{j,1}+\lambda$.  Let $P_n$ be the operator on $E_p$
which keeps $\lambda$ and the coefficients $f_{j,\kappa}$ with $|j|\le n$ and
$\cond_p(\kappa)\le n$ and replaces the other coefficients by $0$.  It has finite
rank, it is a linear combination of the functionals $\mathrm{ev}_0$ and
$\phi_{j,\kappa}$ with coefficients in $E_p$, it satisfies
$v^{E}_{l,N}(P_ng)\le v^{E}_{l,N}(g)$, and $P_ng\to g$ for every $g$.  For a
Fr\'echet space $X$ the operators $P_n\otimes\id$ on $E_p\widehat\otimes X$ are
therefore equicontinuous and converge to the identity on the algebraic tensor
product, hence everywhere.  Consequently an element $e\in E_p\widehat\otimes X$
with $(\mathrm{ev}_0\otimes\id)(e)=0$ and $(\phi_{j,\kappa}\otimes\id)(e)=0$ for
all $j,\kappa$ satisfies $(P_n\otimes\id)(e)=0$ for all $n$, so $e=0$.

Now enumerate $S_f=\{p_1,\dots,p_r\}$, put
$X_k:=E_{\R}\widehat\otimes E_{p_1}\widehat\otimes\cdots\widehat\otimes E_{p_k}$
and $D_k:=\R^{\times}\times\Q_{p_1}\times\cdots\times\Q_{p_k}$, and let
$\iota_k:X_k\to C_b(D_k)$ be the corresponding map.  The map $\iota_0$ is the
inclusion $E_{\R}\subset C_b(\R^{\times})$.  Suppose that $\iota_{k-1}$ is
injective and let $e\in X_k=X_{k-1}\widehat\otimes E_{p_k}$ with
$\iota_k(e)=0$.  For $x'\in D_{k-1}$,
\[
\iota_{k-1}\bigl((\id\otimes\phi_{j,\kappa})e\bigr)(x')
=\int_{\Z_{p_k}^{\times}}\iota_k(e)(x',p_k^{j}u)\,\overline{\kappa(u)}\,du,
\qquad
\iota_{k-1}\bigl((\id\otimes\mathrm{ev}_0)e\bigr)(x')=\iota_k(e)(x',0):
\]
both identities hold for elementary tensors, and both sides are continuous in
$e$.  The right-hand sides vanish, so by the inductive hypothesis
$(\id\otimes\phi_{j,\kappa})e=0$ and $(\id\otimes\mathrm{ev}_0)e=0$, and $e=0$
by the preceding paragraph.

Finally, if $\lambda+\iota(e)=0$ with $\lambda\in\C$ and $e\in E_P$, then
$|\lambda|=|\iota(e)(x)|\le q(e)e^{-|t(x)|}$ for every $x$, so $\lambda=0$ and
then $e=0$.
\end{proof}

From now on we identify the elements of $E_P\oplus\C1$ with functions on $D$.

\subsubsection*{The rapid crossed product}
The \emph{height} of a positive rational written in lowest terms is
$H(a/b):=\max(a,b)$; it is symmetric, $H(\gamma^{-1})=H(\gamma)$, and
submultiplicative, $H(\gamma\delta)\le H(\gamma)H(\delta)$.  For a Fr\'echet
algebra $A$ with an action $\alpha$ of $\Gamma_S$ by automorphisms, the
\emph{rapid crossed product} $A\rtimes_{\mathrm{rapid}}\Gamma_S$ consists of the
sums $\sum_{\gamma}a_{\gamma}U_{\gamma}$ with
$\sum_{\gamma}H(\gamma)^{k}\nu(a_{\gamma})<\infty$ for every $k\ge0$ and every
continuous seminorm $\nu$ of $A$, with the product determined by
$U_{\gamma}U_{\delta}=U_{\gamma\delta}$ and
$U_{\gamma}aU_{\gamma}^{-1}=\alpha_{\gamma}(a)$.  The action is \emph{tempered}
if for every continuous seminorm $\nu$ there are a continuous seminorm $\nu'$
and constants $C$, $M$ with $\nu(\alpha_{\gamma}f)\le C\,H(\gamma)^{M}\nu'(f)$
for all $\gamma$ and $f$.  Temperedness and the submultiplicativity of $H$
make the product of the rapid crossed product continuous, so that it is a
Fr\'echet algebra.

The action of $\Gamma_S$ on $E_P$ is tempered.  At a finite place $p$ the
automorphism $\alpha_{\gamma}$ shifts the shell index by $v_p(\gamma)$ and
multiplies each coefficient by a number of modulus one, so that
$v_{l,N}(\alpha_{\gamma}f)\le p^{l|v_p(\gamma)|}v_{l,N}(f)\le
H(\gamma)^{l}v_{l,N}(f)$ for $f\in F_p$; and $\alpha_{\gamma}\chi_p-\chi_p$ is,
up to sign, the sum of $|v_p(\gamma)|$ shell idempotents of valuation at most
$|v_p(\gamma)|$ in absolute value, of $v_{l,N}$-size at most
$|v_p(\gamma)|\,p^{l|v_p(\gamma)|}\le H(\gamma)^{l+1}$.  At the real place
$\alpha_{\gamma}$ translates $t$ by $\log\gamma$, and
$p_{l,m}(\alpha_{\gamma}f)\le e^{l|\log\gamma|}p_{l,m}(f)\le
H(\gamma)^{l}p_{l,m}(f)$.

The \emph{source algebra} is
\begin{equation}\label{eq:source}
U:=(E_P\oplus\C 1)\rtimes_{\mathrm{rapid}}\Gamma_S ,
\end{equation}
a unital Fr\'echet algebra, and we write
$N:=E_P\rtimes_{\mathrm{rapid}}\Gamma_S$ for its nonunital coefficient ideal.
Complex conjugation preserves $E_P$, and
\[
\Bigl(\sum_{\gamma}a_{\gamma}U_{\gamma}\Bigr)^{*}:=
\sum_{\gamma}\alpha_{\gamma}^{-1}(\overline{a_{\gamma}})\,U_{\gamma^{-1}}
\]
is a continuous involution of $U$, by temperedness and the symmetry of $H$.
For it $U_{\gamma}^{*}=U_{\gamma}^{-1}$; we call the $U_{\gamma}$ the
\emph{gauge unitaries}.  For $a\in A_S^{\times}$ the automorphism $\beta_a$ of
$E_P\oplus\C1$ commutes with every $\alpha_{\gamma}$, multiplication in
$A_S^{\times}$ being commutative; it therefore extends to a continuous
automorphism of $U$, again denoted $\beta_a$, which fixes every $U_{\gamma}$.
It preserves $N$, and $\beta_{\gamma}$ is conjugation by $U_{\gamma}$ for
$\gamma\in\Gamma_S$.

\subsubsection*{Modules}
Modules are unital Fr\'echet modules over $U$, left, right or two-sided.  A
complex of modules is \emph{admissibly exact} if it has a continuous linear
contraction, and a short exact sequence is \emph{admissible} if it splits
continuously and linearly.  The \emph{free} right, left and two-sided modules
on a Fr\'echet space $V$ are $V\widehat\otimes U$, $U\widehat\otimes V$ and
$U\widehat\otimes V\widehat\otimes U$, and a module is \emph{relatively
projective} if it is a module retract of a free one; equivalently, if it has
the lifting property for admissible surjections.  A relatively projective
bimodule is relatively projective on each side.  If $P$ is a relatively
projective right module, then $P\,\widehat\otimes_U\,(-)$ carries admissibly
exact complexes of left modules to admissibly exact complexes: for
$P=V\widehat\otimes U$ this functor is $V\widehat\otimes(-)$, which preserves
linear contractions, and if $i:P\to V\widehat\otimes U$ and
$r:V\widehat\otimes U\to P$ are module maps with $ri=\id$ and $h$ is a
contraction of $V\widehat\otimes C$, then $(r\otimes\id)h(i\otimes\id)$ is a
contraction of $P\widehat\otimes_UC$.

We write $D^{-}_{\mathrm{rel}}(U\text{-}U)$ for the relative derived category of
bounded-above complexes of bimodules, obtained by inverting the maps with
admissibly exact cone.  We use the following standard facts of relative
homological algebra, for which we refer to \cite{Helemskii}.  Every
bounded-above complex $M$ has a resolution $Q(M)\to M$ by a bounded-above
complex of relatively projective bimodules, with admissibly exact cone and
unique up to homotopy equivalence.  The category
$D^{-}_{\mathrm{rel}}(U\text{-}U)$ is triangulated, admissible short exact
sequences of complexes giving distinguished triangles, and
\begin{equation}\label{eq:comp}
M\star M':=Q(M)\,\widehat\otimes_U\,Q(M'),\qquad \Delta:=U,
\end{equation}
is a triangulated functor of each variable, associative up to natural
isomorphism, with unit $\Delta$.  If the terms of $M$ are relatively projective
as right modules, or those of $M'$ as left modules, the canonical map
$M\star M'\to M\widehat\otimes_UM'$ is an isomorphism in the relative derived
category of the remaining module structures.  The same notions and facts apply
to modules over the quotient algebras of $U$ considered in
Section~\ref{sec:cat}, and we write $\widehat\otimes^{\mathbf L}_U$ for the
derived tensor product of one-sided modules.

For a bounded-above complex $M$ of bimodules we call
\begin{equation}\label{eq:shadow}
\Sh(M):=\mathrm{Tot}\bigl[M\,\widehat\otimes\,U^{\widehat\otimes n}\bigr]_{n\ge0},
\end{equation}
with the Hochschild boundary, the \emph{relative Hochschild shadow} of $M$.

\subsection{The valuation splitting}

\begin{lemma}\label{thm:split}
Define
\begin{equation}\label{eq:phi}
\varphi:A_S^{\times}\longrightarrow\Gamma_S,\qquad
\varphi(x):=\prod_{p\in S_f}p^{\,v_p(x_p)} .
\end{equation}
Then $\varphi$ is a continuous group homomorphism with
$\varphi|_{\Gamma_S}=\id$, and therefore, as topological groups,
\begin{equation}\label{eq:splitting}
A_S^{\times}=\Gamma_S\times\ker\varphi,\qquad
\ker\varphi=\R^{\times}\times\prod_{p\in S_f}\Z_p^{\times},\qquad
C_S:=\Gamma_S\backslash A_S^{\times}\cong\ker\varphi .
\end{equation}
Moreover $|x|_S=|x_{\infty}|_{\infty}$ for $x\in\ker\varphi$.
\end{lemma}

\begin{proof}
Each $v_p$ is a homomorphism $\Q_p^{\times}\to\Z$, so $\varphi$ is a
homomorphism, and it is continuous for the discrete topology of $\Gamma_S$
because $v_p$ is locally constant on $\Q_p^{\times}$.  If
$\gamma=\prod_{p\in S_f}p^{e_p}\in\Gamma_S$ is embedded diagonally then
$v_p(\gamma)=e_p$ for $p\in S_f$, so $\varphi(\gamma)=\gamma$.  In particular
$\Gamma_S$ is discrete in $A_S^{\times}$, and
$x\mapsto(\varphi(x),x\varphi(x)^{-1})$ is an isomorphism of topological groups
from $A_S^{\times}$ onto $\Gamma_S\times\ker\varphi$ with inverse the
multiplication map.  The kernel is
$\{x: v_p(x_p)=0\text{ for all }p\in S_f\}$, which is \eqref{eq:splitting}.
For such $x$ one has $|x_p|_p=1$ at every finite place of $S$, so
$|x|_S=|x_{\infty}|$.
\end{proof}

We fix once and for all the resulting section
\begin{equation}\label{eq:section}
\sigma: C_S\;\xrightarrow{\ \cong\ }\;\ker\varphi\;\hookrightarrow\;A_S^{\times}
\end{equation}
of the class map; it is a homomorphism.  We normalize Haar measure
$\dstar w$ on $C_S\cong\R^{\times}\times\prod_{p\in S_f}\Z_p^{\times}$ so that
each compact factor $\Z_p^{\times}$ has total mass one and the archimedean
factor carries
\begin{equation}\label{eq:archmeasure}
\dstar x_{\infty}=\frac{dx}{2|x|},
\end{equation}
so that the two signs together have total mass one per unit of logarithmic
length.  In the coordinate $t=\log|w|_S$ we therefore have
$\dstar w=dt\otimes\mu_{\mathrm{cpt}}$ with $\mu_{\mathrm{cpt}}$ a probability
measure on the compact group $\{\pm1\}\times\prod_{p\in S_f}\Z_p^{\times}$.

\begin{remark}\label{rem:normdict}
Two normalizations of the involution occur in this subject.  The
unnormalized one is $f^{\sharp}(u)=|u|_S^{-1}\overline{f(u^{-1})}$, paired
with the local density $|1-u^{-1}|_v^{-1}$; the half-density one, used
throughout this paper, is $F^{\dagger}(u)=\overline{F(u^{-1})}$, paired with
$\tau_v(u)=|u|_v^{1/2}|1-u|_v^{-1}$.  The substitution
$f\mapsto F=|\cdot|_S^{1/2}f$ is an isomorphism of topological algebras for
convolution, $|\cdot|_S^{1/2}$ being a character, and it carries the first
convention to the second:
$|u|_S^{1/2}f^{\sharp}(u)=|u|_S^{-1/2}\overline{f(u^{-1})}=F^{\dagger}(u)$, and
$|u|_v^{-1/2}\,|1-u^{-1}|_v^{-1}=\tau_v(u)$.  All formulas below are in the
half-density convention.
\end{remark}

\subsection{The test algebra}

For $n\ge0$ put $\Theta_n:=\prod_{q\in S_f}(1+q^{n}\Z_q)$, with the convention
$1+q^{0}\Z_q:=\Z_q^{\times}$; these are open subgroups of the compact part of
$C_S$ and form a neighbourhood basis of $1$ in $\prod_q\Z_q^{\times}$.

\begin{definition}\label{def:DA}
$\DA:=C_c^{\infty}(C_S)$ is the space of functions on $C_S$ which are smooth
in the $\R^{\times}$ coordinate, locally constant in the
$\prod_p\Z_p^{\times}$ coordinates, and compactly supported.  It carries the
strict inductive limit topology of the Fr\'echet spaces of functions supported
in a fixed compact set and invariant under a fixed $\Theta_n$; this topology
is nuclear and of type LF.  We equip $\DA$ with the convolution
\[
(F\star G)(w)=\int_{C_S}F(u)G(u^{-1}w)\dstar u
\]
and the involution $F^{\dagger}(u)=\overline{F(u^{-1})}$.  We write
\[
\DA^{\mathrm{off}}:=\{F\in\DA:\ F\equiv0\ \text{on a neighbourhood of }1\}.
\]
\end{definition}

$\DA$ is a commutative $\dagger$-algebra: associativity and commutativity of
convolution are Fubini's theorem on the locally compact abelian group $C_S$,
and $(F\star G)^{\dagger}=F^{\dagger}\star G^{\dagger}$ follows from the
invariance of Haar measure under inversion.  Each step is a closed subspace of
the later ones with the induced topology, so a bounded subset of $\DA$ lies in
a single step \cite[Ch.~II, \S6]{Schaefer}: there are a compact
$K\subseteq C_S$ and an integer $n\ge0$ such that all its members are supported
in $K$, are invariant under $\Theta_n$, and have uniformly bounded step
seminorms.

\begin{remark}\label{rem:approxid}
$\DA$ has an approximate identity.  Fix $\phi\in C_c^{\infty}(\R)$ with
$\phi\ge0$, $\supp\phi\subseteq(-1,1)$ and $\int\phi=1$, and for
$\varepsilon>0$ and $n\ge1$ put
\[
e_{\varepsilon,n}(w):=2\,\varepsilon^{-1}\phi\bigl(\varepsilon^{-1}\log w_{\infty}\bigr)
\mathbf 1_{w_{\infty}>0}\cdot\vol(\Theta_n)^{-1}\,
\mathbf 1_{\Theta_n}\bigl((w_q)_{q\in S_f}\bigr).
\]
Then $e_{\varepsilon,n}\in\DA$, and $\int e_{\varepsilon,n}\dstar w=1$ by
\eqref{eq:archmeasure}.  If $F\in\DA$ is $\Theta_n$-invariant, then
\[
(F\star e_{\varepsilon,n})(w)=\int_{\R}F(e^{-s}w)\,\varepsilon^{-1}\phi(s/\varepsilon)\,ds ,
\]
where $e^{-s}$ multiplies the archimedean coordinate.  This function is
$\Theta_n$-invariant and, for $\varepsilon\le1$, supported in the compact set
$[e^{-1},e]\cdot\supp F$, and it converges to $F$ uniformly together with all
its derivatives in $t$ as $\varepsilon\to0$.  Hence
$F\star e_{\varepsilon,n}\to F$ in a single Fr\'echet step as
$\varepsilon\to0$ and $n\to\infty$.  The same holds for
$e_{\varepsilon,n}^{\dagger}$, which is of the same form with
$\phi(-\,\cdot\,)$ in place of $\phi$.
\end{remark}

\section{The correspondence category and excision}\label{sec:cat}

\subsection{Graph bimodules}

\begin{convention}\label{conv:twist}
For a continuous automorphism $\beta$ of $U$ we write ${}_{1}U_{\beta}$ for the
bimodule whose underlying space is $U$, with left action by multiplication and
right action twisted by $\beta$:
\[
x\cdot m\cdot y:=x\,m\,\beta(y)\qquad(x,y,m\in U).
\]
The map $m\mapsto\beta^{-1}(m)$ is an isomorphism of ${}_{1}U_{\beta}$ onto the
bimodule ${}_{\beta^{-1}}U_{1}$ with actions $x\cdot m\cdot y=\beta^{-1}(x)\,m\,y$,
so a right twist by $\beta$ is the same as a left twist by $\beta^{-1}$.
\end{convention}

For $a\in A_S^{\times}$ the \emph{graph bimodule} of $a$ is
\begin{equation}\label{eq:graphbimod}
E_a:={}_{1}U_{\beta_a} .
\end{equation}
It plays the role of the structure sheaf of the graph of multiplication by $a$
on $X_S$.  As a left module and as a right module it is free of rank one, with
generator $1$.  For $u,v\in A_S^{\times}$ the map
\begin{equation}\label{eq:starmult}
E_u\,\widehat\otimes_U\,E_v\longrightarrow E_{uv},\qquad
x\otimes y\mapsto x\,\beta_u(y),
\end{equation}
is an isomorphism of bimodules.  Indeed it is balanced, since
$x\beta_u(w)\otimes y$ and $x\otimes wy$ both go to $x\beta_u(w)\beta_u(y)$; it
is right linear, since $x\otimes y\beta_v(w)$ goes to
$x\beta_u(y)\beta_{uv}(w)$; and $z\mapsto z\otimes1$ is its inverse, because
$x\otimes y=x\beta_u(y)\otimes1$.  The isomorphisms \eqref{eq:starmult} are
associative: both bracketings of $x\otimes y\otimes z$ in
$E_u\widehat\otimes_UE_v\widehat\otimes_UE_w$ go to
$x\beta_u(y)\beta_{uv}(z)$.  Since $E_u$ is free on each side,
$E_u\star E_v\cong E_u\widehat\otimes_UE_v\cong E_{uv}$, and
$E_1=\Delta=U$.

\begin{proposition}\label{prop:hom}
For $a,b\in A_S^{\times}$,
\begin{equation}\label{eq:homclass}
\Hom_{U\text{-}U}(E_a,E_b)=
\begin{cases}
\C\cdot U_{\gamma}, & a=\gamma b\ \text{with}\ \gamma\in\Gamma_S,\\
0, & \text{otherwise},
\end{cases}
\end{equation}
the element $U_{\gamma}$ acting by right multiplication; since $\Gamma_S$ acts
freely on $A_S^{\times}$, at most one $\gamma$ occurs.  Consequently $E_a\cong
E_b$ if and only if $a\in\Gamma_Sb$, and every bounded complex of finite
direct sums of bimodules $E_a$, with bimodule maps as differentials, is
homotopy equivalent through bimodule maps to a finite direct sum of shifts of
bimodules $E_a$.
\end{proposition}

\begin{proof}
Let $\Omega:=\R^{\times}\times\prod_{p\in S_f}\Q_p^{\times}$, a dense open
subset of $D$ stable under $A_S^{\times}$.  The algebra $E_P$ separates the
points of $\Omega$: $E_{\R}$ contains $C_c^{\infty}(\R^{\times})$, $E_p$
contains the indicator function of every compact open subset of
$\Q_p^{\times}$, which is a finite combination of functions
$p^{j}u\mapsto\kappa(u)$ supported on one shell, and products of such
functions separate the points of the product.

A bimodule map $f:E_a\to E_b$ satisfies $f(x)=xz$ with $z:=f(1)$, by left
linearity; and $f(x\beta_a(y))=f(x)\beta_b(y)$ for all $x,y$ is equivalent to
\begin{equation}\label{eq:intertw}
\beta_a(y)\,z=z\,\beta_b(y)\qquad(y\in U).
\end{equation}
Conversely every $z$ with \eqref{eq:intertw} defines a continuous bimodule
map.  Write $z=\sum_{\gamma}z_{\gamma}U_{\gamma}$ with $z_{\gamma}\in
E_P\oplus\C1$.  Take $y=c\in E_P$.  Since
$U_{\gamma}\beta_b(c)=\alpha_{\gamma}(\beta_b(c))U_{\gamma}$, comparing the
coefficients of $U_{\gamma}$ in \eqref{eq:intertw} and using that the
identification of Lemma~\ref{lem:functions} is multiplicative gives
\[
\bigl(c(a^{-1}x)-c(b^{-1}\gamma^{-1}x)\bigr)\,z_{\gamma}(x)=0
\qquad(x\in D,\ c\in E_P).
\]
If $a\ne\gamma b$, then for $x\in\Omega$ the points $a^{-1}x$ and
$(\gamma b)^{-1}x$ of $\Omega$ are distinct, because every coordinate of $x$ is
nonzero; choosing $c$ which separates them gives $z_{\gamma}(x)=0$.  So
$z_{\gamma}$ vanishes on $\Omega$, hence on $D$ by continuity and density,
and therefore $z_{\gamma}=0$ by Lemma~\ref{lem:functions}.
Thus $z=0$ unless $a=\gamma b$ for some $\gamma$, which is then unique, and in
that case $z=hU_{\gamma}$ with $h\in E_P\oplus\C1$.

Suppose $a=\gamma b$ and $z=hU_{\gamma}$.  Taking $y=U_{\delta}$ in
\eqref{eq:intertw}, and using that $\beta_a,\beta_b$ fix $U_{\delta}$ and that
$\Gamma_S$ is abelian, gives $\alpha_{\delta}(h)=h$ for every
$\delta\in\Gamma_S$.  Write $h=\lambda+e$ with $\lambda\in\C$ and $e\in E_P$;
then $e$ is invariant.  Fix $p\in S_f$.  Invariance gives $e(x)=e(p^{n}x)$ for
all $n\in\Z$, and $t(p^{n}x)=t(x)+n\log p$, so with the seminorm $q$ of
Lemma~\ref{lem:functions}
\[
|e(x)|\le q(e)\,e^{-|t(x)+n\log p|}\longrightarrow0\qquad(n\to\infty).
\]
Hence $e$ vanishes on $D$, so $e=0$ and $z=\lambda U_{\gamma}$.  Conversely, if
$a=\gamma b$ then $U_{\gamma}\beta_b(y)U_{\gamma}^{-1}=\beta_a(y)$ for
$y\in E_P$, since $\alpha_{\gamma}\beta_b=\beta_{\gamma b}$, and for
$y=U_{\delta}$ trivially; so $z=U_{\gamma}$ satisfies \eqref{eq:intertw}.
This proves \eqref{eq:homclass}.

By \eqref{eq:homclass}, right multiplication by $U_{\gamma}$ is an isomorphism
$E_{\gamma a}\to E_a$ with inverse $U_{\gamma^{-1}}$, the endomorphisms of $E_a$
are the scalars, and there are no nonzero maps between $E_a$ and $E_b$ when
$a\notin\Gamma_Sb$.  Let $K$ be a bounded complex of finite sums of bimodules
$E_a$.  Replacing each summand by an isomorphic one, we may assume that the
summands of a given class in $C_S$ are all equal to one $E_a$.  Then
$K=\bigoplus_{[a]}V_{[a]}^{\bullet}\otimes E_a$, a finite sum, where each
$V_{[a]}^{\bullet}$ is a bounded complex of finite-dimensional vector spaces
and the differentials are scalar matrices.  Each $V^{\bullet}_{[a]}$ is
homotopy equivalent to its cohomology with zero differential, by linear
algebra, and hence
$K\simeq\bigoplus_{[a]}\bigoplus_nH^{n}(V_{[a]})\otimes E_a[-n]$.
\end{proof}

\begin{definition}\label{def:G}
$G$ is the differential graded category whose objects are the bounded
complexes $K$ of finite direct sums of the bimodules $E_a$, $a\in\ker\varphi$,
with bimodule maps as differentials, and whose morphism complexes are
\[
\Hom^{n}(K,L):=\prod_{i}\Hom_{U\text{-}U}(K^{i},L^{i+n}),\qquad
d(f):=d_L\circ f-(-1)^{n}f\circ d_K ,
\]
the product being finite.  By Proposition~\ref{prop:hom} every $E_a$,
$a\in A_S^{\times}$, is isomorphic to exactly one generator
$E_{a\varphi(a)^{-1}}$ of $G$, and the morphisms between generators are
\[
\Hom_{U\text{-}U}(E_a,E_b)=\begin{cases}\C,&a=b,\\0,&a\ne b\end{cases}
\qquad(a,b\in\ker\varphi).
\]
\end{definition}

The set $\ker\varphi=\sigma(C_S)$ of labels is a subgroup of $A_S^{\times}$.
The morphisms of $G$ are bimodule maps; we use the evident functor from $G$ to
$D^{-}_{\mathrm{rel}}(U\text{-}U)$ on objects only.

\begin{definition}\label{def:smear}
A \emph{smeared object} is a function on $C_S$ of the form
\begin{equation}\label{eq:smear}
E_F:\ u\longmapsto F(u)\cdot1\in E_{\sigma(u)},\qquad F\in\DA ,
\end{equation}
where $1$ is the generator of $E_{\sigma(u)}$.  For
$u,v\in C_S$ let
$\mu_{u,v}:E_{\sigma(u)}\widehat\otimes_UE_{\sigma(v)}\to
E_{\sigma(u)\sigma(v)}=E_{\sigma(uv)}$ be the isomorphism \eqref{eq:starmult}.
The product and the involution of smeared objects are
\[
(E_F\star E_G)(w):=\int_{C_S}\mu_{u,u^{-1}w}\bigl(E_F(u)\otimes E_G(u^{-1}w)\bigr)\dstar u
\in E_{\sigma(w)},\qquad
E_F^{\dagger}(u):=\overline{F(u^{-1})}\cdot1\in E_{\sigma(u)} .
\]
\end{definition}

\begin{corollary}\label{cor:trivial}
The section $\sigma$ of \eqref{eq:section} satisfies
$\beta_{\sigma(u)}\beta_{\sigma(v)}=\beta_{\sigma(uv)}$, and
$\mu_{u,v}(1\otimes1)=1$ for all $u,v\in C_S$.  Consequently
\[
E_F\star E_G=E_{F\star G},\qquad E_F^{\dagger}=E_{F^{\dagger}}\qquad(F,G\in\DA),
\]
and $F\mapsto E_F$ is a bijection of $\DA$ onto the smeared objects.
\end{corollary}

\begin{proof}
$\sigma$ is a group homomorphism by Lemma~\ref{thm:split} and
$\beta_a\beta_b=\beta_{ab}$, which gives the first identity and the equality
$E_{\sigma(u)\sigma(v)}=E_{\sigma(uv)}$ used in
Definition~\ref{def:smear}.  By \eqref{eq:starmult},
$\mu_{u,v}(1\otimes1)=1\cdot\beta_{\sigma(u)}(1)=1$.  Hence the integrand in
the definition of $E_F\star E_G$ is $F(u)G(u^{-1}w)\cdot1$, and the integral
is $(F\star G)(w)\cdot1$.  The identity for $\dagger$ is the definition of
$F^{\dagger}$, and $F$ is recovered from $E_F$ as the coefficient of the
generator.
\end{proof}

\begin{remark}\label{rem:types}
For a set-theoretic section $s$ of the class map which is not a homomorphism,
$s(u)s(v)=c(u,v)s(uv)$ with $c(u,v)\in\Gamma_S$, and the bimodules
$E_{s(u)s(v)}$ and $E_{s(uv)}$ are identified only through the gauge unitary
$U_{c(u,v)}$; the product of the generators is then $U_{c(u,v)}$ and not $1$.
The homomorphic section $\sigma$ removes this cocycle.  A smeared object is
not an object of $G$, which contains only finite sums of single twists; the
two are related by the multiplicative map $F\mapsto E_F$ and by nothing else.
\end{remark}

\begin{proposition}\label{prop:cat}
\begin{enumerate}
\item[(i)] $G$ is closed under shifts, finite direct sums and cones, and its
homotopy category is triangulated.
\item[(ii)] For objects $K,L$ of $G$ the total complex of the double complex
$K^{i}\widehat\otimes_UL^{j}$, with each summand
$E_a\widehat\otimes_UE_b$ identified with $E_{ab}$ by \eqref{eq:starmult}, is
an object $K\star L$ of $G$ which represents the product \eqref{eq:comp}.  On
generators $E_a\star E_b=E_{ab}$, the product of the scalar endomorphisms
$\lambda$ of $E_a$ and $\mu$ of $E_b$ is $\lambda\mu$, and $\Delta=E_1$ is the
unit.
\item[(iii)] In $E_a$ one has $\beta_a(c)\cdot1=1\cdot c$ for every
$c\in E_P$: the generator is annihilated by the elements
$\beta_a(c)\otimes1-1\otimes c$, which are the functions of $(y,x)\in D\times D$
of the form $c(a^{-1}y)-c(x)$ and vanish on the graph $y=ax$.
\item[(iv)] The generators $E_a$, $a\in\ker\varphi$, are pairwise
non-isomorphic, and there are infinitely many of them.
\end{enumerate}
\end{proposition}

\begin{proof}
(i)  Shifts and finite sums of objects are objects.  For a closed morphism
$f:K\to L$ of degree zero the cone is $L\oplus K[1]$ with differential
$\left(\begin{smallmatrix}d_L&f\\0&-d_K\end{smallmatrix}\right)$, which
squares to zero because $d_Lf=fd_K$; it is again a bounded complex of finite
sums of generators.  The homotopy category of a differential graded category
closed under shifts and cones is triangulated, with the cone triangles as
distinguished triangles \cite{BondalKapranov}.

(ii)  The labels form the group $\ker\varphi$, so $E_{ab}$ is a generator when
$E_a$ and $E_b$ are.  The tensor product of right multiplication by $\lambda$
on $E_a$ and by $\mu$ on $E_b$ is right multiplication by $\lambda\mu$ on
$E_{ab}$ under \eqref{eq:starmult}, so the differentials of the total complex
are scalar matrices and $K\star L$ is an object of $G$.  The terms of $K$ and
$L$ are free on each side, so the total complex represents
$Q(K)\widehat\otimes_UQ(L)$.  The unit property is $E_1=U$.

(iii)  The right action of $c$ on the generator is $1\cdot c=\beta_a(c)$,
which is also $\beta_a(c)\cdot1$.

(iv)  By \eqref{eq:homclass}, $E_a\cong E_b$ if and only if $a\in\Gamma_Sb$,
and distinct elements of $\ker\varphi$ lie in distinct classes.  The group
$\ker\varphi\cong C_S$ is infinite, since $|\cdot|_S$ maps it onto $\R_{>0}$ by
Lemma~\ref{thm:split}.
\end{proof}

\begin{remark}\label{rem:semisimple}
By Definition~\ref{def:G} the endomorphism algebra of every generator is $\C$
and there are no morphisms between distinct generators.  Hence $G$ is
equivalent to the category of bounded complexes of finite-dimensional
$C_S$-graded complex vector spaces, the grading recording the label of each
summand, with $\star$ given on generators by multiplication of labels and
$\Delta$ the unit label.
\end{remark}

\subsection{Two-sided excision}

\subsubsection*{The localized ideals}

For $J\subseteq S_f$ let
\begin{equation}\label{eq:UJ}
\begin{gathered}
\Omega_J:=\bigl\{x\in D:\ x_p\ne0\ \text{for every }p\in J\bigr\},\\
E_P^{J}:=E_{\R}\,\widehat\otimes\bigotimes_{p\in J}F_p\,\widehat\otimes
\bigotimes_{p\in S_f\setminus J}E_p,\qquad
I_J:=E_P^{J}\rtimes_{\mathrm{rapid}}\Gamma_S .
\end{gathered}
\end{equation}
Since $F_p$ is a closed ideal of $E_p$, complemented by $\C\chi_p$ and stable
under every $\beta_a$, the space $E_P^{J}$ is a closed ideal of $E_P\oplus\C1$
stable under every $\beta_a$, and $I_J$ is a closed two-sided ideal of $U$
with $\beta_a(I_J)=I_J$ for every $a\in A_S^{\times}$.  The elements of
$E_P^{J}$ vanish on $D\setminus\Omega_J$, because elementary tensors do.  We
have $I_{\emptyset}=N$ and $\Omega_{S_f}=\Omega$.

The splittings $E_p=\C\chi_p\oplus F_p$ decompose $E_P$ as a finite direct sum
of closed subspaces indexed by the subsets $L\subseteq S_f$, the summand of
index $L$ being
$E_{\R}\widehat\otimes\bigotimes_{p\in L}F_p\otimes\bigotimes_{p\notin L}\C\chi_p$.
Let $C_L\subseteq U$ be the closed subspace of the sums
$\sum_{\gamma}a_{\gamma}U_{\gamma}$ all of whose coefficients lie in that
summand.  Then
\begin{equation}\label{eq:Ldecomp}
U=\Bigl(\bigoplus_{L\subseteq S_f}C_L\Bigr)\oplus\bigl(\C1
\rtimes_{\mathrm{rapid}}\Gamma_S\bigr),\qquad
I_J=\bigoplus_{L\supseteq J}C_L ,
\end{equation}
and the projection onto any set of summands is continuous.  In particular
$I_J\cap I_K=I_{J\cup K}$, and $0\to I_J\to U\to U/I_J\to0$ is admissible.
The decomposition \eqref{eq:Ldecomp} is one of topological vector spaces; the
summands $C_L$ are not stable under multiplication or under $\Gamma_S$.

\subsubsection*{The adapted partition}

For $l,m,k\ge0$ let $\pi^{J}_{l,m}$ be the projective tensor product, on
$E_P^{J}$, of the seminorms $p_{l,m}$ on $E_{\R}$, $v_{l,m}$ on $F_p$
$(p\in J)$ and $v^{E}_{l,m}$ on $E_p$ $(p\notin J)$, and put
$\pi^{J}_{l,m,k}(a):=\sum_{\gamma}H(\gamma)^{k}\pi^{J}_{l,m}(a_{\gamma})$ for
$a=\sum_{\gamma}a_{\gamma}U_{\gamma}\in I_J$.  These seminorms increase with
$l$, $m$ and $k$ and define the topology of $I_J$, so every continuous
seminorm on $I_J$ is dominated by a multiple of one of them.

\begin{lemma}\label{lem:partition}
Fix $J\subseteq S_f$.  There are families $(e_i)_{i\in\Xi}$ and
$(\rho_i)_{i\in\Xi}$ of elements of $E_P^{J}\subseteq I_J$ such that
\begin{equation}\label{eq:partition}
e_i\rho_i=\rho_i\ \ \text{for every }i,\qquad
\sum_i\rho_i=1\ \text{ pointwise on }\Omega_J ,
\end{equation}
the sum being locally finite on $\Omega_J$, every $\rho_i$ vanishing on
$D\setminus\Omega_J$, and such that for every pair of continuous seminorms
$\alpha,\beta$ on $I_J$ there are a continuous seminorm $\nu$ on $I_J$ and a
constant $C$ with
\begin{equation}\label{eq:summable}
\sum_i\alpha(e_i)\,\beta(\rho_i a)\ \le\ C\,\nu(a)\qquad(a\in I_J).
\end{equation}
Moreover, for every $a\in I_J$ the series $\sum_i\rho_ia$ converges absolutely
in $I_J$ to $a$.
\end{lemma}

\begin{proof}
\emph{Construction.}  Choose $\psi_0\in C_c^{\infty}(\R)$ with $\psi_0\ge0$,
$\supp\psi_0\subseteq[-R,R]$ and $\sum_{n\in\Z}\psi_0(t-n)=1$, and
$\varepsilon_0\in C_c^{\infty}(\R)$ with $\varepsilon_0\equiv1$ on
$\supp\psi_0$ and $\supp\varepsilon_0\subseteq[-R-1,R+1]$.  For $n\in\Z$ and a
sign $\pm$ let $\psi_{n,\pm}$ and $\varepsilon_{n,\pm}$ be the elements of
$E_{\R}$ supported on the half-line of sign $\pm$ and equal there to
$\psi_0(t-n)$ and $\varepsilon_0(t-n)$.  At a place $p\in J$ put
$\omega_{p,j}:=1_{v_p=j}$ for $j\in\Z$; these idempotents of $F_p$ partition
$\Q_p^{\times}$.  At a place $p\in S_f\setminus J$ put $\omega_{p,0}:=\chi_p$
and $\omega_{p,j}:=1_{v_p=j}$ for $j<0$; these idempotents of $E_p$ partition
$\Q_p$.  Let $\Xi$ be the set of indices
$i=(n,\pm,(j_p)_{p\in S_f})$ with $j_p\in\Z$ for $p\in J$ and $j_p\le0$ for
$p\notin J$, and put
\[
\rho_i:=\psi_{n,\pm}\otimes\bigotimes_p\omega_{p,j_p},\qquad
e_i:=\varepsilon_{n,\pm}\otimes\bigotimes_p\omega_{p,j_p}.
\]
Then $e_i\rho_i=\rho_i$, and $\sum_i\rho_i=1$ on $\Omega_J$, locally finitely,
because at each place the chosen family partitions the relevant set and the
$\psi_{n,\pm}$ partition $\R^{\times}$.  If $x_p=0$ for some $p\in J$ then
$\omega_{p,j}(x_p)=0$ for all $j$, so $\rho_i(x)=0$.  The idempotent $\chi_p$
is used at the places outside $J$ because a nonzero value at $x_p=0$ cannot be
written as a convergent sum over shells.

\emph{The estimate.}  It suffices to prove \eqref{eq:summable} for
$\alpha=\beta=\pi^{J}_{l,m,k}$.  Write $|i|_l:=e^{l|n|}\prod_pp^{l|j_p|}$.
Since $p_{l,m}(\varepsilon_{n,\pm})\le e^{l(|n|+R+1)}\max_{m'\le m}
\sup|D^{m'}\varepsilon_0|$, $v_{l,m}(1_{v_p=j})=p^{l|j|}$ and
$v^{E}_{l,m}(\chi_p)=1$, there is a constant $C_0$ with
\begin{equation}\label{eq:sizeofe}
\pi^{J}_{l,m,k}(e_i)=\pi^{J}_{l,m}(e_i)\le C_0\,|i|_l\qquad(i\in\Xi).
\end{equation}
Put $l':=2l+2$.  We claim that there is a constant $C_1$ with
\begin{equation}\label{eq:window}
\pi^{J}_{l,m}(\rho_ic)\le C_1\,|i|_{l+2}^{-1}\,\pi^{J}_{l',m}(c)
\qquad(c\in E_P^{J},\ i\in\Xi).
\end{equation}
By the definition of the projective tensor seminorm it suffices to prove
\eqref{eq:window} for elementary tensors
$c=c_{\infty}\otimes\bigotimes_pc_p$, for which both sides are products over
the places.  At the real place, on the support of $\psi_{n,\pm}$ one has
$|t-n|\le R$, hence $e^{l|t|}\le e^{l'|t|}e^{-(l+2)(|n|-R)}$, and the Leibniz
rule gives
\[
p_{l,m}(\psi_{n,\pm}c_{\infty})\le2^{m}\max_{m'\le m}\sup|D^{m'}\psi_0|\;
e^{(l+2)R}\,e^{-(l+2)|n|}\,p_{l',m}(c_{\infty}).
\]
At a place $p\in J$, where $c_p\in F_p$,
\[
v_{l,m}(1_{v_p=j}c_p)=p^{l|j|}\sum_{\kappa}(1+\cond_p(\kappa))^{m}|(c_p)_{j,\kappa}|
\le p^{-(l+2)|j|}\,v_{l',m}(c_p).
\]
At a place $p\notin J$, for $c_p=\lambda\chi_p+f$ and $j<0$ one has
$1_{v_p=j}c_p=1_{v_p=j}f$, with the same bound, and
$v^{E}_{l,m}(\chi_pc_p)=|\lambda|+v_{l,m}(\chi_pf)\le v^{E}_{l',m}(c_p)$ for the
index $j_p=0$.  This proves \eqref{eq:window}.  For
$a=\sum_{\gamma}a_{\gamma}U_{\gamma}\in I_J$ we have
$\rho_ia=\sum_{\gamma}(\rho_ia_{\gamma})U_{\gamma}$, so \eqref{eq:sizeofe} and
\eqref{eq:window} give
\[
\sum_i\pi^{J}_{l,m,k}(e_i)\,\pi^{J}_{l,m,k}(\rho_ia)\le
C_0C_1\Bigl(\sum_{i\in\Xi}|i|_2^{-1}\Bigr)\,\pi^{J}_{l',m,k}(a),
\]
and $\sum_i|i|_2^{-1}\le2\sum_ne^{-2|n|}\prod_p\sum_jp^{-2|j|}<\infty$.

\emph{Convergence.}  Since $\pi^{J}_{l,m}(e_i)\ge1$ for every $i$, the estimate
gives $\sum_i\pi^{J}_{l,m,k}(\rho_ia)<\infty$, so $\sum_i\rho_ia$ converges
absolutely in $I_J$, to an element $b$ with coefficients
$b_{\gamma}=\sum_i\rho_ia_{\gamma}$ in $E_P^{J}$.  By
Lemma~\ref{lem:functions} this series also converges uniformly on $D$.  On
$\Omega_J$ its sum is $a_{\gamma}$ by \eqref{eq:partition}, and on
$D\setminus\Omega_J$ both $b_{\gamma}$ and $a_{\gamma}$ vanish.  Hence
$b_{\gamma}=a_{\gamma}$ as functions, and $b=a$ by Lemma~\ref{lem:functions}.
\end{proof}

\begin{lemma}\label{lem:projective}
Each $I_J$ is relatively projective as a right $U$-module and as a left
$U$-module.  Moreover the bar complex of $I_J$ admits a continuous linear
contraction, so $I_J$ is admissibly H-unital.
\end{lemma}

\begin{proof}
Put $h_J(a):=\sum_ie_i\otimes\rho_ia$.  By \eqref{eq:summable} the sum
converges absolutely in the completed projective tensor product
$I_J\widehat\otimes I_J$ and $h_J$ is continuous.  Composed with multiplication
it is the identity, by \eqref{eq:partition} and the last assertion of
Lemma~\ref{lem:partition}.  Regarded as a map into
$I_J\widehat\otimes U$, it is right $U$-linear, since
$\rho_i(ab)=(\rho_ia)b$.  Its retraction, multiplication, is also right
$U$-linear.  Hence $I_J$ is a right-module retract of the free right module
$I_J^{\flat}\widehat\otimes U$, where $I_J^{\flat}$ denotes the underlying
Fr\'echet space of $I_J$, and $I_J$ is relatively projective on the right.

For the left, put $h_J^{\mathrm{op}}(a):=\sum_i a\rho_i\otimes e_i$.  For
$a=\sum_{\gamma}a_{\gamma}U_{\gamma}$ we have
$a\rho_i=\sum_{\gamma}a_{\gamma}\alpha_{\gamma}(\rho_i)U_{\gamma}$ and
$a_{\gamma}\alpha_{\gamma}(\rho_i)=\alpha_{\gamma}\bigl(\alpha_{\gamma}^{-1}
(a_{\gamma})\rho_i\bigr)$.  By temperedness, applied to $\alpha_{\gamma}$ and
to $\alpha_{\gamma}^{-1}$, and by \eqref{eq:window}, for every $l,m$ there are
$l_1\ge l$, $l_2$, $m_2$, $C$ and $M$ with
\[
\pi^{J}_{l,m}\bigl(a_{\gamma}\alpha_{\gamma}(\rho_i)\bigr)\le
C\,H(\gamma)^{M}\,|i|_{l_1+2}^{-1}\,\pi^{J}_{l_2,m_2}(a_{\gamma}),
\]
so that, by \eqref{eq:sizeofe},
$\sum_i\pi^{J}_{l,m,k}(a\rho_i)\,\pi^{J}_{l,m,k}(e_i)\le
C'\pi^{J}_{l_2,m_2,k+M}(a)$.  Thus $h_J^{\mathrm{op}}$ is a continuous map into
$I_J\widehat\otimes I_J$, and it is left $U$-linear as a map into
$U\widehat\otimes I_J^{\flat}$.  Its composition with multiplication is
$a\mapsto\sum_ia\rho_ie_i=\sum_ia\rho_i$, because $\rho_i$ and $e_i$ lie in
the commutative algebra $E_P$; and $\sum_ia\rho_i=a$ by the argument of
Lemma~\ref{lem:partition}, since $\sum_i\alpha_{\gamma}(\rho_i)=1$ on the
$\Gamma_S$-stable set $\Omega_J$ and vanishes off it.  Hence $I_J$ is a retract
of the free left module $U\widehat\otimes I_J^{\flat}$.

Finally, let $B'(I_J)$ be the bar complex, $B'_n=I_J^{\widehat\otimes n}$ with
$b'(a_1\otimes\cdots\otimes a_n)=\sum_{j=1}^{n-1}(-1)^{j-1}a_1\otimes\cdots\otimes
a_ja_{j+1}\otimes\cdots\otimes a_n$, and put
$s(a_1\otimes\cdots\otimes a_n):=h_J(a_1)\otimes a_2\otimes\cdots\otimes a_n$.
The first term of $b's$ is $\sum_ie_i\rho_ia_1\otimes a_2\otimes\cdots=
a_1\otimes a_2\otimes\cdots$, and the remaining terms cancel $sb'$ because
$h_J(a_1a_2)=\sum_ie_i\otimes\rho_ia_1a_2$.  So $b's+sb'=\id$, and $s$ is
continuous by \eqref{eq:summable}.
\end{proof}

\begin{lemma}\label{lem:idealproduct}
For $J,K\subseteq S_f$, multiplication induces a topological isomorphism
$I_J\widehat\otimes_UI_K\cong I_{J\cup K}$ of bimodules, and consequently
\begin{equation}\label{eq:idealprod}
I_J\star I_K\cong I_{J\cup K}.
\end{equation}
\end{lemma}

\begin{proof}
Tensor the admissible sequence $0\to I_K\to U\to U/I_K\to0$ of left modules
with $I_J$, which is relatively projective on the right
(Lemma~\ref{lem:projective}).  The result is admissible, so multiplication
embeds $I_J\widehat\otimes_UI_K$ as a closed subspace of
$I_J\widehat\otimes_UU=I_J$.  Its image lies in the closure of $I_JI_K$, hence
in $I_J\cap I_K=I_{J\cup K}$ by \eqref{eq:Ldecomp}.

For surjectivity, apply $h_{J\cup K}$ of Lemma~\ref{lem:projective} to
$a\in I_{J\cup K}$: each $e_i$ lies in $I_{J\cup K}\subseteq I_J$ and each
$\rho_ia$ lies in $I_{J\cup K}\subseteq I_K$.  So $h_{J\cup K}$ followed by the
quotient map $I_J\widehat\otimes I_K\to I_J\widehat\otimes_UI_K$ is a continuous
right inverse of multiplication onto $I_{J\cup K}$.  Together with injectivity
this proves the isomorphism, including the continuity of the inverse.

Since $I_J$ is relatively projective on the right,
$I_J\star I_K\cong I_J\widehat\otimes_UI_K$, which gives \eqref{eq:idealprod}.
\end{proof}

\begin{lemma}\label{lem:cech}
Let $J_1,\dots,J_m\subseteq S_f$ be a finite nonempty list, put
$O=\bigcup_{a}\Omega_{J_a}$, $I_O=\sum_aI_{J_a}$ and $B_O=U/I_O$; for
$\emptyset\ne T\subseteq\{1,\dots,m\}$ put $J_T=\bigcup_{a\in T}J_a$.  Then:
\begin{enumerate}
\item[(i)] $I_O$ is a closed two-sided ideal stable under every $\beta_w$,
$w\in A_S^{\times}$, the sequence $0\to I_O\to U\to B_O\to0$ is admissible,
and $I_O$ is the direct sum of those summands $C_L$ of \eqref{eq:Ldecomp} for
which $J_a\subseteq L$ for some $a$;
\item[(ii)] the augmented complex
$P_O^{0}=U$, $P_O^{-n}=\bigoplus_{|T|=n}I_{J_T}$ $(1\le n\le m)$, with the
alternating inclusions and the augmentation $U\to B_O$, is admissibly exact;
its terms are relatively projective on either side, and the terms of negative
degree form an admissibly exact resolution of $I_O$;
\item[(iii)] $I_{J_T}\widehat\otimes_UB_O=0$ and
$B_O\widehat\otimes_UI_{J_T}=0$ for every $T\ne\emptyset$;
\item[(iv)] $B_O\widehat\otimes^{\mathbf L}_UB_O\cong B_O$.
\end{enumerate}
\end{lemma}

\begin{proof}
(i)  By \eqref{eq:Ldecomp}, $I_{J_a}=\bigoplus_{L\supseteq J_a}C_L$, so the sum
$I_O$ is the direct sum of the $C_L$ with $J_a\subseteq L$ for some $a$, a
finite direct sum of closed summands; the projection onto it is continuous, so
$I_O$ is closed and the quotient map is continuously linearly split.  It is a
sum of two-sided ideals each of which is stable under $\beta_w$.

(ii)  The differential squares to zero because deleting two indices in the two
possible orders gives opposite signs.  The inclusions respect the
decomposition \eqref{eq:Ldecomp}, so the complex is the direct sum of its
parts on the summands $C_L$ and on the scalar summand.  Put
$\mathcal A_L:=\{a:J_a\subseteq L\}$.  The part of $P_O^{-n}$ on $C_L$ is the
sum of copies of $C_L$ indexed by the $n$-element subsets of $\mathcal A_L$.
If $\mathcal A_L\ne\emptyset$, then $C_L\subseteq I_O$ and the part on $C_L$ is
$C_L$ tensored with the augmented simplicial chain complex of the full simplex
on the vertex set $\mathcal A_L$, which is contracted by insertion of the
least vertex.  If $\mathcal A_L=\emptyset$, the part on $C_L$ is the identity
map from $P_O^{0}$ to $B_O$, and the same holds on the scalar summand.  These
contractions are continuous and are combined by the continuous projections of
\eqref{eq:Ldecomp}.  Every term of negative degree is relatively projective on
either side by Lemma~\ref{lem:projective}, and $P_O^{0}=U$ is free.  The image
of $P_O^{-1}$ in $U$ is $I_O$, which gives the last assertion.

(iii)  Choose $a\in T$.  Since $J_a\subseteq J_T$, the map $h_{J_T}$ of
Lemma~\ref{lem:projective} takes values in
$I_{J_T}\widehat\otimes I_{J_a}$ and is a continuous linear right inverse of
multiplication $I_{J_T}\widehat\otimes I_{J_a}\to I_{J_T}$.  Hence the map
\[
I_{J_T}\widehat\otimes I_{J_a}\widehat\otimes B_O\longrightarrow
I_{J_T}\widehat\otimes_UB_O,\qquad y\otimes z\otimes b\mapsto yz\otimes b,
\]
is surjective.  It vanishes on elementary tensors, because
$yz\otimes b=y\otimes zb$ and $I_{J_a}$ annihilates $B_O$; being continuous, it
vanishes identically.  So $I_{J_T}\widehat\otimes_UB_O=0$.  The other identity
is proved in the same way with $h_{J_T}^{\mathrm{op}}$.

(iv)  By (ii), $P_O$ is a resolution of $B_O$ by relatively projective right
modules, so $B_O\widehat\otimes^{\mathbf L}_UB_O$ is represented by
$P_O\widehat\otimes_UB_O$.  By (iii) every term of negative degree vanishes, and
the term of degree zero is $U\widehat\otimes_UB_O=B_O$.
\end{proof}

For an automorphism $\beta$ of $U$ with $\beta(I_O)=I_O$ we write
${}_{\beta}B_O$ for $B_O$ with the left action $y\cdot b:=\beta(y)b$.

\begin{lemma}\label{lem:excvanish}
Let $T\ne\emptyset$, $K\subseteq S_f$ and $w\in A_S^{\times}$.  Then
$I_{K\cup J_T}\,\widehat\otimes_U\,{}_{\beta_w}B_O=0$.
\end{lemma}

\begin{proof}
Choose $a\in T$.  Since $J_a\subseteq K\cup J_T$, the map $h_{K\cup J_T}$
takes values in $I_{K\cup J_T}\widehat\otimes I_{J_a}$ and is a continuous
linear right inverse of multiplication.  Hence
$y\otimes z\otimes b\mapsto yz\otimes b$ maps
$I_{K\cup J_T}\widehat\otimes I_{J_a}\widehat\otimes B_O$ onto
$I_{K\cup J_T}\widehat\otimes_U{}_{\beta_w}B_O$.  On elementary tensors
$yz\otimes b=y\otimes\beta_w(z)b=0$, because
$\beta_w(z)\in I_{J_a}\subseteq I_O$; so the map vanishes.
\end{proof}

\subsubsection*{Restriction to a closed stratum}

For a bounded-above complex $M$ of $U$-bimodules put
\[
\mathcal R_O(M):=B_O\,\widehat\otimes_U\,Q(M)\,\widehat\otimes_U\,B_O ,
\]
a bounded-above complex of relatively projective $B_O$-bimodules, and for
complexes $X,Y$ of $B_O$-bimodules write
$X\star_{B_O}Y:=X\widehat\otimes^{\mathbf L}_{B_O}Y$.  The functor
$\mathcal R_O$ is the two-sided derived restriction to the complement of $O$;
it is triangulated.  For a complex $Z$ of $U$-bimodules put also
\[
\Phi_{M,M'}(Z):=B_O\,\widehat\otimes_U\,Q(M)\,\widehat\otimes_U\,Z\,
\widehat\otimes_U\,Q(M')\,\widehat\otimes_U\,B_O ,\qquad
\mathcal E_O(M,M'):=\Phi_{M,M'}(I_O).
\]

\begin{theorem}\label{thm:exctriangle}
For bounded-above complexes $M,M'$ of $U$-bimodules there is a distinguished
triangle
\begin{equation}\label{eq:exctriangle}
\mathcal E_O(M,M')\longrightarrow \mathcal R_O(M\star M')\longrightarrow
\mathcal R_O(M)\star_{B_O}\mathcal R_O(M')\longrightarrow\mathcal E_O(M,M')[1]
\end{equation}
in $D^{-}_{\mathrm{rel}}(B_O\text{-}B_O)$.  In particular
$\mathcal R_O(M\star M')\to\mathcal R_O(M)\star_{B_O}\mathcal R_O(M')$ is an
isomorphism if and only if $\mathcal E_O(M,M')\cong0$.
\end{theorem}

\begin{proof}
A term of $Q(M)$ is a retract of a free bimodule $U\widehat\otimes V\widehat\otimes
U$, and a term of $Q(M')$ is a retract of some $U\widehat\otimes V'\widehat\otimes
U$.  For a $U$-bimodule $Z$,
\begin{multline*}
B_O\widehat\otimes_U(U\widehat\otimes V\widehat\otimes U)\widehat\otimes_UZ
\widehat\otimes_U(U\widehat\otimes V'\widehat\otimes U)\widehat\otimes_UB_O\\
=B_O\widehat\otimes V\widehat\otimes Z\widehat\otimes V'\widehat\otimes B_O ,
\end{multline*}
which carries admissible sequences in $Z$ to admissible sequences.  Applying $\Phi_{M,M'}$ to the admissible sequence $0\to I_O\to U\to B_O\to0$ of
Lemma~\ref{lem:cech}(i) therefore gives an admissible sequence of complexes of
$B_O$-bimodules, hence a distinguished triangle.  Its first term is
$\mathcal E_O(M,M')$.  Its middle term is
$B_O\widehat\otimes_U\bigl(Q(M)\widehat\otimes_UQ(M')\bigr)\widehat\otimes_UB_O$,
and $Q(M)\widehat\otimes_UQ(M')$ is a bounded-above complex of relatively
projective bimodules representing $M\star M'$; so the middle term is
$\mathcal R_O(M\star M')$.  Its third term is
\[
B_O\widehat\otimes_UQ(M)\widehat\otimes_UB_O\widehat\otimes_UQ(M')\widehat\otimes_UB_O
=\mathcal R_O(M)\,\widehat\otimes_{B_O}\,\mathcal R_O(M'),
\]
which represents $\mathcal R_O(M)\star_{B_O}\mathcal R_O(M')$ because the terms
of $\mathcal R_O(M)$ are relatively projective $B_O$-bimodules.
\end{proof}

\begin{theorem}[Two-sided excision]\label{thm:excision}
Let $\mathcal T$ be the thick subcategory of $D^{-}_{\mathrm{rel}}(U\text{-}U)$
generated by the graph bimodules $E_a$, $a\in A_S^{\times}$, and the ideals
$I_K$, $K\subseteq S_f$.  For $M,M'\in\mathcal T$,
\[
\mathcal R_O(M\star M')\cong\mathcal R_O(M)\star_{B_O}\mathcal R_O(M'),
\]
and $\mathcal R_O(\Delta)\cong B_O$, the unit of $\star_{B_O}$.
\end{theorem}

\begin{proof}
By Theorem~\ref{thm:exctriangle} we must show $\mathcal E_O(M,M')\cong0$.  The
functor $\mathcal E_O$ is triangulated in each variable, so it suffices to
take for $M$ and $M'$ generators of $\mathcal T$.  By the computation in the
proof of Theorem~\ref{thm:exctriangle}, the functor $\Phi_{M,M'}$ preserves admissible exactness, so by Lemma~\ref{lem:cech}(ii) we may replace
$I_O$ by the terms $I_{J_T}$, $T\ne\emptyset$, of its resolution: it suffices
to show that $\Phi_{M,M'}(I_{J_T})$ is admissibly exact.  This complex represents
$B_O\widehat\otimes^{\mathbf L}_U\bigl(X\widehat\otimes^{\mathbf L}_UB_O\bigr)$
with $X:=M\star I_{J_T}\star M'$.  We compute $X$ for the four kinds of pairs
of generators, using that the generators and the ideals are relatively
projective on each side, so that the products are represented by ordinary
tensor products.  If $M=E_u$ then $M\widehat\otimes_UI_{J_T}$ is $I_{J_T}$ with
the left action twisted by $\beta_u^{-1}$, by $x\otimes m=1\otimes\beta_u^{-1}(x)m$;
if $M=I_K$ it is $I_{K\cup J_T}$ by Lemma~\ref{lem:idealproduct}.  If
$M'=E_v$ then $Y\widehat\otimes_UE_v$ is $Y$ with the right action twisted by
$\beta_v$, by $m\otimes x=mx\otimes1$; if $M'=I_L$ then
$I_{K'}\widehat\otimes_UI_L=I_{K'\cup L}$.  In every case $X$ is, as a right
module, an ideal $I_{K'}$ with $K'\supseteq J_T$ whose right action is twisted
by $\beta_w$, with $w=v$ or $w=1$.  Such a module is relatively projective on
the right, because $\beta_w$ is an automorphism, so
\[
X\,\widehat\otimes^{\mathbf L}_U\,B_O\cong X\,\widehat\otimes_U\,B_O
\cong I_{K'}\,\widehat\otimes_U\,{}_{\beta_w^{-1}}B_O=0
\]
by Lemma~\ref{lem:excvanish}, applied with $K'$ in place of $K$ and $w^{-1}$ in
place of $w$.

Finally $\mathcal R_O(\Delta)=B_O\widehat\otimes^{\mathbf L}_UB_O\cong B_O$ by
Lemma~\ref{lem:cech}(iv).
\end{proof}

\begin{lemma}\label{lem:nested}
Let $O$ and $O'$ be given by lists $(J_a)$ and $(J'_b)$ with
$I_O\subseteq I_{O'}$.  Then
$B_O\widehat\otimes^{\mathbf L}_UB_{O'}\cong B_{O'}\cong
B_{O'}\widehat\otimes^{\mathbf L}_UB_O$, and
$\mathcal R_{O'}(M)\cong B_{O'}\widehat\otimes^{\mathbf L}_{B_O}\mathcal R_O(M)
\widehat\otimes^{\mathbf L}_{B_O}B_{O'}$ for every bounded-above complex $M$ of
$U$-bimodules.
\end{lemma}

\begin{proof}
Resolve $B_O$ by the complex $P_O$ of Lemma~\ref{lem:cech}(ii).  For
$T\ne\emptyset$ and $a\in T$ the ideal $I_{J_a}\subseteq I_{O'}$ annihilates
$B_{O'}$, so the argument of Lemma~\ref{lem:cech}(iii) gives
$I_{J_T}\widehat\otimes_UB_{O'}=0=B_{O'}\widehat\otimes_UI_{J_T}$.  Only the term
$U\widehat\otimes_UB_{O'}=B_{O'}$ survives, on either side.  For the last
assertion, $\mathcal R_O(M)$ consists of relatively projective
$B_O$-bimodules, so the right-hand side is represented by
$B_{O'}\widehat\otimes_{B_O}B_O\widehat\otimes_UQ(M)\widehat\otimes_UB_O
\widehat\otimes_{B_O}B_{O'}=B_{O'}\widehat\otimes_UQ(M)\widehat\otimes_UB_{O'}$.
\end{proof}

\section{The involution}\label{sec:dagger}

The involution is defined on the category $G$ and not on a larger category
containing the coefficient ideal $N$, for the following reason.  We have
\begin{equation}\label{eq:HomUN}
\Hom_{U\text{-}U}(U,N)=0 .
\end{equation}
Indeed, a bimodule map $U\to U$ is right multiplication by an element $z$
satisfying \eqref{eq:intertw} with $a=b=1$; by Proposition~\ref{prop:hom}, $z$
is a scalar.  A nonzero scalar does not lie in $N$, so a bimodule map $U\to N$
is zero.

\begin{remark}\label{prop:obstruction}
Let $\mathcal C$ be a category of $U$-bimodules, with bimodule maps as
morphisms, which contains $U$ and a nonzero ideal $I\subseteq N$.  There is no
faithful contravariant functor $D:\mathcal C\to\mathcal C$ with $D(U)\cong U$
and $D(I)\cong I$: it would send the inclusion $I\hookrightarrow U$ to a
nonzero map $U\to I\subseteq N$, contradicting \eqref{eq:HomUN}.  The ideal $N$
is not an object of $G$, and on $G$ only $\Delta^{\dagger}\cong\Delta$ is
required.
\end{remark}

\begin{definition}\label{def:dagger}
For $a\in A_S^{\times}$ put $E_a^{\dagger}:=E_{a^{-1}}$, and for a scalar
morphism $\lambda:E_a\to E_a$ put
$\lambda^{\dagger}:=\bar\lambda:E_a^{\dagger}\to E_a^{\dagger}$.  For an object
$K$ of $G$ let $K^{\dagger}$ be the complex with
$(K^{\dagger})^{n}:=(K^{-n})^{\dagger}$, the direct sum of the daggers of the
summands, and with differential
$d_{K^{\dagger}}^{\,n}:=\bigl(d_K^{-n-1}\bigr)^{\dagger}$, where the dagger of
a matrix of scalar morphisms is its conjugate transpose.  For
$f=(f^{i})_i\in\Hom^{n}(K,L)$, $f^{i}:K^{i}\to L^{i+n}$, let
$f^{\dagger}\in\Hom^{n}(L^{\dagger},K^{\dagger})$ be the family with components
\[
(f^{\dagger})^{i}:=(-1)^{n(n+1)/2}\,\bigl(f^{-i-n}\bigr)^{\dagger}:
(L^{-i})^{\dagger}\longrightarrow(K^{-i-n})^{\dagger}.
\]
\end{definition}

For $a,b\in A_S^{\times}$ with $a=\gamma b$, the morphism
$\lambda U_{\gamma}:E_a\to E_b$ of \eqref{eq:homclass} has the dagger
$\bar\lambda\,U_{\gamma}:E_{b^{-1}}\to E_{a^{-1}}$, which is a morphism of
\eqref{eq:homclass} because $b^{-1}=\gamma a^{-1}$.

\begin{theorem}\label{thm:daggerthm}
The rule of Definition~\ref{def:dagger} is a conjugate-linear contravariant
differential graded functor $\dagger:G\to G$ with
$\dagger\circ\dagger=\id$, with $\Delta^{\dagger}=\Delta$, and with natural
isomorphisms $(K\star L)^{\dagger}\cong L^{\dagger}\star K^{\dagger}$.
\end{theorem}

\begin{proof}
The labels form the group $\ker\varphi$, so $E_{a^{-1}}$ is a generator.  On
the scalar morphisms of generators, $\lambda\mapsto\bar\lambda$ is
conjugate-linear, involutive, and compatible with composition, composition
being multiplication of scalars.  Conjugate transposition of scalar matrices
reverses the order of composition, so $\dagger$ is a contravariant
conjugate-linear operation on matrices.  The differential of $K^{\dagger}$ squares to
zero and $K^{\dagger\dagger}=K$.  For $f\in\Hom^{n}(K,L)$ and
$g\in\Hom^{m}(L,L')$ a direct computation with the components gives
\[
d(f^{\dagger})=d(f)^{\dagger},\qquad
(g\circ f)^{\dagger}=(-1)^{mn}f^{\dagger}\circ g^{\dagger},\qquad
f^{\dagger\dagger}=f:
\]
without the sign $\epsilon_n:=(-1)^{n(n+1)/2}$ the first identity holds up to
the factor $(-1)^{n+1}$ and the second without sign, and
$\epsilon_{n+1}=(-1)^{n+1}\epsilon_n$,
$\epsilon_{m+n}=(-1)^{mn}\epsilon_m\epsilon_n$, $\epsilon_n^{2}=1$.  So
$\dagger$ is a contravariant conjugate-linear differential graded functor with
$\dagger\circ\dagger=\id$.  Since $1^{-1}=1$,
$\Delta^{\dagger}=E_1=\Delta$.  On generators,
$(E_a\star E_b)^{\dagger}=E_{(ab)^{-1}}=E_{b^{-1}a^{-1}}=E_b^{\dagger}\star
E_a^{\dagger}$ by \eqref{eq:starmult}, and on scalar morphisms both sides give
$\overline{\lambda\mu}$.  For complexes $K$, $L$ the summands of
$(K\star L)^{\dagger}$ and of $L^{\dagger}\star K^{\dagger}$ in each degree
therefore coincide, and the identity maps of these summands, multiplied by the
Koszul sign $(-1)^{ij}$ on the summand coming from $K^{i}$ and $L^{j}$, form
the natural isomorphism; this is the isomorphism
$(V\otimes W)^{*}\cong W^{*}\otimes V^{*}$ for bounded complexes of
finite-dimensional vector spaces, under the equivalence of
Remark~\ref{rem:semisimple}.
\end{proof}

By Corollary~\ref{cor:trivial} the involution of smeared objects corresponds
to $F\mapsto F^{\dagger}$ of Definition~\ref{def:DA}; it has the same form as
the involution of $G$, namely inversion of the label and conjugation of the
coefficient.

\subsubsection*{The realization on $L^{2}(A_S)$}
For $a\in A_S^{\times}$ the operator
\[
(\Lambda_af)(x):=|a|_S^{-1/2}f(a^{-1}x)
\]
on $L^{2}(A_S,\mu)$ is unitary, because $d\mu(ax)=|a|_S\,d\mu(x)$, and
$\Lambda_a\Lambda_b=\Lambda_{ab}$.  By Lemma~\ref{lem:prodformula},
$(\Lambda_{\gamma}f)(x)=f(\gamma^{-1}x)$ for $\gamma\in\Gamma_S$.  For
$b=\sum_{\gamma}b_{\gamma}U_{\gamma}\in U$ put
$\pi(b):=\sum_{\gamma}b_{\gamma}\Lambda_{\gamma}$, the coefficient $b_{\gamma}$
acting by multiplication.  The sum converges in operator norm, since
$\sum_{\gamma}\sup_D|b_{\gamma}|$ is a continuous seminorm on $U$ by
Lemma~\ref{lem:functions}, and $\pi$ is a continuous $*$-representation of $U$:
$\Lambda_{\gamma}c\Lambda_{\gamma}^{-1}=\alpha_{\gamma}(c)$ for a multiplication
operator $c$, and the adjoint of $b_{\gamma}\Lambda_{\gamma}$ is
$\alpha_{\gamma}^{-1}(\overline{b_{\gamma}})\Lambda_{\gamma^{-1}}$.  Moreover
$\Lambda_a\pi(c)\Lambda_a^{-1}=\pi(\beta_a(c))$ for $c\in U$ and
$a\in A_S^{\times}$.  With $T_{b,a}f:=\pi(b)\bigl(f(a^{-1}\,\cdot\,)\bigr)$, the
\emph{graph operators} are
\begin{equation}\label{eq:graphops}
\Pi(b,a):=\pi(b)\Lambda_a=|a|_S^{-1/2}T_{b,a},\quad
\Pi(b,a)\Pi(c,d)=\Pi(b\beta_a(c),ad),\quad
\Pi(b,a)^{*}=\Pi(\beta_a^{-1}(b^{*}),a^{-1}),
\end{equation}
for $b,c\in U$ and $a,d\in A_S^{\times}$; the second identity follows from
$\Lambda_a\pi(c)=\pi(\beta_a(c))\Lambda_a$ and the third from
$\Lambda_a^{*}=\Lambda_{a^{-1}}$.  In particular
$\Pi(1,u)\Pi(1,v)=\Pi(1,uv)$ and $\Pi(1,a)^{*}=\Pi(1,a^{-1})$, which are the
operator forms of \eqref{eq:starmult} and of $E_a^{\dagger}=E_{a^{-1}}$.
\section{The Hochschild shadow of a graph bimodule}\label{sec:shadow}

In this section we compute the relative Hochschild shadow \eqref{eq:shadow} of
a graph bimodule $E_w$.  The complex splits over the lattice labels
$\gamma\in\Gamma_S$, and the piece of label $\gamma$ is governed by the twisted
Hochschild complexes of the local algebras $E_v$ for the dilation $\gamma w$.
At a finite place the twisted complex is homotopy equivalent to $\C[0]$ for
every dilation other than the identity, and at the real place it is
contractible.  The shadow therefore does not determine the dilation, and the
pairing of Section~\ref{sec:trace} is defined analytically.

\subsection{Conventions and homotopy lemmas}

All complexes in this section are complexes of complete locally convex spaces
with continuous differentials, concentrated in degrees $n\ge0$.  Chain maps and
homotopies are continuous and linear.  A complex is \emph{contractible} if its
identity is homotopic to zero, and a chain map is a \emph{homotopy
equivalence} if it has an inverse up to homotopy.  A subcomplex $C'\subseteq
C$ is \emph{graded-split} if each $C'_n$ has a closed complement in $C_n$.

For a complete locally convex algebra $A$ with jointly continuous product, not
necessarily unital, and a complete $A$-bimodule $M$, the Hochschild complex is
$C_n(A,M)=M\widehat\otimes A^{\widehat\otimes n}$ with $b=\sum_{i=0}^{n}(-1)^{i}d_i$,
\[
d_0(m\otimes a_1\otimes\cdots)=m\cdot a_1\otimes a_2\otimes\cdots,\qquad
d_n(m\otimes a_1\otimes\cdots\otimes a_n)=a_n\cdot m\otimes a_1\otimes\cdots\otimes a_{n-1},
\]
and $d_i$ multiplying $a_i$ and $a_{i+1}$ for $0<i<n$.  For a continuous
endomorphism $\rho$ of $A$ we write $A_{\rho}$ for $A$ with the bimodule
structure $a\cdot m\cdot b=\rho(a)mb$, and ${}_{1}A_{\rho}$ for $A$ with
$a\cdot m\cdot b=am\rho(b)$, as in \eqref{eq:graphbimod}.  In the notation
of Convention~\ref{conv:twist} the bimodule $A_{\rho}$ is ${}_{\rho}A_{1}$.  If $A$ is
commutative, the reversal
\begin{equation}\label{eq:reversal}
m\otimes a_1\otimes\cdots\otimes a_n\ \longmapsto\ (-1)^{n(n+1)/2}\,m\otimes
a_n\otimes\cdots\otimes a_1
\end{equation}
is an isomorphism of complexes $C(A,{}_{1}A_{\rho})\to C(A,A_{\rho})$: it
exchanges the faces $d_i$ and $d_{n-i}$, and the sign changes by $(-1)^{n}$
from degree $n-1$ to degree $n$.  It commutes with every automorphism of $A$
that commutes with $\rho$, applied to all letters.

If $A$ and $M$ are unital, $C(A,M)$ is a simplicial complete space, with
degeneracies inserting the unit, and the projection of $C(A,M)$ onto the
normalized complex $\bar C(A,M)$, $\bar C_n=M\widehat\otimes(A/\C1)^{\widehat\otimes
n}$, is a homotopy equivalence.  The homotopy inverse and the homotopies of
the normalization theorem \cite[Chapter~VIII, \S6]{MacLane} are finite integral
combinations of composites of faces and degeneracies, so they are continuous,
and they preserve every family of closed subspaces stable under the faces and
degeneracies.

At a finite place $p$ we write $1_j$ for the indicator function of the shell
$p^{j}\Z_p^{\times}$, $e_{j,\kappa}(x)=1_j(x)\kappa(p^{-j}x)$ for the basis in
which the coefficients $f_{j,\kappa}$ of \eqref{eq:seminorms} are taken,
$F^{(j)}:=1_jF_p$ for the closed span of the $e_{j,\kappa}$ with $j$ fixed,
$F_p^{\ge0}$ for the closed span of the $F^{(j)}$ with $j\ge0$, and
$A_*:=\C\chi_p\oplus F_p^{\ge0}=\chi_pE_p$.  Each $F^{(j)}$ is a unital
commutative algebra with unit $1_j$, since $e_{j,\kappa}e_{j,\lambda}=e_{j,\kappa\lambda}$
and $\cond_p(\kappa\lambda)\le\max(\cond_p(\kappa),\cond_p(\lambda))$; one has
$v_{l,N}(1_j)=p^{l|j|}$ and $v_{l,N}(1_j)\,v_{l,N}(1_jf)=v_{2l,N}(1_jf)$.
The algebra $A_*$ is unital with unit $\chi_p$, and
$E_p=A_*\oplus\widehat\bigoplus_{j<0}F^{(j)}$.

\begin{lemma}\label{lem:split}
Let $C'\subseteq C$ be a graded-split subcomplex and $C'':=C/C'$.
\begin{enumerate}
\item[(a)] If $C''$ is contractible, the inclusion $C'\to C$ is a homotopy
equivalence.
\item[(b)] If $C'$ and $C''$ are contractible, so is $C$, and a contraction of
$C$ can be chosen to extend a given contraction of $C'$.
\item[(c)] If $C$ and $C''$ are contractible, so is $C'$.
\item[(d)] If $C$ has an increasing filtration by graded-split subcomplexes
$0=F_{-1}\subseteq F_0\subseteq F_1\subseteq\cdots$ which in each degree $n$
satisfies $(F_r)_n=C_n$ for $r$ large, and every $F_r/F_{r-1}$ is contractible,
then $C$ is contractible.
\end{enumerate}
\end{lemma}

\begin{proof}
Choose complements, so that $C_n=C'_n\oplus C''_n$ and
$d(x',x'')=(d'x'+\tau x'',d''x'')$ with $\tau:C''_n\to C'_{n-1}$ continuous;
$d^{2}=0$ gives $d'\tau+\tau d''=0$.

(a) Let $h''$ be a contraction of $C''$ and put $r(x',x''):=x'-\tau h''x''$.
Then $d'r(x',x'')=d'x'+\tau d''h''x''$ and
$r\,d(x',x'')=d'x'+\tau x''-\tau h''d''x''$, which agree because
$d''h''+h''d''=\id$; so $r$ is a chain map, and $r$ restricts to the identity
on $C'$.  With $H(x',x''):=(0,h''x'')$ one finds
$(dH+Hd)(x',x'')=(\tau h''x'',x'')=(x',x'')-r(x',x'')$.

(b) Let $h'$ contract $C'$ and put $h(x',x''):=(h'x'-h'\tau h''x'',h''x'')$.
The $C''$-component of $(dh+hd)(x',x'')$ is $x''$, and its $C'$-component is
\[
x'+\bigl(-d'h'\tau h''+\tau h''+h'\tau-h'\tau h''d''\bigr)x''
=x'+h'\tau\bigl(\id-d''h''-h''d''\bigr)x''=x',
\]
where we used $d'h'=\id-h'd'$ and $d'\tau=-\tau d''$.

(c) By (a) the inclusion $C'\to C$ is a homotopy equivalence, and a complex
homotopy equivalent to a contractible one is contractible.

(d) By (b) and induction on $r$ there are contractions $h_r$ of $F_r$ with
$h_r|_{F_{r-1}}=h_{r-1}$.  In degree $n$ they are independent of $r$ for $r$
large, and they define a contraction of $C$.
\end{proof}

\begin{lemma}[Central homotopy]\label{lem:central}
Let $A$ be a complete locally convex algebra, $M$ a complete $A$-bimodule and
$z$ a central element of $A$.  Put
\[
h_z(m\otimes a_1\otimes\cdots\otimes a_n):=\sum_{i=0}^{n}(-1)^{i}\,m\otimes
a_1\otimes\cdots\otimes a_i\otimes z\otimes a_{i+1}\otimes\cdots\otimes a_n .
\]
Then $bh_z+h_zb=R_z-L_z$ on $C(A,M)$, where $R_z$ and $L_z$ replace the
coefficient $m$ by $m\cdot z$ and by $z\cdot m$.  In particular, if $A$ is
unital and commutative, $\rho$ is a continuous endomorphism of $A$, and
$z-\rho(z)$ is invertible in $A$ for some $z\in A$, then $C(A,A_{\rho})$ is
contractible, by the contraction $G\,h_z$, where $G$ multiplies the coefficient
by $(z-\rho(z))^{-1}$.
\end{lemma}

\begin{proof}
Let $t_i$ insert $z$ after the $i$-th slot, $0\le i\le n$, so that
$h_z=\sum_i(-1)^{i}t_i$.  Then $d_jt_i=t_{i-1}d_j$ for $j<i$ and
$d_jt_i=t_id_{j-1}$ for $j>i+1$, while $d_it_i$ multiplies the $i$-th letter
on the right by $z$ and $d_{i+1}t_i$ multiplies the $(i+1)$-st letter on the
left by $z$; for $i=n$ the latter is $z\cdot m$, and for $i=0$ the former is
$m\cdot z$.  For $1\le i\le n$ the terms $d_it_i$ and $d_it_{i-1}$ are equal,
because $z$ is central, and they occur in $bh_z$ with opposite signs.  The
terms with $j<i$ or $j>i+1$ sum to $-h_zb$.  What remains is
$d_0t_0-d_{n+1}t_n=R_z-L_z$.

For $M=A_{\rho}$ with $A$ commutative, $R_z-L_z$ multiplies the coefficient by
$g:=z-\rho(z)$.  It is a chain map, being $bh_z+h_zb$, and it is bijective with
inverse $G$; so $G$ is a chain map and $b(Gh_z)+(Gh_z)b=G(R_z-L_z)=\id$.
\end{proof}

\subsection{Excision at the origin and the K\"unneth decomposition}

For a complete locally convex algebra $I$ let $B'_n(I)=I^{\widehat\otimes n}$,
$n\ge1$, with
$b'(a_1\otimes\cdots\otimes a_n)=\sum_{j=1}^{n-1}(-1)^{j-1}a_1\otimes\cdots\otimes
a_ja_{j+1}\otimes\cdots\otimes a_n$.  A \emph{right splitting} of $I$ is a
continuous linear map $h:I\to I\widehat\otimes I$ with $\mu\circ h=\id_I$,
$\mu$ the multiplication, and $h(ab)=h(a)b$ for $a,b\in I$.  If $h$ is a right
splitting and $Y$ is a complete space, then $s:=h\otimes\id$ satisfies
\begin{equation}\label{eq:barcontraction}
b's+sb'=\id\quad\text{on }B'(I)\widehat\otimes Y .
\end{equation}
Indeed, on $I\widehat\otimes I^{\widehat\otimes m}$ the differential is
$\mu\otimes\id-\id_I\otimes b'$, so $b's=\mu h\otimes\id-(\id_I\otimes b')s$ on
$I^{\widehat\otimes m}$.  For $m\ge2$,
\[
(\id_I\otimes b')s(a_0\otimes a_1\otimes\cdots)=h(a_0)a_1\otimes a_2\otimes\cdots
-h(a_0)\otimes b'(a_1\otimes\cdots)=sb'(a_0\otimes a_1\otimes\cdots),
\]
because $h(a_0)a_1=h(a_0a_1)$; hence $b's=\id-sb'$.  For $m=1$ one has
$b's=\mu h=\id$.

\begin{lemma}[Right splittings]\label{lem:hunital}
Each of the following algebras has a right splitting $h$, and therefore a
contractible complex $B'$:
\begin{enumerate}
\item[(i)] $E_V:=\widehat\bigotimes_{v\in V}E_v$ for every nonempty
$V\subseteq S$, where $E_{\infty}:=E_{\R}$; in particular $E_P=E_S$;
\item[(ii)] every ideal $I_J$, in particular $N=I_{\emptyset}$;
\item[(iii)] $F_p$ and $F_p^{\ge0}$ for every $p\in S_f$.
\end{enumerate}
In each case $h(a)=\sum_ie_i\otimes\rho_ia$ for families $e_i,\rho_i$ in the
algebra, and if the algebra is an ideal in a complete locally convex algebra
$A$, then $h(ab)=h(a)b$ for all $b\in A$.
\end{lemma}

\begin{proof}
(ii) The map $h_J(a)=\sum_ie_i\otimes\rho_ia$ of the proof of
Lemma~\ref{lem:projective}, with the families of Lemma~\ref{lem:partition},
is a right splitting.

(i) Run the construction of Lemma~\ref{lem:partition} at the places of $V$
only: the functions $\psi_n,\varepsilon_n$ on both half-lines if $\infty\in V$,
and $\chi_p$ together with the shell idempotents $1_j$, $j<0$, at each finite
$p\in V$.  This gives families $e_i,\rho_i\in E_V$ with $e_i\rho_i=\rho_i$ and
$\sum_i\rho_i=1$ pointwise, the sum being locally finite.  The estimate
\eqref{eq:summable} holds for them, because its proof is a product over the
places of one-place estimates (the window weight $e^{-|n|}$ or $p^{-|j|}$
against the size $e^{l|n|}$ or $p^{l|j|}$ of the cutoff) and involves neither
the lattice weights nor the places outside $V$.  Hence
$h(a)=\sum_ie_i\otimes\rho_ia$ converges absolutely in $E_V\widehat\otimes E_V$,
is continuous, and $\mu h(a)=\sum_i\rho_ia=a$.

(iii) Take $e_j=\rho_j=1_j$, with $j\in\Z$ for $F_p$ and $j\ge0$ for
$F_p^{\ge0}$.  For seminorms $v_{l,N}$ and $v_{l',N'}$,
\[
\sum_jv_{l,N}(1_j)\,v_{l',N'}(1_ja)=\sum_jp^{l|j|}v_{l',N'}(1_ja)=v_{l+l',N'}(a),
\]
so $h(a)=\sum_j1_j\otimes1_ja$ converges absolutely and is continuous, and
$\mu h(a)=\sum_j1_ja=a$.

In all cases right linearity is associativity, $\rho_i(ab)=(\rho_ia)b$, applied
termwise and extended by continuity; it holds for $b$ in any algebra $A$
containing the given one as an ideal.  The contraction of $B'$ is
\eqref{eq:barcontraction}.
\end{proof}

\begin{theorem}[Split augmented excision]\label{thm:B2}
Let $R=\C1\oplus I$ be a unital complete locally convex algebra with a closed
ideal $I$ that has a right splitting, and let $\sigma$ be a continuous unital
endomorphism of $R$ with $\sigma(I)\subseteq I$.  Let $C^{\mathrm{sc}}\subseteq
C(R,R_{\sigma})$ be the span of the words all of whose letters are scalar
multiples of $1$, and $C^{+}(R,R_{\sigma})$ the closed span of the words with
at least one letter in $I$.  Then
\[
C(R,R_{\sigma})=C^{\mathrm{sc}}\oplus C^{+}(R,R_{\sigma})
\]
as complexes, $C^{\mathrm{sc}}\cong C(\C,\C)$ is homotopy equivalent to
$\C[0]$, and the inclusion $C(I,I_{\sigma})\to C^{+}(R,R_{\sigma})$ is a
homotopy equivalence.  In particular
$C(R,R_{\sigma})\simeq\C[0]\oplus C(I,I_{\sigma})$, the summand $\C[0]$ being
spanned by the class of the unit.
\end{theorem}

\begin{proof}
Decompose each letter of a word according to $R=\C1\oplus I$.  A face
multiplies two adjacent letters or applies $\sigma$ to the last letter and
multiplies it with the first; since $\sigma(1)=1$, $\sigma(I)\subseteq I$ and
$I$ is an ideal, a word of scalars goes to a word of scalars and a word with a
letter in $I$ to a word with a letter in $I$.  So $C^{\mathrm{sc}}$ and
$C^{+}$ are complementary subcomplexes, both stable under the degeneracies.
The complex $C^{\mathrm{sc}}=C(\C,\C)$ is $\C$ in every degree with
differential alternately $0$ and $\id$, and its normalized complex is $\C[0]$.

By the normalization theorem applied to the simplicial space $C^{+}$, the
projection $C^{+}\to\bar C^{+}$ is a homotopy equivalence, where
$\bar C^{+}_n=P_n\oplus Q_n$ with $P_n:=\C1\otimes I^{\widehat\otimes n}$ for
$n\ge1$, $P_0:=0$, and $Q_n:=I\widehat\otimes I^{\widehat\otimes n}$.  The
boundary is
\[
b(a_0\otimes\cdots\otimes a_n)=\sum_{j<n}(-1)^{j}a_0\otimes\cdots\otimes
a_ja_{j+1}\otimes\cdots+(-1)^{n}\sigma(a_n)a_0\otimes a_1\otimes\cdots\otimes
a_{n-1}.
\]
On $Q$ all products lie in $I$, so $Q$ is a subcomplex and $Q=C(I,I_{\sigma})$.
Applied to $1\otimes a_1\otimes\cdots\otimes a_n\in P_n$, the terms with
$1\le j<n$ lie in $P_{n-1}$ and add up to $-1\otimes b'(a_1\otimes\cdots\otimes
a_n)$, while the term $j=0$ and the last term lie in $Q_{n-1}$.  Hence $Q$ is a
graded-split subcomplex of $\bar C^{+}$ with quotient the complex $B'(I)$, with
differential $-b'$ and degrees shifted by one, which is contractible by
\eqref{eq:barcontraction}.  By Lemma~\ref{lem:split}(a) the inclusion
$Q\to\bar C^{+}$ is a homotopy equivalence.  It is the composite of the
inclusion $C(I,I_{\sigma})\to C^{+}$ with the projection $C^{+}\to\bar C^{+}$,
so the former is a homotopy equivalence as well.
\end{proof}

The following is a twisted form of excision in Hochschild homology.  For
untwisted coefficients, excision for H-unital ideals is due to Wodzicki
\cite{Wodzicki}, and for Fr\'echet algebras to Brodzki and Lykova
\cite{BrodzkiLykova}.

\begin{theorem}[Twisted excision]\label{thm:excision-input}
Let $A$ be a unital complete locally convex algebra, $I\subseteq A$ a closed
two-sided ideal with a closed complement $L\ni1$, and $\sigma$ a continuous
unital endomorphism of $A$ with $\sigma(I)\subseteq I$.  Suppose that $I$ has a
right splitting $h$ with $h(ab)=h(a)b$ for $a\in I$, $b\in A$.  Let
\[
\mathcal D(A,I):=\ker\bigl(C(A,A_{\sigma})\to C(A/I,(A/I)_{\sigma})\bigr).
\]
Then the inclusions
\[
C(I,I_{\sigma})\ \hookrightarrow\ C(A,I_{\sigma})\ \hookrightarrow\ \mathcal D(A,I)
\]
are homotopy equivalences, and the quotient complex
$\mathcal D(A,I)/C(I,I_{\sigma})$ is contractible.
\end{theorem}

\begin{proof}
Write each letter of $A^{\widehat\otimes(n+1)}$ according to $A=L\oplus I$.
Then $A^{\widehat\otimes(n+1)}$ is the finite topological direct sum of the
$2^{n+1}$ \emph{patterns} $X_0\widehat\otimes\cdots\widehat\otimes X_n$,
$X_j\in\{L,I\}$.  The map to $C_n(A/I,(A/I)_{\sigma})\cong L^{\widehat\otimes(n+1)}$
is the projection onto the pattern $L\cdots L$, so $\mathcal D(A,I)_n$ is the
sum of all other patterns; $C_n(A,I_{\sigma})$ is the sum of the patterns with
$X_0=I$, and $C_n(I,I_{\sigma})$ is the pattern $I\cdots I$.  All three are
subcomplexes because $I$ is an ideal and $\sigma(I)\subseteq I$, and each is
graded-split in the next.  After a face, the new letter is decomposed by the
continuous projections $\pi_L,\pi_I$.  By Lemma~\ref{lem:split}(a),(b) it
suffices to show that the two successive quotient complexes are contractible,
and by
Lemma~\ref{lem:split}(d) it suffices to filter each of them by sums of
patterns, finitely many steps in each degree, with contractible graded pieces.

\emph{Step 1: $\mathcal Q_1:=C(A,I_{\sigma})/C(I,I_{\sigma})$.}  It is the sum of
the patterns with $X_0=I$ and some $X_j=L$, $j\ge1$.  Let $m\ge1$ be the first
such $j$ and $q:=n-m+1$.  The faces $d_j$, $j<m-1$, preserve $q$.  The face
$d_{m-1}$ replaces the first $L$ by a letter of $I$, since $I\cdot L\subseteq
I$; the faces $d_j$, $m\le j<n$, shorten the tail; and the last face lowers $q$
if $m<n$ and lands in $C(I,I_{\sigma})$ if $m=n$.  All of these lower $q$ or
give zero in $\mathcal Q_1$.  So $G_s:=\{q\le s\}$ is a filtration by sums of
patterns, with $1\le q\le n$ in degree $n$, and $G_s/G_{s-1}\cong
B'(I)\widehat\otimes Y_s$ with $Y_s=L\widehat\otimes A^{\widehat\otimes(s-1)}$
and differential $b'\otimes\id$ on the head $a_0\otimes\cdots\otimes a_{m-1}$.
It is contractible by \eqref{eq:barcontraction}.

\emph{Step 2: $\mathcal Q_2:=\mathcal D(A,I)/C(A,I_{\sigma})$.}  It is the sum of
the patterns with $X_0=L$ and some $X_j=I$, $j\ge1$.  Let $m\ge1$ be the first
such $j$.  A face $d_j$ with $j\le m-1$ lowers $m$ or sends the word into
$C(A,I_{\sigma})$: a product $l_jl_{j+1}$ has an $L$-part, which shifts the
first $I$ to the left, and an $I$-part, which puts a letter of $I$ at the
position $j<m$; and $l_{m-1}a_m\in I$.  The faces $d_j$, $m\le j<n$, preserve
$m$, and so does the $L$-part of the last face when $m<n$; its $I$-part, and
the whole last face when $m=n$ (then $\sigma(a_n)\in I$), land in
$C(A,I_{\sigma})$.  So $\{m\le r\}$ is a filtration by sums of patterns, with
$1\le m\le n$ in degree $n$, and its $m$-th graded piece has the terms
$L^{\widehat\otimes m}\widehat\otimes I\widehat\otimes A^{\widehat\otimes k}$,
$n=m+k$, $k\ge0$.  Moving the head $X:=L^{\widehat\otimes m}$ to the right
identifies it, up to the global sign $(-1)^{m}$ of the differential, with the
complex $\mathcal B_k=I\widehat\otimes A^{\widehat\otimes k}\widehat\otimes X$,
\[
\partial(a\otimes b_1\otimes\cdots\otimes b_k\otimes x)=ab_1\otimes b_2\otimes\cdots
+\sum_{j=1}^{k-1}(-1)^{j}\cdots\otimes b_jb_{j+1}\otimes\cdots
+(-1)^{k}a\otimes b_1\otimes\cdots\otimes b_{k-1}\otimes b_k\cdot x ,
\]
where $X$ is the unital left $A$-module with
$b\cdot(l_0\otimes y)=\pi_L(\sigma(b)l_0)\otimes y$.  This is a module
structure because $\sigma(b)\pi_I(\,\cdot\,)\in I$, and $I\cdot X=0$ because
$\sigma(I)\subseteq I$.  Let $\mathrm{Bar}_k:=A\widehat\otimes
A^{\widehat\otimes k}\widehat\otimes X$ be the bar resolution of $X$, with
differential $\partial'$, augmentation $\varepsilon:\mathrm{Bar}_0\to X$ given by
the action, and $c(\xi):=1\otimes\xi$; then $\partial'c+c\partial'=\id$ in
degrees $k\ge1$ and $\partial'c=\id-c_{-1}\varepsilon$ on $\mathrm{Bar}_0$, where
$c_{-1}(x)=1\otimes x$.  Define
\[
\Phi:=h\otimes\id:\mathcal B_k\to I\widehat\otimes\mathrm{Bar}_k,\qquad
\Psi:I\widehat\otimes\mathrm{Bar}_k\to\mathcal B_k,\quad
\Psi(a\otimes b_0\otimes\xi)=ab_0\otimes\xi,
\]
where $h$ is regarded as a map $I\to I\widehat\otimes A$.  Both are chain maps
for the differentials $\partial$ and $\id_I\otimes\partial'$: for $\Psi$ this is
associativity, and for $\Phi$ it is $h(ab_1)=h(a)b_1$.  Moreover
$\Psi\Phi=\mu h\otimes\id=\id$.  Put $H:=\Psi(\id_I\otimes c)\Phi$.  In degrees
$k\ge1$, $\partial H+H\partial=\Psi\Phi=\id$.  In degree $0$,
$\partial H=\Psi\Phi-\Psi(\id_I\otimes c_{-1}\varepsilon)\Phi$, and the last term
vanishes: for $h(a)=\sum u\otimes u'$ with $u'\in I$ it sends $a\otimes x$ to
$\sum u\otimes u'\cdot x=0$.  So $H$ is a contraction of $\mathcal B$.
\end{proof}

\begin{theorem}[K\"unneth over places]\label{thm:kunneth}
Let $\sigma=\bigotimes_{v\in S}\sigma_v$ be a continuous endomorphism of $E_P$
acting coordinatewise.  Then there are homotopy equivalences
\begin{equation}\label{eq:kunneth}
C\bigl(E_P,(E_P)_{\sigma}\bigr)\ \simeq\
\mathrm{Tot}\,\widehat\bigotimes_{v\in S}C\bigl(E_v,(E_v)_{\sigma_v}\bigr),
\end{equation}
where $E_{\infty}:=E_{\R}$.  The same holds for $E_V$, $\emptyset\ne
V\subseteq S$.
\end{theorem}

\begin{proof}
We argue by induction on $|V|$, the case $|V|=1$ being empty.  Let $|V|\ge2$,
fix $v\in V$, and put $I:=E_v$, $J:=E_{V\setminus\{v\}}$,
$\sigma':=\bigotimes_{u\ne v}\sigma_u$, $K:=I\widehat\otimes J=E_V$.  Let
$I^{+}=\C1\oplus I$, $J^{+}=\C1\oplus J$ and
\[
T:=I^{+}\widehat\otimes J^{+}=\C\oplus I\oplus J\oplus K ,
\]
with $\sigma$ extended unitally.  Call the \emph{support} of a letter in
$\C,I,J,K$ the set $\emptyset,\{1\},\{2\},\{1,2\}$, and the support of a word
the union of the supports of its letters.  A product of letters has support the
union of their supports or vanishes, and $\sigma$ preserves the four summands;
so $C(T,T_{\sigma})$ is the direct sum of its subcomplexes of given support.

\emph{Step 1.}  The space $C_n(T,T_{\sigma})$ is
$C_n(I^{+},I^{+}_{\sigma_v})\widehat\otimes C_n(J^{+},J^{+}_{\sigma'})$, by
associativity and commutativity of the completed projective tensor product,
and its faces and degeneracies are $d_i\otimes d_i$ and $s_i\otimes s_i$.  The
Eilenberg--Zilber theorem \cite[Chapter~VIII, \S8]{MacLane} provides, for
simplicial modules $X$ and $Y$, a chain homotopy equivalence between the
normalized complex of $(X_n\otimes Y_n)_n$ and the total complex of the tensor
product of the normalized complexes of $X$ and $Y$, natural in $X$ and $Y$.
Natural maps $X_p\otimes Y_q\to X_n\otimes Y_m$ are, by the Yoneda lemma,
finite integral combinations of maps $\theta\otimes\theta'$ with
$\theta,\theta'$ composites of faces and degeneracies, and an identity between
such combinations which holds for all simplicial modules holds formally.  The
equivalence, its inverse and the homotopies are therefore defined, with the
same identities, for simplicial complete spaces and the completed tensor
product, and they are continuous.  Together with the normalization theorem
this gives a homotopy equivalence
\[
C(T,T_{\sigma})\ \simeq\ \mathrm{Tot}\bigl(C(I^{+},I^{+}_{\sigma_v})\widehat\otimes
C(J^{+},J^{+}_{\sigma'})\bigr),
\]
all of whose maps preserve subspaces of the form $X'\widehat\otimes Y'$ with
$X'$, $Y'$ stable under faces and degeneracies.  A word of $T$ has support
$\{1,2\}$ if and only if its two components lie in $C^{+}(I^{+})$ and
$C^{+}(J^{+})$, in the notation of Theorem~\ref{thm:B2}.  Hence, by
Theorem~\ref{thm:B2} and the right splittings of
Lemma~\ref{lem:hunital}(i),
\[
C(T,T_{\sigma})^{(12)}\ \simeq\ \mathrm{Tot}\bigl(C^{+}(I^{+})\widehat\otimes
C^{+}(J^{+})\bigr)\ \simeq\ \mathrm{Tot}\bigl(C(I,I_{\sigma_v})\widehat\otimes
C(J,J_{\sigma'})\bigr).
\]

\emph{Step 2.}  The inclusion $C(K,K_{\sigma})\to C(T,T_{\sigma})^{(12)}$ is a
homotopy equivalence.  Indeed, $K$ is a closed ideal of $T$ with closed
complement $\C\oplus I\oplus J\ni1$, and its right splitting is $T$-linear on
the right by Lemma~\ref{lem:hunital}.  By Theorem~\ref{thm:excision-input},
$C(K,K_{\sigma})\to\mathcal D(T,K)$ is a homotopy equivalence.  The complex
$\mathcal D(T,K)$ is spanned by the words with a letter in $K$, so it is a
graded-split subcomplex of $C(T,T_{\sigma})^{(12)}$, and the quotient is the
support-$\{1,2\}$ part $C(Q',Q'_{\sigma})^{(12)}$ of the complex of
$Q':=T/K=\C\oplus I\oplus J$, in which $IJ=JI=0$.  By
Lemma~\ref{lem:split}(a) it remains to show that
$C(Q',Q'_{\sigma})^{(12)}$ is contractible.

Now $I$ is a closed ideal of $Q'$ with closed complement $J^{+}\ni1$, and its
right splitting is $Q'$-linear on the right.  The complex $\mathcal D(Q',I)$ is
spanned by the words with a letter in $I$; it is the direct sum of its parts
of support $\{1\}$ and $\{1,2\}$, the latter being all of
$C(Q',Q'_{\sigma})^{(12)}$, and $C(I,I_{\sigma})$ lies in the part of support
$\{1\}$.  By Theorem~\ref{thm:excision-input} the quotient
$\mathcal D(Q',I)/C(I,I_{\sigma})$ is contractible.  It contains
$C(Q',Q'_{\sigma})^{(12)}$ as a direct summand complex, and a direct summand of
a contractible complex is contractible.

\emph{Step 3.}  By Steps 1 and 2, $C(K,K_{\sigma})\simeq
\mathrm{Tot}(C(I,I_{\sigma_v})\widehat\otimes C(J,J_{\sigma'}))$.  By the
induction hypothesis $C(J,J_{\sigma'})$ is homotopy equivalent to
$\mathrm{Tot}\,\widehat\bigotimes_{u\ne v}C(E_u,(E_u)_{\sigma_u})$, and the
completed tensor product of a homotopy equivalence with a fixed complex is a
homotopy equivalence.
\end{proof}

\subsection{The local twisted computations}

\begin{lemma}[Block lemma]\label{lem:block}
Let $A$ be a complete locally convex algebra and $(e_a)_{a\in\mathcal I}$ a
family of central, pairwise orthogonal idempotents of $A$ such that
$x=\sum_ae_ax$ for all $x\in A$ and such that for every continuous seminorm
$\nu$ there is a continuous seminorm $\nu'$ with
\begin{equation}\label{eq:blockA}
\sum_a\bigl(1+\nu(e_a)\bigr)\,\nu(e_ax)\le\nu'(x)\qquad(x\in A).
\end{equation}
Let $M$ be a complete $A$-bimodule such that $m=\sum_{a,b}e_a\cdot m\cdot e_b$
for all $m$, and such that for all continuous seminorms $\mu$ on $M$ and $\nu$
on $A$ there is a continuous seminorm $\mu'$ with
\begin{equation}\label{eq:blockM}
\sum_{a,b}\bigl(1+\nu(e_b)\bigr)\,\mu(e_a\cdot m\cdot e_b)\le\mu'(m)\qquad(m\in M).
\end{equation}
Put $A_a:=e_aA$ and $M_{ab}:=e_a\cdot M\cdot e_b$.  Then
\[
C(A,M)\ =\ \widehat\bigoplus_{a}C(A_a,M_{aa})\ \oplus\ X
\]
as complexes, with $X$ contractible.  In particular, if $M_{aa}=0$ for every
$a$, then $C(A,M)$ is contractible.
\end{lemma}

\begin{proof}
By \eqref{eq:blockA} and \eqref{eq:blockM}, applied in each tensor factor,
every chain is an absolutely convergent sum of its \emph{homogeneous}
components, which are limits of combinations of words $w=m\otimes
a_1\otimes\cdots\otimes a_n$ with $m\in M_{ab}$ and $a_r\in A_{k_r}$, and the
sum over all multi-indices of a projective seminorm of the components is
bounded by a projective seminorm of the chain.  Put $k_0:=b$ and
$k_{n+1}:=a$, and say that the edge $r$, $0\le r\le n$, of $w$ is
\emph{matched} if $k_r=k_{r+1}$.  The face $d_r$ multiplies across the edge
$r$, and it vanishes on $w$ if the edge is unmatched, because the product
contains $e_{k_r}e_{k_{r+1}}=0$.  A matched face merges two letters and leaves
the remaining edges, matched or not, unchanged.  Therefore the fully matched
words, those with all indices equal to some $a$, span the subcomplex
$\widehat\bigoplus_aC(A_a,M_{aa})$, and the words with at least one unmatched
edge span a complementary subcomplex $X$.

On $X$ let $r=r(w)$ be the first unmatched edge and define
$s(w):=(-1)^{r}s_{r}(w)$, where $s_r$ inserts $e_{k_r}$ immediately after the
$r$-th slot (after $m$ if $r=0$).  Then $d_rs_rw=w$, because $e_{k_r}$ is a
unit for the $r$-th letter; $d_{r+1}s_rw=0$, because $e_{k_r}$ meets the next
letter (or, for $r=n$, acts on $m$) across an unmatched edge; for $t<r$ one has
$d_ts_r=s_{r-1}d_t$, and $d_tw$ has first unmatched edge $r-1$, with the same
idempotent to insert; and for $t>r+1$ one has $d_ts_r=s_rd_{t-1}$, while
$d_{t-1}w$ is either $0$ or has first unmatched edge $r$.  Also $d_rw=0$.
Summing with the signs, all terms of $bs+sb$ cancel in pairs except
$(-1)^{r}(-1)^{r}d_rs_rw=w$, so $bs+sb=\id$ on $X$ in every degree.  The map
$s$ is continuous: inserting $e_{k_r}$ multiplies a projective seminorm of a
homogeneous component by $\nu(e_{k_r})$, and by \eqref{eq:blockA} and
\eqref{eq:blockM}, applied with the factor $\nu(e_{k_r})$ in the $r$-th slot,
the sum over the components is bounded by $n+1$ times a projective seminorm of
the chain.
\end{proof}

The hypotheses of Lemma~\ref{lem:block} hold for $A=E_p$ with the idempotents
$\chi_p$ and $1_j$, $j<0$, for $A=F_p^{\ge0}$ with the $1_j$, $j\ge0$, and for
$A=F_p$ with all the $1_j$: one has
$\sum_j(1+p^{l|j|})v_{l,N}(1_jx)\le2v_{2l,N}(x)$, and $\chi_p$ is a single
idempotent.  They hold also for the bimodules $A_{\rho}$ below, in which
$\rho(e_a)e_b$ runs over the indicator functions of pairwise disjoint sets,
each a single shell or a set $p^{k}\Z_p$ with $0\le k\le|v|$.

\begin{lemma}[Non-unit twists]\label{lem:nonunit}
Let $\sigma$ be the dilation $f\mapsto f(w^{-1}\,\cdot\,)$ by
$w\in\Q_p^{\times}$ with $v:=v_p(w)\ne0$.  Then
\begin{equation}\label{eq:D1}
C\bigl(E_p,(E_p)_{\sigma}\bigr)\ \simeq\ \C[0],
\end{equation}
the class being that of $\chi_p$.  Moreover $C(F_p,(F_p)_{\sigma})$ is
contractible.
\end{lemma}

\begin{proof}
$\sigma$ maps the shell $j$ to the shell $j+v$, so $\sigma(1_j)=1_{j+v}$ and
$\sigma(\chi_p)=\mathbf 1_{w\Z_p}$.  Apply Lemma~\ref{lem:block} to $A=E_p$ with
the blocks $A_*$ and $F^{(j)}$, $j<0$, and $M=(E_p)_{\sigma}$, for which
$M_{ab}=\sigma(e_a)E_pe_b$.  For a shell block, $M_{jj}=1_{j+v}1_jE_p=0$.  Hence
$C(E_p,M)\simeq C(A_*,M_{**})$ with $M_{**}=\mathbf 1_{w\Z_p}\chi_pE_p$.

If $v<0$, then $\Z_p\subset w\Z_p$ and $M_{**}=A_*$ with $a\cdot m\cdot
b=\sigma'(a)mb$, where $\sigma'(a)(x):=a(w^{-1}x)$ for $x\in\Z_p$.  Since
$w^{-1}\Z_p\subseteq\Z_p$, $\sigma'$ is a continuous unital endomorphism of
$A_*$ which preserves the value at the origin, so
$\sigma'(F_p^{\ge0})\subseteq F_p^{\ge0}$.  If $v>0$, then $M_{**}=\sigma(A_*)$
consists of the functions on $w\Z_p$, and $\psi:A_*\to M_{**}$,
$\psi(c)=\sigma(c)$, is an isomorphism with
$a\cdot\psi(c)\cdot b=\psi\bigl(ac\,\rho(b)\bigr)$, $\rho(b)(y):=b(wy)$.  So
$M_{**}\cong{}_{1}(A_*)_{\rho}$ for the unital endomorphism $\rho$ of $A_*$,
which again preserves the value at the origin, and by \eqref{eq:reversal} its
complex is isomorphic to $C(A_*,(A_*)_{\rho})$.

In both cases let $\rho'$ be $\sigma'$ or $\rho$.  Theorem~\ref{thm:B2} applies
to $R=A_*=\C\chi_p\oplus F_p^{\ge0}$, the ideal $F_p^{\ge0}$ having a right
splitting by Lemma~\ref{lem:hunital}(iii), and gives
$C(A_*,(A_*)_{\rho'})\simeq\C[0]\oplus C(F_p^{\ge0},(F_p^{\ge0})_{\rho'})$, the
summand $\C[0]$ being the class of $\chi_p$.  The endomorphism $\rho'$ sends
$1_j$ to $1_{j-|v|}$, or to $0$ when $j<|v|$, and never to $1_j$; so
Lemma~\ref{lem:block}, applied to $F_p^{\ge0}$ with the idempotents $1_j$,
$j\ge0$, shows that $C(F_p^{\ge0},(F_p^{\ge0})_{\rho'})$ is contractible.  The
last assertion is Lemma~\ref{lem:block} applied to $F_p$ with all the $1_j$,
where $M_{jj}=1_{j+v}1_jF_p=0$.
\end{proof}

\begin{lemma}[Unit twists]\label{lem:unit}
Let $\sigma$ be the dilation by a unit $u\in\Z_p^{\times}$, $u\ne1$.  Then
$C(F^{(j)},(F^{(j)})_{\sigma})$ is contractible for every $j\in\Z$,
$C(F_p,(F_p)_{\sigma})$ is contractible, and
\[
C\bigl(E_p,(E_p)_{\sigma}\bigr)\ \simeq\ \C[0],
\]
the class being that of $\chi_p$.
\end{lemma}

\begin{proof}
Choose a character $\kappa$ of $\Z_p^{\times}$ with $\kappa(u)\ne1$ and put
$z_j:=e_{j,\kappa}\in F^{(j)}$.  Then $\sigma(z_j)=\lambda z_j$ with
$\lambda=\kappa(u)^{-1}\ne1$, and $g_j:=z_j-\sigma(z_j)=(1-\lambda)z_j$ is
invertible in the unital commutative algebra $F^{(j)}$, with inverse
$(1-\lambda)^{-1}e_{j,\kappa^{-1}}$.  By Lemma~\ref{lem:central},
$H_j:=G_jh_{z_j}$ is a contraction of $C(F^{(j)},(F^{(j)})_{\sigma})$, where
$G_j$ multiplies the coefficient by $g_j^{-1}$.  Inserting $z_j$ multiplies a
projective seminorm built from $v_{l,N}$ by
$v_{l,N}(z_j)=p^{l|j|}(1+\cond_p(\kappa))^{N}$, and multiplication by $g_j^{-1}$
inside the shell $j$ multiplies $v_{l,N}$ by at most
$|1-\lambda|^{-1}(1+\cond_p(\kappa))^{N}$.  Since $p^{l|j|}v_{l,N}(m)=v_{2l,N}(m)$
for $m\in F^{(j)}$, the maps $H_j$ are bounded from the seminorms of weight
$2l$ to those of weight $l$ with a constant independent of $j$, and
$\widehat\bigoplus_jH_j$ is a contraction of
$\widehat\bigoplus_jC(F^{(j)},(F^{(j)})_{\sigma})$ over any set of shells.

The dilation fixes $\chi_p$ and every $1_j$, so for each of the algebras
$E_p$, $F_p^{\ge0}$, $F_p$ with its idempotents one has $M_{aa}=A_a$ with the
restricted twist.  Lemma~\ref{lem:block} for $F_p$ gives the contractibility
of $C(F_p,(F_p)_{\sigma})$, and likewise of
$C(F_p^{\ge0},(F_p^{\ge0})_{\sigma})$.  Lemma~\ref{lem:block} for $E_p$ gives
\[
C\bigl(E_p,(E_p)_{\sigma}\bigr)\ \simeq\ C\bigl(A_*,(A_*)_{\sigma}\bigr)\oplus
\widehat\bigoplus_{j<0}C\bigl(F^{(j)},(F^{(j)})_{\sigma}\bigr),
\]
and Theorem~\ref{thm:B2} gives $C(A_*,(A_*)_{\sigma})\simeq\C[0]\oplus
C(F_p^{\ge0},(F_p^{\ge0})_{\sigma})$, the summand $\C[0]$ being the class of
$\chi_p$.
\end{proof}

\begin{lemma}[The real place]\label{lem:arch}
Let $\sigma$ be the dilation $f\mapsto f(a^{-1}\,\cdot\,)$ of $E_{\R}$ by
$a\in\R^{\times}$, $a\ne1$.  Then $C(E_{\R},(E_{\R})_{\sigma})$ is
contractible.
\end{lemma}

\begin{proof}
We embed $E_{\R}$ as an ideal in a unital commutative algebra $A=E_{\R}\oplus
L$ of smooth functions on $\R^{\times}$, stable under $\sigma$, where $L$ is a
unital subalgebra containing an element $z$ with $z-\sigma(z)$ invertible in
$L$.

If $a<0$, let $z$ be the sign function and $L:=\C1\oplus\C z$.  Then $z^{2}=1$,
and $\sigma(z)=-z$ because $\sigma$ exchanges the two half-lines; so
$z-\sigma(z)=2z$, with inverse $z/2$.

If $a>0$, put $c:=\log a\ne0$, let $L$ be the Fr\'echet algebra of the
functions $x\mapsto\varphi(\log|x|)$ with $\varphi$ smooth and $2|c|$-periodic,
with the seminorms $\sup|\varphi^{(m)}|$, and let $z(x):=e^{i\pi\log|x|/c}$.
Since $\sigma$ translates $t=\log|x|$ by $c$, it preserves $L$ and
$\sigma(z)=-z$; so $z-\sigma(z)=2z$, with inverse $\bar z/2$.

In both cases $L\cap E_{\R}=0$, since a nonzero element of $L$ does not tend to
$0$ as $|t|\to\infty$, and $L\cdot E_{\R}\subseteq E_{\R}$ continuously, by the
Leibniz rule.  So $A$, with the direct sum topology, is a unital commutative
complete locally convex algebra, $E_{\R}$ is a closed ideal with closed
complement $L\ni1$, $A/E_{\R}\cong L$ as algebras, and $\sigma$ is a continuous
unital endomorphism of $A$ preserving $E_{\R}$ and $L$.  By
Lemma~\ref{lem:central}, $C(A,A_{\sigma})$ and $C(L,L_{\sigma})$ are
contractible.  The kernel $\mathcal D(A,E_{\R})$ of the projection between them
is a graded-split subcomplex, so it is contractible by
Lemma~\ref{lem:split}(c).  By Lemma~\ref{lem:hunital}(i) with $V=\{\infty\}$,
$E_{\R}$ has a right splitting which is $A$-linear on the right, and
Theorem~\ref{thm:excision-input} gives
$C(E_{\R},(E_{\R})_{\sigma})\simeq\mathcal D(A,E_{\R})$.
\end{proof}

The complexes of Lemma~\ref{lem:unit} must be distinguished from the
coinvariants of the dilation, which are not Hochschild groups of
$(F^{(j)},\sigma)$ and which we need for the comparison with the local densities
in Proposition~\ref{prop:haarrank}.

\begin{lemma}[Coinvariants of a unit]\label{lem:coinv}
Let $\sigma$ be the dilation by a unit $u\in\Z_p^{\times}$, $u\ne1$, let
$H_u:=\overline{\langle u\rangle}\subseteq\Z_p^{\times}$ be the closed subgroup
generated by $u$, and put $N(u):=[\Z_p^{\times}:H_u]\in\N\cup\{\infty\}$.  For
every shell $j$:
\begin{enumerate}
\item[(i)] the closure of $(1-\sigma)F^{(j)}$ is the closed span of the
$e_{j,\kappa}$ with $\kappa(u)\ne1$, and
$F^{(j)}/\overline{(1-\sigma)F^{(j)}}$ is isomorphic to the space of
$H_u$-invariant functions in $F^{(j)}$, the closed span of the $e_{j,\kappa}$
with $\kappa$ trivial on $H_u$;
\item[(ii)] $(1-\sigma)F^{(j)}$ is closed if and only if $u$ has finite order;
\item[(iii)] $\dim F^{(j)}/\overline{(1-\sigma)F^{(j)}}=N(u)$, and $N(u)$ is
finite if and only if $u$ has infinite order.  With $q_p=p$ for $p$ odd and
$q_2=4$,
\[
N(u)=\begin{cases}(p-1)p^{k-1}, & u\in1+q_p\Z_p,\ k:=v_p(1-u),\\
2^{k'-1}, & p=2,\ u\in3+4\Z_2,\ u\ne-1,\ k':=v_2(1+u).
\end{cases}
\]
\end{enumerate}
\end{lemma}

\begin{proof}
One has $(1-\sigma)e_{j,\kappa}=(1-\kappa(u)^{-1})e_{j,\kappa}$.

(i) If $\kappa(u)\ne1$ then $e_{j,\kappa}$ lies in the image, so the closure of
the image contains the closed span of these elements.  Conversely the
coefficient of $e_{j,\kappa}$ with $\kappa(u)=1$ is a continuous functional
vanishing on the image.  The coefficient projections split $F^{(j)}$
topologically into the two closed spans.  A character is trivial at $u$ if
and only if it is trivial on $H_u$, by continuity.

(ii) If $u$ has finite order $e$, every $\kappa(u)\ne1$ is a nontrivial $e$-th
root of unity, so $|1-\kappa(u)^{-1}|\ge|1-e^{2\pi i/e}|$ and division by
$1-\kappa(u)^{-1}$ is continuous on the closed span of (i); the image is
closed.  Let $u$ have infinite order.  Write $\Z_p^{\times}=\mu\times(1+q_p\Z_p)$
with $\mu$ the torsion subgroup, and $u=\zeta u_1$ accordingly, so that
$u_1\ne1$; let $k:=v_p(1-u_1)$.  The $p$-adic logarithm identifies $1+q_p\Z_p$
with $q_p\Z_p$ and $u_1$ with an element of valuation $k$.  Hence for every
$m>k$ there is a character $\kappa_m$ of $\Z_p^{\times}$, trivial on $\mu$ and
on $1+p^{m}\Z_p$, with $\kappa_m(u)=\kappa_m(u_1)=e^{-2\pi i/p^{m-k}}$; its
conductor exponent is at most $m$, and
$|1-\kappa_m(u)^{-1}|\le2\pi p^{-(m-k)}$.  The element
$g:=\sum_{m>k}e^{-\sqrt m}e_{j,\kappa_m}$ lies in $F^{(j)}$, because the
seminorms \eqref{eq:seminorms} weight $e_{j,\kappa_m}$ by at most
$p^{l|j|}(1+m)^{N}$, and it lies in the closure of the image by (i).  A
preimage $f$ would have coefficients
$|f_{j,\kappa_m}|=e^{-\sqrt m}|1-\kappa_m(u)^{-1}|^{-1}\ge
e^{-\sqrt m}p^{m-k}/2\pi$, which are unbounded.  So $g$ is not in the image.

(iii) The $H_u$-invariant functions in $F^{(j)}$ contain the indicator
functions of the cosets of every open subgroup containing $H_u$, and they are
spanned by the characters of $\Z_p^{\times}/H_u$; so the dimension is the
index.  If $u$ has finite order, $H_u$ is finite and the index is infinite.
If $u$ has infinite order then, in the notation of (ii), $H_u$ contains the
closure of $\langle u^{|\mu|}\rangle=\langle u_1^{|\mu|}\rangle$, which is open,
and the index is finite.  For $u\in1+q_p\Z_p$ the logarithm identifies
$\overline{\langle u\rangle}$ with $p^{k}\Z_p$, so $H_u=1+p^{k}\Z_p$ and the
index is $|\Z_p^{\times}/(1+p^{k}\Z_p)|=(p-1)p^{k-1}$.  For $p=2$ and
$u\in3+4\Z_2$, $u\ne-1$, one has $k'\ge2$ and
$v_2(1-u^{2})=v_2(1-u)+v_2(1+u)=1+k'$, so
$\overline{\langle u^{2}\rangle}=1+2^{k'+1}\Z_2$, and $H_u$ is the union of this
group and its translate by $u$, of index $2^{k'}/2=2^{k'-1}$.
\end{proof}

\begin{remark}\label{rem:closure}
On $1+q_p\Z_p\setminus\{1\}$ one has
$N(u)=(p-1)p^{k-1}=(1-p^{-1})|1-u|_p^{-1}$.  Outside this set $N(u)$ is not a
function of $|1-u|_p$.  For $p=7$ the units $u=-8$ and $u=3$ both have
$|1-u|_7=1$; but $H_{-8}=\{\pm1\}\cdot(1+7\Z_7)$ has index $3$, while $3$ is a
primitive root modulo $49$, so $H_3=\Z_7^{\times}$ has index $1$.  For $p=2$
the unit $u=3$ has $|1-u|_2^{-1}=2$ and
$H_3=(1+8\Z_2)\cup3(1+8\Z_2)$, of index $2$, whereas $(1-2^{-1})\cdot2=1$.
\end{remark}

\subsection{The shadow of a graph}

The decomposition of the Hochschild complex of a crossed product over the
group elements is classical in the algebraic setting; see Nistor
\cite{Nistor}, Getzler--Jones \cite{GetzlerJones} and Brylinski
\cite{Brylinski}.  We prove a continuous version for the rapid crossed
product by a free abelian group, by means of a Koszul resolution of the
crossed product as a bimodule.

In the next proposition $p_1,\dots,p_d$ are the primes of $S_f$,
$\Gamma=\Gamma_S\cong\Z^{d}$ with basis $p_1,\dots,p_d$, and $H$ is the height
of Section~\ref{sec:carrier}.  For a complete Fr\'echet space $V$ we write
$\ell^{1}_{H}(X;V)$, $X=\Gamma^{m}$, for the space of families
$(v_x)_{x\in X}$ with $\sum_xH(x_1)^{k}\cdots H(x_m)^{k}\nu(v_x)<\infty$ for all
$k$ and all continuous seminorms $\nu$; by Grothendieck's description of
$\ell^{1}\widehat\otimes_{\pi}V$ \cite{Treves}, one has
$\ell^{1}_{H}(X;V)=\ell^{1}_{H}(X)\widehat\otimes V$ and
$\ell^{1}_{H}(X)\widehat\otimes\ell^{1}_{H}(Y)=\ell^{1}_{H}(X\times Y)$.  For a
complex $W$ with commuting automorphisms $\alpha_{p_1},\dots,\alpha_{p_d}$ we
write $\mathcal K_{\Gamma}(\alpha;W)$ for the total complex of
$\Lambda^{\bullet}(\C^{d})\otimes W$ with the Koszul differential
\[
e_I\otimes x\longmapsto\sum_{r}(-1)^{r-1}e_{I\setminus i_r}\otimes
(\alpha_{p_{i_r}}-1)x,\qquad I=\{i_1<\dots<i_m\},
\]
plus $(-1)^{|I|}$ times the differential of $W$.

\begin{proposition}\label{prop:crossed}
Let $A$ be a unital complete Fr\'echet algebra with an action $\alpha$ of
$\Gamma$ which is tempered with respect to $H$, let
$B=A\rtimes_{\mathrm{rapid}}\Gamma$, and let $\beta$ be a continuous
automorphism of $B$ with $\beta(A)=A$ and $\beta(U_{\gamma})=U_{\gamma}$ for
every $\gamma$.  For a family of complete spaces $X_{\gamma}$, $\gamma\in\Gamma$, the
\emph{rapid direct sum} $\widehat\bigoplus_{\gamma}X_{\gamma}$ is the space of
families $(x_{\gamma})$ with $\sum_{\gamma}H(\gamma)^{k}\nu_{\gamma}(x_{\gamma})
<\infty$ for every $k$ and every compatible choice of continuous seminorms
$\nu_{\gamma}$.  Let the \emph{label} of an elementary tensor
$a_0U_{g_0}\otimes\cdots\otimes a_nU_{g_n}$ be $g_0\cdots g_n$.  Then
$C(B,{}_{1}B_{\beta})$ is the rapid direct sum of its subcomplexes of given
label, and the subcomplex of label $\gamma$ is homotopy equivalent to
$\mathcal K_{\Gamma}\bigl(\alpha;C(A,{}_{1}A_{\alpha_{\gamma}\beta})\bigr)$,
where $\alpha$ acts on all letters of $C(A,{}_{1}A_{\alpha_{\gamma}\beta})$.
\end{proposition}

The faces of $C(B,{}_{1}B_{\beta})$ preserve labels because $\Gamma$ is abelian
and $\beta$ fixes every $U_{\gamma}$, and the label projections are continuous
by the definition of the rapid crossed product.  The action $\alpha$ commutes
with the boundary of $C(A,{}_{1}A_{\alpha_{\gamma}\beta})$ because $\beta$
commutes with every $\alpha_g$:
$\beta(\alpha_g(a))=\beta(U_gaU_g^{-1})=U_g\beta(a)U_g^{-1}=\alpha_g(\beta(a))$.

\begin{proof}
\emph{Step 0: freeness of $B$ over $A$.}  The maps
$\ell^{1}_{H}(\Gamma)\widehat\otimes A\to B$, $\delta_g\otimes a\mapsto
aU_g$ and $\delta_g\otimes a\mapsto U_ga$, are topological isomorphisms of
left, respectively right, $A$-modules: the first is the definition of $B$,
and the second differs from it by $a\mapsto\alpha_g(a)$ on the $g$-th
component, since $U_ga=\alpha_g(a)U_g$; this map and its inverse cost at most
$CH(g)^{M}$ by temperedness and are absorbed by the rapid weights.
Consequently, for a complete $A$-bimodule $X$, the map $U_g\otimes x\otimes
U_h\mapsto\delta_g\otimes x\otimes\delta_h$ identifies
$B\widehat\otimes_AX\widehat\otimes_AB$ with $\ell^{1}_{H}(\Gamma)\widehat
\otimes X\widehat\otimes\ell^{1}_{H}(\Gamma)$ as a topological vector space.
Under this identification the map induced by an $A$-bimodule map $\phi$ is
$1\otimes\phi\otimes1$.  In particular, if a complex of $A$-bimodules has a
continuous linear contraction $s$, not necessarily $A$-linear, then
$1\otimes s\otimes1$ is a contraction of the induced complex, and it preserves
the labels $g,h$.

\emph{Step 1: a Koszul resolution of $B$.}  Put
$K_m:=B\widehat\otimes_AB\otimes\Lambda^{m}(\C^{d})$, a $B$-bimodule through
the outer factors, and for $i=1,\dots,d$ let
\[
D_i(x\otimes y):=xU_{p_i}\otimes U_{p_i}^{-1}y-x\otimes y .
\]
This is well defined on the balanced tensor product, because
$xaU_{p_i}\otimes U_{p_i}^{-1}y=xU_{p_i}\alpha_{p_i}^{-1}(a)\otimes
U_{p_i}^{-1}y=xU_{p_i}\otimes U_{p_i}^{-1}ay$; it is a bimodule map, the
$D_i$ commute, and multiplication $\mu:K_0\to B$ satisfies $\mu D_i=0$.  We
give $K_{\bullet}$ the differential $\partial(\xi\otimes e_I):=
\sum_{r}(-1)^{r-1}D_{i_r}\xi\otimes e_{I\setminus i_r}$,
$I=\{i_1<\dots<i_m\}$; it squares to zero because the $D_i$ commute, and
$K_{\bullet}$ is augmented to $B$ by $\mu$.  We claim that the augmented
complex is contractible.  By Step 0, $B\widehat\otimes_AB
\cong\ell^{1}_{H}(\Gamma\times\Gamma)\widehat\otimes A$ through
$U_g\otimes U_ha\leftrightarrow\delta_{(g,h)}\otimes a$, and in these
coordinates $D_i$ replaces $\delta_{(g,h)}$ by
$\delta_{(gp_i,hp_i^{-1})}-\delta_{(g,h)}$ while $\mu$ sends
$\delta_{(g,h)}\otimes a$ to $U_{gh}a$.  Change variables to $(g,s)$ with
$s=gh$; the weight families $H(g)^kH(h)^k$ and $H(g)^kH(s)^k$ are cofinal in
each other, since $H(s)\le H(g)H(h)$ and $H(h)\le H(g)H(s)$.  So the augmented
complex is the tensor product of $\ell^{1}_{H}(\Gamma)\widehat\otimes A$, in
the variable $s$ and with zero differential, and the augmented Koszul complex
\begin{equation}\label{eq:koszulell}
0\to\ell^{1}_{H}(\Gamma)\otimes\Lambda^{d}\to\cdots\to
\ell^{1}_{H}(\Gamma)\otimes\Lambda^{0}\xrightarrow{\ \Sigma\ }\C\to0
\end{equation}
of the translations $\tau_i:\delta_g\mapsto\delta_{gp_i}$, augmented by
$\Sigma(\phi)=\sum_g\phi(g)$.  It suffices to give \eqref{eq:koszulell} a
continuous contraction.  Writing $\gamma=a/b$ in lowest terms, the height
satisfies $H(\gamma)\le\prod_ip_i^{|v_{p_i}(\gamma)|}=ab\le H(\gamma)^{2}$, so
$\ell^{1}_{H}(\Gamma)=\widehat\bigotimes_i\ell^{1}_{H_i}(\Z)$ with
$H_i(n)=p_i^{|n|}$, and \eqref{eq:koszulell} is the tensor product of the $d$
augmented complexes $0\to\ell^{1}_{H_i}(\Z)\xrightarrow{\tau-1}
\ell^{1}_{H_i}(\Z)\xrightarrow{\Sigma}\C\to0$.  Each of these is contracted
by $s_{-1}(1)=\delta_0$ and
\[
s_0(\delta_n)=\sum_{0\le j<n}\delta_j\quad(n>0),\qquad
s_0(\delta_0)=0,\qquad
s_0(\delta_n)=-\sum_{n\le j<0}\delta_j\quad(n<0);
\]
indeed $(\tau-1)s_0(\delta_n)=\delta_n-\delta_0=\delta_n-s_{-1}\Sigma(\delta_n)$
and $s_0((\tau-1)\delta_n)=s_0(\delta_{n+1})-s_0(\delta_n)=\delta_n$ for
every $n$.  Since $H_i(j)\le H_i(n)$ for $j$ between $0$ and $n$,
$\lVert s_0\delta_n\rVert_k\le|n|p_i^{k|n|}\le p_i^{(k+1)|n|}$, so $s_0$ is
continuous from the weight $k+1$ to the weight $k$.  If $X$ and $X'$ are
complexes with augmentations $\varepsilon,\varepsilon'$ to $\C$, sections
$\iota,\iota'$ of them, and maps $s,s'$ with
$\partial s+s\partial=1-\iota\varepsilon$ and
$\partial s'+s'\partial=1-\iota'\varepsilon'$, where $\iota\varepsilon$ is
extended by zero to positive degrees, then
$s\otimes1+\iota\varepsilon\otimes s'$ is a homotopy from the identity of
$X\otimes X'$ to $\iota\varepsilon\otimes\iota'\varepsilon'$.  Applying this
$d-1$ times gives a continuous contraction of \eqref{eq:koszulell}, which
proves the claim.  The contraction changes $g$ at fixed $s=gh$.

\emph{Step 2: a free resolution.}  Let $\mathrm{Bar}(A)\to A$ be the bar
resolution, $\mathrm{Bar}_n(A)=A\widehat\otimes A^{\widehat\otimes
n}\widehat\otimes A$, with the continuous contraction
$\xi\mapsto1\otimes\xi$.  By Step 0 the complexes
$P_{m,\bullet}:=B\widehat\otimes_A\mathrm{Bar}_{\bullet}(A)\widehat\otimes_A
B\otimes\Lambda^{m}$, augmented to $K_m$, are contractible, and
$P_{m,n}\cong B\widehat\otimes A^{\widehat\otimes n}\widehat\otimes
B\otimes\Lambda^{m}$ is a free $B$-bimodule.  The maps
\[
\widetilde D_i(x\otimes a_1\otimes\cdots\otimes a_n\otimes y):=
xU_{p_i}\otimes\alpha_{p_i}^{-1}(a_1)\otimes\cdots\otimes
\alpha_{p_i}^{-1}(a_n)\otimes U_{p_i}^{-1}y-x\otimes a_1\otimes\cdots\otimes
a_n\otimes y
\]
are continuous bimodule maps lifting $D_i$; they commute with the bar
differential, because $xU_{p_i}\alpha_{p_i}^{-1}(a_1)=xa_1U_{p_i}$ and
$\alpha_{p_i}^{-1}(a_n)U_{p_i}^{-1}y=U_{p_i}^{-1}a_ny$, and they commute with
one another.  Let $P$ be the total complex of $P_{\bullet,\bullet}$, with
differential $\sum_r(-1)^{r-1}\widetilde D_{i_r}$ on the factor $e_I$ plus
$(-1)^{|I|}$ times the bar differential.  The augmentations of the columns
give a map of complexes of bimodules $P\to K_{\bullet}$.  Its cone is filtered
by the $d+1$ columns, the filtration is graded-split, and the graded pieces
are the augmented columns, which are contractible; so the cone is contractible
by Lemma~\ref{lem:split}(d).  Together with Step 1, the complex $P$ augmented
to $B$ is contractible.  Thus $P$ is a resolution of $B$ by free
$B$-bimodules with a continuous linear contraction, and all the maps
involved preserve the total label, in which $U_g$ has label $g$ and the
elements of $A$ have label $1$.

\emph{Step 3: comparison with the bar resolution.}  The bar resolution
$\mathrm{Bar}(B)\to B$ is also a resolution by free bimodules with a
continuous linear contraction $\xi\mapsto1\otimes\xi$, which preserves labels.
The comparison theorem of relative homological algebra
\cite[Chapter~IX]{MacLane} applies: a continuous bimodule map out of a free
bimodule $B\widehat\otimes V\widehat\otimes B$ is determined by its restriction
to $1\otimes V\otimes1$, which may be any continuous linear map, and one
defines chain maps $\mathrm{Bar}(B)\rightleftarrows P$ over $\id_B$, and
homotopies between their composites and the identities, degree by degree on
the generating spaces by composing the previously constructed map, the
differential and the contraction of the target.  The resulting maps are
continuous bimodule maps, and they preserve labels because the contractions
do.

\emph{Step 4: the Hochschild complexes.}  Let $M={}_{1}B_{\beta}$.  Then
$C(B,M)=M\widehat\otimes_{B^{e}}\mathrm{Bar}(B)$, and for a complete space
$V$ one has $M\widehat\otimes_{B^{e}}(B\widehat\otimes V\widehat\otimes B)\cong
M\widehat\otimes V$, by $m\otimes(x\otimes v\otimes y)\mapsto y\cdot m\cdot
x\otimes v$.  Under this identification $M\widehat\otimes_{B^{e}}P_{m,\bullet}$
is $C(A,M)\otimes\Lambda^{m}$, the Hochschild complex of $A$ with coefficients
in $M$ regarded as an $A$-bimodule, and $\widetilde D_i$ becomes
\[
m\otimes a_1\otimes\cdots\otimes a_n\mapsto
U_{p_i}^{-1}\cdot m\cdot U_{p_i}\otimes\alpha_{p_i}^{-1}(a_1)\otimes\cdots
\otimes\alpha_{p_i}^{-1}(a_n)-m\otimes a_1\otimes\cdots\otimes a_n .
\]
Applying $M\widehat\otimes_{B^{e}}(-)$ to Step 3 gives a label-preserving
homotopy equivalence between $C(B,M)$ and the total complex of
$C(A,M)\otimes\Lambda^{\bullet}$ with the Koszul differential built from these
operators.

\emph{Step 5: the label-$\gamma$ piece.}  As an $A$-bimodule,
$M=\widehat\bigoplus_{\gamma}AU_{\gamma}$, and $a\cdot(cU_{\gamma})\cdot a'=
acU_{\gamma}\beta(a')=ac\,\alpha_{\gamma}(\beta(a'))U_{\gamma}$; so
$AU_{\gamma}\cong{}_{1}A_{\alpha_{\gamma}\beta}$ by $cU_{\gamma}\mapsto c$.  The
operator $m\mapsto U_{p_i}^{-1}\cdot m\cdot U_{p_i}$ sends $cU_{\gamma}$ to
$U_{p_i}^{-1}cU_{\gamma}U_{p_i}=\alpha_{p_i}^{-1}(c)U_{\gamma}$, since
$\beta(U_{p_i})=U_{p_i}$ and $\Gamma$ is abelian.  So the label-$\gamma$ part
of the complex of Step 4 is the total complex of $\Lambda^{\bullet}\otimes
C(A,{}_{1}A_{\alpha_{\gamma}\beta})$ with the Koszul differential of the
operators $\alpha_{p_i}^{-1}-1$ acting on all letters.  Finally
$\alpha_{p_i}^{-1}-1=-\alpha_{p_i}^{-1}(\alpha_{p_i}-1)$, and the
$\alpha_{p_i}$ are invertible and commute with one another and with the
Hochschild boundary, so $e_I\otimes x\mapsto e_I\otimes
\prod_{i\in I}(-\alpha_{p_i})x$ is an isomorphism from
$\mathcal K_{\Gamma}(\alpha;C(A,{}_{1}A_{\alpha_{\gamma}\beta}))$ onto this
complex.
\end{proof}

In the next theorem $w\in C_S$ is identified with its lift \eqref{eq:section}
in $A_S^{\times}$, and for $a\in A_S^{\times}$ we write $(E_P^{u})_{a}$ for the
bimodule twisted by the dilation $\beta_a$ of $E_P^{u}:=E_P\oplus\C1$, and
likewise $(E_v)_{x}$ for the local algebra $E_v$ twisted by the dilation by
$x\in\Q_v^{\times}$.

\begin{theorem}\label{thm:shadow}
Let $w\in C_S$, and let $d=|S_f|$.
\begin{enumerate}
\item[(a)] \emph{(label splitting)}  There is a homotopy equivalence
\begin{equation}\label{eq:labelsplit}
\Sh(E_w)\ \simeq\ \widehat\bigoplus_{\gamma\in\Gamma_S}
\mathcal K_{\Gamma}\Bigl(\alpha;\,C\bigl(E_P^{u},(E_P^{u})_{\gamma w}\bigr)\Bigr),
\end{equation}
which carries the part of $\Sh(E_w)$ of lattice label $\gamma$ to the
$\gamma$-th summand; here $\widehat\bigoplus_{\gamma}$ is the $H$-rapid sum.
\item[(b)] \emph{(K\"unneth)}  For every $a\in A_S^{\times}$,
\[
C\bigl(E_P^{u},(E_P^{u})_{a}\bigr)=C^{\mathrm{sc}}\oplus C^{+},\qquad
C^{\mathrm{sc}}\simeq\C[0],\qquad
C^{+}\simeq\mathrm{Tot}\,\widehat\bigotimes_{v\in S}C\bigl(E_v,(E_v)_{a_v}\bigr),
\]
the decomposition being stable under $\Gamma_S$, which acts trivially on
$C^{\mathrm{sc}}$.
\item[(c)] \emph{(local factors)}  For $p\in S_f$ and $x\in\Q_p^{\times}$,
$x\ne1$, the complex $C(E_p,(E_p)_{x})$ is homotopy equivalent to $\C[0]$, the
class being that of $\chi_p$.  For $a\in\R^{\times}$, $a\ne1$, the complex
$C(E_{\R},(E_{\R})_{a})$ is contractible.  For the trivial twist,
$\HH_0(E_p,E_p)=E_p$.
\item[(d)] \emph{(label pieces)}  If $\gamma\in\Gamma_S$ and $\gamma
w_{\infty}\ne1$, the part of $\Sh(E_w)$ of label $\gamma$ is homotopy
equivalent to $\Lambda^{\bullet}(\C^{d})$ with zero differential,
$\Lambda^{m}$ in degree $m$.
\end{enumerate}
\end{theorem}

\begin{proof}
(a)  By \eqref{eq:shadow} and \eqref{eq:graphbimod},
$\Sh(E_w)=C(U,{}_{1}U_{\beta_w})$.  Proposition~\ref{prop:crossed} applies with
$A=E_P^{u}$, $\Gamma=\Gamma_S$ and $\beta=\beta_w$: $E_P^{u}$ is a unital
complete Fr\'echet algebra with a tempered $\Gamma_S$-action by
Section~\ref{sec:carrier}, and $\beta_w$ is continuous, preserves $E_P^{u}$ and
fixes every $U_{\gamma}$.  The equivalences for the individual labels are the
restrictions of the single label-preserving homotopy equivalence of Step 4 of
that proof, and the isomorphism of Step 5 is defined on all labels at once; so
they assemble into a homotopy equivalence onto the rapid sum.  The twist of the
$\gamma$-th summand is $\alpha_{\gamma}\beta_w=\beta_{\gamma w}$ on the right.
Since $E_P^{u}$ is commutative, the reversal \eqref{eq:reversal} identifies
its complex with $C(E_P^{u},(E_P^{u})_{\gamma w})$, compatibly with $\alpha$.

(b) is Theorem~\ref{thm:B2} for $R=E_P^{u}$, $I=E_P$, with the right splitting
of Lemma~\ref{lem:hunital}(i), followed by Theorem~\ref{thm:kunneth}.  The
action of $\Gamma_S$ fixes $1$ and preserves $E_P$.

(c)  The first statement is Lemma~\ref{lem:nonunit} if $v_p(x)\ne0$ and
Lemma~\ref{lem:unit} if $x$ is a unit.  The second is Lemma~\ref{lem:arch}.
The third holds because $E_p$ is commutative, so that $b=0$ on $C_1(E_p,E_p)$.

(d)  Let $a=\gamma w$.  By (c) the real factor of $C^{+}$ in (b) is
contractible, by a contraction $h$; then $h\otimes\id$ contracts the tensor
product, so $C^{+}$ is contractible.  The complex
$\mathcal K_{\Gamma}(\alpha;C^{+})$ is filtered by the exterior degree, the
graded pieces are the complexes $\Lambda^{m}\otimes C^{+}$, and it is
contractible by Lemma~\ref{lem:split}(d).  On $C^{\mathrm{sc}}$ the operators
$\alpha_{p_i}-1$ vanish, so $\mathcal K_{\Gamma}(\alpha;C^{\mathrm{sc}})=
\Lambda^{\bullet}\otimes C^{\mathrm{sc}}$ with the differential of
$C^{\mathrm{sc}}$ alone, which is homotopy equivalent to
$\Lambda^{\bullet}\otimes\C[0]$.  The claim follows from (a).
\end{proof}

\begin{remark}\label{rem:notorsion}
At a finite place the fixed point $x_p=0$ of a dilation $x\ne1$ contributes
the complex $\C[0]$, spanned by the class of $\chi_p$, and the homotopy
equivalences of Lemmas~\ref{lem:nonunit} and \ref{lem:unit} do not involve $x$
beyond the condition $x\ne1$.  No quantity depending on $x$, such as
$|1-x|_p$, can be recovered from the twisted Hochschild complex.  The
dependence on $x$ enters the theory through the densities of
Section~\ref{sec:trace}.
\end{remark}

\begin{corollary}\label{cor:blind}
Let $w\in C_S$, $\gamma\in\Gamma_S$ and $a=\gamma w$.
\begin{enumerate}
\item[(i)] At every finite place $p$ with $a_p\ne1$ the local factor
$C(E_p,(E_p)_{a_p})$ of the label-$\gamma$ part of $\Sh(E_w)$ is homotopy
equivalent to $\C[0]$.
\item[(ii)] If $a_{\infty}\ne1$, the label-$\gamma$ part of $\Sh(E_w)$ is
homotopy equivalent to $\Lambda^{\bullet}(\C^{|S_f|})$ with zero differential.
Its homology has dimension $\binom{|S_f|}{m}$ in degree $m$ and Euler
characteristic $0$.
\item[(iii)] If $w_{\infty}$ is not a positive $S$-unit, then (ii) holds for
every label $\gamma$.  In particular, for any two such classes $w,w'$ and any
$\gamma$, the label-$\gamma$ parts of $\Sh(E_w)$ and $\Sh(E_{w'})$ are homotopy
equivalent.
\end{enumerate}
\end{corollary}

\begin{proof}
(i) and (ii) are Theorem~\ref{thm:shadow}(c),(d); the Euler characteristic is
$\sum_m(-1)^{m}\binom{|S_f|}{m}=0$ since $S_f\ne\emptyset$.  For (iii),
$\gamma w_{\infty}=1$ would give $w_{\infty}=\gamma^{-1}\in\Gamma_S$.
\end{proof}

Thus the ranks and Euler characteristics of the label parts of the shadow are
the same for all generators $E_w$ with $w_{\infty}\notin\Gamma_S$, and by
Remark~\ref{rem:notorsion} the local complexes do not depend on the dilation.
The trace of the next section depends on the class through the densities
$\tau_v$; see Theorem~\ref{thm:separate}.  The relation between the coinvariant count
$N(u)$ of Lemma~\ref{lem:coinv} and the density $\tau_p$ is recorded in
Proposition~\ref{prop:haarrank}.
\section{The trace}\label{sec:trace}

\subsection{The local densities}

\begin{definition}\label{def:tau}
For a place $v$ of $\Q$ and $x\in\Q_v^{\times}$ with $x\ne1$ put
\begin{equation}\label{eq:tau}
\tau_v(x):=\frac{|x|_v^{1/2}}{|1-x|_v}.
\end{equation}
\end{definition}

\begin{lemma}\label{lem:tau}
Let $p$ be finite and $x=p^{k}u$ with $u\in\Z_p^{\times}$.
\begin{enumerate}
\item[(i)] If $k\ne0$ then $\tau_p(x)=p^{-|k|/2}$.
\item[(ii)] If $k=0$ and $x\ne1$ then $\tau_p(x)=|1-x|_p^{-1}=p^{\,v_p(1-x)}$,
a finite number.
\end{enumerate}
Moreover, for an integer $N\ge0$ the truncated density
$\tau_p^{(N)}(x):=|x|_p^{1/2}\min\bigl(p^{N},|1-x|_p^{-1}\bigr)$ equals
$\tau_p(x)$ for every $N\ge v_p(1-x)$; in case (i) this holds for every
$N\ge0$.
\end{lemma}

\begin{proof}
(i) $|x|_p^{1/2}=p^{-k/2}$.  If $k>0$ then $|x|_p<1$, so $|1-x|_p=1$ and
$\tau_p=p^{-k/2}$.  If $k<0$ then $|x|_p>1$, so $|1-x|_p=|x|_p=p^{-k}$ and
$\tau_p=p^{-k/2}p^{k}=p^{k/2}=p^{-|k|/2}$.  In both cases $|1-x|_p\ge1$, so
$\min(p^{N},|1-x|_p^{-1})=|1-x|_p^{-1}$ for every $N\ge0$.

(ii) For a unit $x\ne1$ the integer $j:=v_p(1-x)\ge0$ is finite and
$|1-x|_p^{-1}=p^{j}$; hence $\min(p^{N},p^{j})=p^{j}$ for every $N\ge j$.
\end{proof}

\begin{remark}\label{rem:unitdiv}
The function $u\mapsto|1-u|_p^{-1}=p^{v_p(1-u)}$ is finite at every unit
$u\ne1$ but is not integrable on $\Z_p^{\times}$: for the probability Haar
measure the set $\{v_p(1-u)=n\}$ has measure $p^{-n}$ for $n\ge1$, and
$\int_{\Z_p^{\times}}\min(p^{N},|1-u|_p^{-1})\,du=N+1$.  The regularization
below therefore subtracts the value at the identity.
\end{remark}

\begin{lemma}\label{lem:inv}
$\tau_v(x^{-1})=\tau_v(x)$ for every place $v$ and every $x\ne1$.
\end{lemma}

\begin{proof}
$|x^{-1}|_v^{1/2}=|x|_v^{-1/2}$ and
$|1-x^{-1}|_v=|x|_v^{-1}|x-1|_v$; dividing, the factors $|x|_v^{-1/2}$ and
$|x|_v^{-1}$ combine to $|x|_v^{1/2}$.
\end{proof}

At the archimedean place, in the multiplicative coordinate $x=\pm e^{t}$,
\begin{equation}\label{eq:tauinf}
\tau_{\infty}(+e^{t})=\frac{1}{2|\sinh(t/2)|},\qquad
\tau_{\infty}(-e^{t})=\frac{1}{2\cosh(t/2)} .
\end{equation}

\subsection{The trace}

For $v\in S$ let $j_v:\Q_v^{\times}\to C_S$ be the map sending $x$ to the class
of the idele having $x$ at $v$ and $1$ at every other place of $S$.  By
Lemma~\ref{thm:split}, the canonical representative of that class in
$C_S\cong\R^{\times}\times\prod_{q\in S_f}\Z_q^{\times}$ is obtained by
dividing by $\varphi$ of the idele, and therefore
\begin{equation}\label{eq:slices}
\begin{aligned}
j_p(p^{k}w)&=\bigl(p^{-k}\text{ at }\infty,\ w\text{ at }p,\
p^{-k}\text{ at every other place of }S\bigr),\qquad p\in S_f,\\
j_{\infty}(x)&=\bigl(x\text{ at }\infty,\ 1\text{ at every finite place of }S
\bigr).
\end{aligned}
\end{equation}

\begin{definition}\label{def:trace}
Fix a real constant $c$.  The \emph{trace} is the linear functional on $\DA$
defined by its local decomposition
\begin{equation}\label{eq:tracedef}
T(F):=\sum_{v\in S}W_v(F)+c\,F(1),\qquad F\in\DA ,
\end{equation}
where, at a finite place $p$, with $a:=F\circ j_p$ and with $\dstar x$ the Haar
measure of $\Q_p^{\times}$ giving every shell $p^{k}\Z_p^{\times}$ the volume
$\log p$ (this is $\log p$ times the measure used on the factor
$\Z_p^{\times}$ of $C_S$),
\begin{equation}\label{eq:Wp}
W_p(F)=\int_{v_p(x)\neq0}a(x)\,\tau_p(x)\dstar x
\ +\ (\log p)\int_{\Z_p^{\times}}\frac{a(u)-a(1)}{|1-u|_p}\,du,
\end{equation}
the second integral against the probability Haar measure of
$\Z_p^{\times}$; and at the archimedean place, with the two profiles
\[
b_{+}(t):=F\bigl(j_{\infty}(e^{t})\bigr),\qquad
b_{-}(t):=F\bigl(j_{\infty}(-e^{t})\bigr),\qquad t\in\R,
\]
of $F$ on the two real half-lines (so that $b_{+}(0)=F(1)$),
\begin{equation}\label{eq:Winf}
\begin{aligned}
W_{\infty}(F)={}&(\gamma+\log4\pi)\,b_{+}(0)
+\int_0^{\infty}\Bigl[K_{s}(t)\bigl(b_{+}(t)+b_{+}(-t)\bigr)
-2K(t)\,b_{+}(0)\,e^{-t/2}\Bigr]dt\\
&+\int_0^{\infty}K_{c}(t)\bigl(b_{-}(t)+b_{-}(-t)\bigr)dt,
\end{aligned}
\end{equation}
with the kernels
\[
K_{s}(t)=\frac{1}{4\sinh(t/2)},\qquad K_{c}(t)=\frac{1}{4\cosh(t/2)},\qquad
K(t)=K_{s}(t)+K_{c}(t)=\frac{e^{-t/2}}{1-e^{-2t}} .
\]
\end{definition}

\begin{convention}[The archimedean term for all test functions]\label{conv:Winf}
Formula \eqref{eq:Winf} is the definition of $W_{\infty}$ on all of $\DA$,
with no parity assumption: the positive half-line enters only through $b_{+}$
and the kernel $K_{s}$, the negative half-line only through $b_{-}$ and the
kernel $K_{c}$, and the regularization (the constant $\gamma+\log4\pi$ and the
counterterm $2K(t)b_{+}(0)e^{-t/2}$) involves only the value $b_{+}(0)=F(1)$ at
the identity, which is the only point where $\tau_{\infty}$ is singular.  Both
integrals in \eqref{eq:Winf} converge absolutely: near $t=0$ one has
$2K_{s}(t)=t^{-1}+O(t)$ and $2K(t)e^{-t/2}=t^{-1}+O(1)$, so the first bracket is
$b_{+}(t)+b_{+}(-t)-2b_{+}(0)$ times $O(t^{-1})$ plus a bounded function times
$b_{+}(0)$, hence bounded; at infinity $K_{s},K_{c}=O(e^{-t/2})$ and
$K(t)e^{-t/2}=O(e^{-t})$, while $b_{\pm}$ have compact support.  Writing
$F=F^{\mathrm{ev}}+F^{\mathrm{odd}}$ for the decomposition under translation by
$j_{\infty}(-1)$, the even part has $b_{+}=b_{-}=:b$ and \eqref{eq:Winf}
becomes
\begin{equation}\label{eq:Winfeven}
W_{\infty}(F^{\mathrm{ev}})=(\gamma+\log4\pi)\,b(0)
+\int_0^{\infty}K(t)\bigl[b(t)+b(-t)-2b(0)e^{-t/2}\bigr]dt,
\end{equation}
while for the odd part $b_{-}=-b_{+}$ and the $K_{c}$-integral enters with the
opposite sign.  If $F$ vanishes on a neighbourhood of $1$, then $b_{+}(0)=0$,
the counterterm drops, and since $\tau_{\infty}(\pm e^{t})$ equals
$2K_{s}(|t|)$, respectively $2K_{c}(|t|)$, by \eqref{eq:tauinf}, while
$\dstar x=dt/2$ on each half-line by \eqref{eq:archmeasure},
\begin{equation}\label{eq:Winfint}
\begin{aligned}
W_{\infty}(F)&=\int_{\R^{\times}}F\bigl(j_{\infty}(x)\bigr)\,\tau_{\infty}(x)\,
\frac{dx}{2|x|}\\
&=\int_0^{\infty}K_{s}(t)\bigl(b_{+}(t)+b_{+}(-t)\bigr)dt
+\int_0^{\infty}K_{c}(t)\bigl(b_{-}(t)+b_{-}(-t)\bigr)dt .
\end{aligned}
\end{equation}
Thus $W_{\infty}$ is the regularized integral of $F$ against $\tau_{\infty}$
over $\R^{\times}$, with the measure \eqref{eq:archmeasure}, regularized at the
identity only; this is the meaning of $\int'_{\R^{\times}}$ wherever it
appears below.
\end{convention}

At a finite place the two terms of \eqref{eq:Wp} are likewise the integral of
$a\,\tau_p$ over the shells $v_p(x)\ne0$ and the integral of
$(a-a(1))\,\tau_p$ over $\Z_p^{\times}$, both for the measure $\dstar x$, since
$\tau_p(u)=|1-u|_p^{-1}$ on the units.  The integrand of the second vanishes
near $u=1$ because $a$ is locally constant.

$T$ is linear, and it is continuous on each Fr\'echet step of $\DA$, that is,
on the functions supported in a fixed compact set $K$ and invariant under
$\Theta_N=\prod_{q\in S_f}(1+q^{N}\Z_q)$, $N\ge1$.  Indeed, the shell integrals are
bounded by $\sup|F|\,(\log p)\sum_{k\ne0}p^{-|k|/2}$; the unit integral in
\eqref{eq:Wp} is bounded by $2p^{N}\sup|F|$, because $a(u)=a(1)$ on
$1+p^{N}\Z_p$ and $|1-u|_p^{-1}<p^{N}$ elsewhere; and at the archimedean place
the bracket in \eqref{eq:Winf} equals
$K_{s}(t)\bigl(b_{+}(t)+b_{+}(-t)-2b_{+}(0)\bigr)
+b_{+}(0)\bigl(2K_{s}(t)-2K(t)e^{-t/2}\bigr)$, in which
$|b_{+}(t)+b_{+}(-t)-2b_{+}(0)|\le\min\bigl(t^{2}\sup|b_{+}''|,4\sup|b_{+}|\bigr)$
by Taylor's formula and $2K_{s}(t)-2K(t)e^{-t/2}$ is integrable on
$(0,\infty)$, so that
$|W_{\infty}(F)|\le C\bigl(\sup|F|+\sup|\partial_t^{2}F|\bigr)$ with an
absolute constant $C$, where $\partial_t=x_{\infty}\partial_{x_{\infty}}$.

\begin{theorem}[Absolute convergence off the identity class]\label{thm:absconv}
For $F\in\DA^{\mathrm{off}}$,
\begin{equation}\label{eq:absconv}
T(F)=\sum_{v\in S}\int_{\Q_v^{\times}}F(j_v(x))\,\tau_v(x)\dstar x ,
\end{equation}
a finite sum of absolutely convergent integrals; in particular $T(F)$ does not
depend on $c$.
\end{theorem}

\begin{proof}
For $F\in\DA^{\mathrm{off}}$ one has $F(1)=0$, so the term $cF(1)$ drops
whatever $c$ is; and $a=F\circ j_v$ vanishes on a neighbourhood of $1$ in
$\Q_v^{\times}$, because $j_v$ is continuous with $j_v(1)=1$.  Hence the
counterterms $(\log p)a(1)$, $(\gamma+\log4\pi)b_{+}(0)$ and $2K(t)b_{+}(0)e^{-t/2}$ all
vanish.  What remains at a finite place is
$\int_{v_p\ne0}a\tau_p\dstar x+(\log p)\int_{\Z_p^{\times}}a(u)|1-u|_p^{-1}du$.
The first is a finite sum over the finitely many shells met by the compact
support of $a$, each term bounded by $\sup|a|\cdot p^{-|k|/2}\log p$; the
second is the integral of a compactly supported locally constant function
against a function bounded on $\supp a$, by Lemma~\ref{lem:tau}(ii), since
$|1-u|_p^{-1}\le p^{j_0}$ with $j_0$ determined by the distance of $\supp a$
from $1$.  At the archimedean place \eqref{eq:Winfint} applies: $K_{s}$ is continuous
on the compact support of $b_{+}$, which avoids $t=0$, and $K_{c}$ is
continuous everywhere.
\end{proof}

\begin{theorem}[Uniqueness off the identity class]\label{thm:unique}
Let $D$ be a continuous linear functional on $\DA$ vanishing on
$\DA^{\mathrm{off}}$.  Then there are an integer $n\ge0$ and
$c_0,\dots,c_n\in\C$ with
\begin{equation}\label{eq:pointsupport}
D(F)=\sum_{k=0}^{n}c_k\,(\partial_t^{k}F)(1),\qquad F\in\DA,
\end{equation}
where $\partial_t=x_{\infty}\partial_{x_{\infty}}$ is the derivative in the
archimedean coordinate.  In particular, if $T'$ is a continuous linear
functional on $\DA$ agreeing with \eqref{eq:absconv} on $\DA^{\mathrm{off}}$,
then $T'-T$ has the form \eqref{eq:pointsupport}; and $T'-T$ involves no
derivative if and only if $T'=T+c'F(1)$ for some $c'\in\C$.
\end{theorem}

\begin{proof}
Let $N\ge1$, $\Theta_N=\prod_{q\in S_f}(1+q^{N}\Z_q)$, and let $F\in\DA$ be
$\Theta_N$-invariant.  Then $F=\sum_{g}F\cdot\mathbf 1_{\R^{\times}\times g\Theta_N}$, a
finite sum over the cosets $g\Theta_N$ of $\Theta_N$ in $\prod_q\Z_q^{\times}$.  For
$g\notin \Theta_N$ the summand vanishes on the open neighbourhood
$\R^{\times}\times \Theta_N$ of $1$, so $D$ annihilates it.  Hence
$D(F)=D(f\otimes\mathbf 1_{\Theta_N})$ with $f(x):=F(x,1)$.  The map
$f\mapsto f\otimes\mathbf 1_{\Theta_N}$ is continuous from $C_c^{\infty}(\R^{\times})$
to $\DA$, so $D_N(f):=D(f\otimes\mathbf 1_{\Theta_N})$ is a distribution on
$\R^{\times}$.  It vanishes on every $f$ vanishing near $1$, hence is supported
at $\{1\}$, and a distribution supported at a point is a finite linear
combination of derivatives of the Dirac measure at that point
\cite[\S2.3]{Hormander}; thus
$D_N(f)=\sum_{k\le n_N}c^{(N)}_k(\partial_t^{k}f)(1)$.  Finally
$f\otimes(\mathbf 1_{\Theta_N}-\mathbf 1_{\Theta_{N+1}})$ vanishes on
$\R^{\times}\times \Theta_{N+1}$ and so lies in $\DA^{\mathrm{off}}$; therefore
$D_N=D_{N+1}$, and the coefficients do not depend on $N$.  This proves
\eqref{eq:pointsupport}.  The second assertion follows from
Theorem~\ref{thm:absconv}, by which $T$ is given by \eqref{eq:absconv} on
$\DA^{\mathrm{off}}$.
\end{proof}

\begin{remark}\label{rem:diagonal}
The unit for convolution is the Dirac measure $\delta_1$ at the identity
class, which is not an element of $\DA$; under $u\mapsto E_u$ it corresponds to
the diagonal $\Delta=E_1$.  The functional $T$ is not defined on it, the
densities $\tau_v$ being singular at $1$.
\end{remark}

\begin{convention}\label{conv:c}
From here on $c$ denotes a fixed real constant, the same for every set $S$.
Every statement below holds for every $c$.  For $c=0$ the function
$\Th$ of Theorem~\ref{thm:symbol} is the symbol of $T$, by
\eqref{eq:plancherel}.
\end{convention}

\subsection{Separation}

The following computation involves no regularization.  It exhibits two
sections with the same pushforward to the modulus group, hence the same values
of the two functionals $\lambda_{\mp}$ of \eqref{eq:lambdadef}, on which $T$
differs.

\begin{theorem}\label{thm:separate}
Fix $p\in S_f$ and units $u,u'\in1+p\Z_p$, $u,u'\ne1$, with
$j:=v_p(1-u)\ne j':=v_p(1-u')$, and fix $m>\max(j,j')$.  Let
$\theta:=\min_{q\in S_f}q^{1/2}>1$ and let $h\in C_c^{\infty}(\R^{\times})$
satisfy $h\ge0$, $\supp h\subset(\theta^{-1},\theta)$ and $h(1)>0$.  In the
coordinates of Lemma~\ref{thm:split} define
\begin{equation}\label{eq:sepsections}
F\bigl(x_{\infty},(w_q)\bigr):=h(x_{\infty})\cdot
\mathbf 1_{u(1+p^{m}\Z_p)}(w_p)\cdot\prod_{q\ne p}\mathbf 1_{\Z_q^{\times}}(w_q),
\end{equation}
and $F'$ by the same formula with $u$ replaced by $u'$.  Then $F,F'\in
\DA^{\mathrm{off}}$ and, with $\nu:=h(1)\vol(1+p^{m}\Z_p)>0$,
\begin{equation}\label{eq:sepvalues}
T(F)=\nu(\log p)p^{\,j},\qquad T(F')=\nu(\log p)p^{\,j'},\qquad
T(F)-T(F')=\nu(\log p)\bigl(p^{\,j}-p^{\,j'}\bigr)\neq0,
\end{equation}
while $F$ and $F'$ have the same pushforward under
$|\cdot|_S:C_S\to\R_{>0}$; in particular $\lambda_-(F)=\lambda_-(F')$ and
$\lambda_+(F)=\lambda_+(F')$.
\end{theorem}

\begin{proof}
\emph{Membership.}  $F$ and $F'$ are finite products of a compactly supported
smooth function of $x_{\infty}$ with characteristic functions of compact open
subsets of the compact groups $\Z_q^{\times}$, hence lie in $\DA$.  The
identity class is $1=(1,(1)_q)$; since $u\in1+p\Z_p$ with $v_p(1-u)=j<m$, we
have $1\notin u(1+p^{m}\Z_p)$, and that set is open and closed, so $F$ vanishes
on a neighbourhood of $1$.  By Theorem~\ref{thm:absconv} the value $T(F)$ is
the sum \eqref{eq:absconv}.

\emph{The archimedean term vanishes.}  By \eqref{eq:slices} the archimedean
slice is $j_{\infty}(x)=(x,1,\dots,1)$, whose $p$-coordinate is
$1\notin u(1+p^{m}\Z_p)$; hence $F\circ j_{\infty}\equiv0$ and $W_{\infty}(F)=0$.

\emph{The other finite terms vanish.}  Let $q\in S_f$, $q\ne p$.  By
\eqref{eq:slices} the $p$-coordinate of $j_q(q^{k}w)$ is $q^{-k}$.  For $k=0$
this is $1\notin u(1+p^{m}\Z_p)$, so $F\circ j_q$ vanishes on the units.  For
$k\ne0$ the
archimedean coordinate of $j_q(q^kw)$ is $q^{-k}$, which is not in $\supp h$
because $|k|\ge1$ and $\supp h\subset(\theta^{-1},\theta)$ with
$\theta\le q^{1/2}<q$.  So $F\circ j_q\equiv0$ and $W_q(F)=0$.

\emph{The term at $p$.}  By \eqref{eq:slices},
$a(p^{k}w)=h(p^{-k})\mathbf 1_{u(1+p^{m}\Z_p)}(w)$.  For $k\ne0$ we have
$p^{-k}\notin\supp h$, so the shell integral vanishes.  For $k=0$,
$a(w)=h(1)\mathbf 1_{u(1+p^{m}\Z_p)}(w)$ and $a(1)=0$, so by \eqref{eq:Wp}
\[
W_p(F)=(\log p)\,h(1)\int_{u(1+p^{m}\Z_p)}|1-w|_p^{-1}\,dw .
\]
On the coset write $w=u(1+p^{m}y)$ with $y\in\Z_p$; then $1-w=(1-u)-up^{m}y$,
and since $v_p(1-u)=j<m$ we get $v_p(1-w)=j$, so $|1-w|_p^{-1}=p^{j}$
is constant on the coset.  Therefore
$W_p(F)=(\log p)h(1)p^{j}\vol(1+p^{m}\Z_p)=\nu(\log p)p^{j}$, with $\nu>0$
because $h(1)>0$.  The same computation with $u'$, and subtraction, gives
\eqref{eq:sepvalues}; the difference is nonzero because $j\ne j'$.

\emph{The pushforwards agree.}  By Lemma~\ref{thm:split}, $|w|_S=|w_\infty|$
on $C_S$, and $h$ is supported in $\R_{>0}$, so the pushforward of
$F\dstar w$ under $|\cdot|_S$ is
$h(x)\,\vol\bigl(u(1+p^{m}\Z_p)\bigr)\prod_{q\ne p}\vol(\Z_q^{\times})\,
dx/(2x)$, and $\vol(u(1+p^{m}\Z_p))=\vol(1+p^{m}\Z_p)=\vol(u'(1+p^{m}\Z_p))$ by
translation invariance of Haar measure on $\Z_p^{\times}$.  Since
$\lambda_{\mp}(F)$ is the integral of $|x|^{\mp1/2}$ against this pushforward,
the values of $\lambda_{\mp}$ agree as well.
\end{proof}

\begin{proposition}[The local weight and the coinvariants]\label{prop:haarrank}
Let $p\in S_f$ and let $x\in\Z_p^{\times}$, $x\ne1$.  Let $H_x$ be the closure
of the subgroup of $\Z_p^{\times}$ generated by $x$, and let
$N(x):=[\Z_p^{\times}:H_x]\in\N\cup\{\infty\}$, which by
Lemma~\ref{lem:coinv} is the dimension of the space of coinvariants
$F^{(j)}/\overline{(1-\sigma)F^{(j)}}$ of the dilation $\sigma$ by $x$ on each
shell.  Put $q_p=p$ for odd $p$ and $q_2=4$.
\begin{enumerate}
\item[(a)] If $x\in1+q_p\Z_p$ and $k:=v_p(1-x)$, then $N(x)=(p-1)p^{k-1}$ and
\[
\tau_p(x)=|1-x|_p^{-1}=\frac{N(x)}{1-p^{-1}} .
\]
\item[(b)] If $x\notin1+q_p\Z_p$, the relation of (a) fails in general.  For
$p$ odd one has $\tau_p(x)=1$, which is not of the form $n/(1-p^{-1})$ with
$n\in\N\cup\{\infty\}$.  For $p=2$ and $x=3$ one has $\tau_2(3)=2$ and
$N(3)=2$, so $N(3)/(1-2^{-1})=4$.  For $x=-1$ and every $p$ one has
$N(-1)=\infty$.
\end{enumerate}
\end{proposition}

\begin{proof}
(a)  The $p$-adic logarithm is an isomorphism of topological groups
$1+q_p\Z_p\to q_p\Z_p$ with $v_p(\log x)=v_p(x-1)=k$.  It carries $H_x$ to the
closure of $\Z\log x$, which is $p^{k}\Z_p$, so $H_x=1+p^{k}\Z_p$ and
$N(x)=|(\Z/p^{k})^{\times}|=(p-1)p^{k-1}=(1-p^{-1})p^{k}$, while
$\tau_p(x)=p^{k}$ by Lemma~\ref{lem:tau}(ii).

(b)  For $p$ odd and $x\notin1+p\Z_p$ one has $|1-x|_p=1$, so $\tau_p(x)=1$,
and $n/(1-p^{-1})=np/(p-1)$ is never $1$ for $n\in\N\cup\{\infty\}$.  For $p=2$
and $x=3$, $|1-3|_2^{-1}=2$; and $9\in1+8\Z_2$ with $v_2(1-9)=3$, so by (a) the
closure of $\langle9\rangle$ is $1+8\Z_2$ and $H_3=\{1,3\}\cdot(1+8\Z_2)$, of
index $2$ in $\Z_2^{\times}$.  Finally $H_{-1}=\{\pm1\}$ is finite, hence of
infinite index in $\Z_p^{\times}$.
\end{proof}

The trace of Definition~\ref{def:trace} is defined by the densities $\tau_v$
directly.  By Proposition~\ref{prop:haarrank}, on the units of $1+q_p\Z_p$ the
density equals the number $N(x)$ of Lemma~\ref{lem:coinv} divided by
$1-p^{-1}$; on the remaining units, which contribute to the unit integral of
\eqref{eq:Wp}, it does not.  At a non-unit $x$ the dilation permutes the shells
and $\tau_p(x)=p^{-|v_p(x)|/2}$ by Lemma~\ref{lem:tau}(i).

\section{The Tate characters and the symbol of the trace}\label{sec:symbol}

\subsection{The two Tate characters}

\begin{theorem}\label{thm:tate}
The map
\begin{equation}\label{eq:rhoT}
\rho_T(u):=\bigl(|u|_S^{-1/2},\,|u|_S^{+1/2}\bigr)
\end{equation}
is well defined on the class group $C_S$, independent of the lift, and
satisfies $\rho_T(u^{-1})=\rho_T(u)^{\sharp}$, where $\sharp$ is the
conjugate-linear swap
$(z_-,z_+)^{\sharp}:=(\overline{z_+},\overline{z_-})$ on
$\BT:=\C_-\oplus\C_+$.  The two \emph{Tate functionals}
\begin{equation}\label{eq:lambdadef}
\lambda_{\mp}(F):=\int_{C_S}F(u)\,|u|_S^{\mp1/2}\dstar u,\qquad F\in\DA,
\end{equation}
are continuous and satisfy
$\lambda_{\mp}(F\star G)=\lambda_{\mp}(F)\lambda_{\mp}(G)$ and
$\lambda_{\mp}(F^{\dagger})=\overline{\lambda_{\pm}(F)}$.  Consequently
$M:=(\lambda_-,\lambda_+):\DA\to\BT$ is a continuous algebra homomorphism with
$M(F^{\dagger})=M(F)^{\sharp}$, and it is surjective.  The Hermitian form
$\beta_T(z,w):=z_-\overline{w_+}+z_+\overline{w_-}$ on $\BT$ is nondegenerate;
in the basis $e_-=(1,0)$, $e_+=(0,1)$ its matrix is
$\left(\begin{smallmatrix}0&1\\1&0\end{smallmatrix}\right)$.
\end{theorem}

\begin{proof}
$|\gamma|_S=1$ for $\gamma\in\Gamma_S$ by Lemma~\ref{lem:prodformula}, so
$|\cdot|_S^{\pm1/2}$ is constant on $\Gamma_S$-orbits and \eqref{eq:rhoT}
descends to $C_S$; its values are real, so
$\rho_T(u^{-1})=(|u|_S^{1/2},|u|_S^{-1/2})=\rho_T(u)^{\sharp}$.  On the
functions supported in a fixed compact set $K$ one has
$|\lambda_{\mp}(F)|\le\sup|F|\int_K|u|_S^{\mp1/2}\dstar u$, which gives
continuity.  For $s=\mp\tfrac12$, Fubini's theorem and the substitution $w=uv$
give
\[
\int_{C_S}\int_{C_S}F(u)G(u^{-1}w)\,|w|_S^{s}\dstar u\dstar w
=\int_{C_S}F(u)|u|_S^{s}\dstar u\int_{C_S}G(v)|v|_S^{s}\dstar v ,
\]
the integrand being continuous with compact support; this is
multiplicativity.  For the involution, substituting $v=u^{-1}$ and using
inversion invariance,
\[
\lambda_{\mp}(F^{\dagger})=\int\overline{F(u^{-1})}|u|_S^{\mp1/2}\dstar u
=\overline{\int F(v)|v|_S^{\pm1/2}\dstar v}=\overline{\lambda_{\pm}(F)},
\]
which is $M(F^{\dagger})=M(F)^{\sharp}$.  Surjectivity: for
$f\in C_c^{\infty}(\R)$ the function $F(w):=f(\log|w|_S)$ lies in $\DA$, and
$\lambda_{\mp}(F)=\int_{\R}f(t)e^{\mp t/2}dt$ since $\dstar w=dt\otimes
\mu_{\mathrm{cpt}}$.  If $a\lambda_-+b\lambda_+=0$ on $\DA$, then
$\int f(t)(ae^{-t/2}+be^{t/2})dt=0$ for all $f$, so $ae^{-t/2}+be^{t/2}=0$ and
$a=b=0$.  Thus the image of $M$ is a subspace of $\C^{2}$ annihilated by no
nonzero linear form, that is, all of $\C^{2}$.  The determinant of the matrix
of $\beta_T$ is $-1$, so $\beta_T$ is nondegenerate.
\end{proof}

The functionals $\lambda_{\mp}$ are Mellin transforms at the exponents
$\mp\tfrac12$, of the kind studied in \cite{Tate}.

\subsection{The dual group}

By Lemma~\ref{thm:split}, $C_S\cong\R^{\times}\times\prod_{q\in
S_f}\Z_q^{\times}$, so its Pontryagin dual is
\begin{equation}\label{eq:dual}
\Ch=\bigl\{\chi=(\varepsilon,\tau;(\kappa_q))\bigr\},\qquad
\chi(w)=\sgn(w_{\infty})^{\varepsilon}|w_{\infty}|^{i\tau}
\prod_{q\in S_f}\kappa_q(w_q),
\end{equation}
with $\varepsilon\in\{0,1\}$, $\tau\in\R$ and $\kappa_q$ a finite-order
character of $\Z_q^{\times}$ of conductor exponent $\cond(\kappa_q):=\cond_q(\kappa_q)\ge0$, in the notation of
Section~\ref{sec:carrier}.  As a
topological space $\Ch$ is a disjoint union of copies of $\R$, one for each of
the countably many pairs $(\varepsilon,(\kappa_q))$; we call these its
\emph{components}, and we write $\mathfrak c_0$ for the \emph{principal
component}, $\varepsilon=0$ and $\kappa_q=1$ for all $q$.  Set
\begin{equation}\label{eq:FT}
\widehat F(\chi):=\int_{C_S}F(w)\overline{\chi(w)}\dstar w .
\end{equation}
If $F\in\DA$ is invariant under $\prod_q(1+q^{N}\Z_q)$ then $\widehat F$ is
supported on the finitely many components with $\cond(\kappa_q)\le N$ for all $q$,
and on each of those it is the Fourier transform
$\widehat F(\mathfrak c,\tau)=\int_{\R}f_{\mathfrak c}(t)e^{-i\tau t}dt$ of a
function $f_{\mathfrak c}\in C_c^{\infty}(\R)$ of $t=\log|w|_S$, namely the
integral of $F\overline{\chi}$ over the compact factor.  Hence $\widehat
F(\mathfrak c,\cdot)$ is a Schwartz function of $\tau$ and, by the
Paley--Wiener theorem \cite[\S7.3]{Hormander}, extends to an entire function
of $\tau$ with
\begin{equation}\label{eq:PW}
|\widehat F(\mathfrak c,\sigma+iy)|\le C_n\,e^{R|y|}(1+|\sigma|)^{-n}
\qquad(n\ge0),
\end{equation}
where $\supp f_{\mathfrak c}\subset[-R,R]$.  The two Tate functionals are
values of this continuation off the real axis:
\begin{equation}\label{eq:lambdaFT}
\lambda_{\mp}(F)=\widehat F\bigl(\varepsilon=0,\ \tau=\mp i/2;\ (1)\bigr),
\end{equation}
because $|w|_S^{-i\tau}=|w|_S^{\mp1/2}$ at $\tau=\mp i/2$.  They are not
evaluations at points of $\Ch$.  We normalize
$d\chi|_{\mathfrak c}=d\tau/2\pi$ on each component, the Plancherel measure
dual to $\dstar w$.

\subsection{The symbol}

\begin{theorem}\label{thm:symbol}
For $\chi\in\Ch$ put
\begin{equation}\label{eq:symboldef}
\Th(\chi):=\Wh_{\infty}(\chi)+\sum_{p\in S_f}\Wh_p(\chi),\qquad
\Wh_v(\chi):=\int'_{\Q_v^{\times}}\chi(j_v(x))\,\tau_v(x)\dstar x,
\end{equation}
with the regularization and the measures of Definition~\ref{def:trace}.  Then
each integral converges and:
\begin{enumerate}
\item[(a)] \emph{(finite places)}  Put
$\eta_p(\chi):=p^{-i\tau}\prod_{q\in S_f,\,q\ne p}\kappa_q(p)^{-1}$, of modulus
one, and $\zeta_p(\chi):=p^{-1/2}\overline{\eta_p(\chi)}$.  Then
\begin{equation}\label{eq:Wphat}
\Wh_p(\chi)=
\begin{cases}
2(\log p)\,\Rel\bigl(\zeta_p/(1-\zeta_p)\bigr), & \kappa_p=1,\\[2pt]
-\cond(\kappa_p)\log p, & \kappa_p\neq1 .
\end{cases}
\end{equation}
In both cases $\Wh_p$ is real, real-analytic in $\tau$, and bounded on each
component by $2(\log p)p^{-1/2}(1-p^{-1/2})^{-1}$, respectively by
$\cond(\kappa_p)\log p$.
\item[(b)] \emph{(the real place)}  $\Wh_{\infty}$ is real, real-analytic in
$\tau$ on each component, bounded above, and
\begin{equation}\label{eq:Winfhat}
\Wh_{\infty}(\varepsilon,\tau)=-\log|\tau|+O(1)\qquad(|\tau|\to\infty),
\end{equation}
uniformly in $\varepsilon$.  For $\varepsilon=0$ one has the closed form
\begin{equation}\label{eq:Winfclosed}
\Wh_{\infty}(0,\tau)=\log\pi-\Rel\,\psi\bigl(\tfrac14+\tfrac{i\tau}{2}\bigr),
\end{equation}
$\psi$ the digamma function.  For $\varepsilon=1$, with $W_{\infty}$ as in
Convention~\ref{conv:Winf}, $\Wh_{\infty}(1,\tau)=\log\pi-\Rel\,\psi(\tfrac34+
\tfrac{i\tau}2)=\Wh_{\infty}(0,\tau)-\pi/\cosh(\pi\tau)$, see
\eqref{eq:Winfodd}.
\item[(c)] $\Th$ is real-valued, real-analytic on each component, bounded above
by a finite constant $m:=\sup_{\Ch}\Th<\infty$, and
$\Th(\varepsilon,\tau;(\kappa_q))\to-\infty$ as $|\tau|\to\infty$.  Hence for
every $c\in\R$ the function $\Th+c$ vanishes identically on no component, and
its zero set is a discrete subset of each component.
\item[(d)] \emph{(Plancherel form)}  For every $H\in\DA$,
\begin{equation}\label{eq:plancherel}
T(H)=\int_{\Ch}\widehat H(\chi)\bigl(\Th(\chi)+c\bigr)d\chi ,
\end{equation}
the integral being absolutely convergent; and the interchange is justified
termwise, so that for every $v\in S$,
\begin{equation}\label{eq:plancherelv}
W_v(H)=\int_{\Ch}\widehat H(\chi)\,\Wh_v(\chi)\,d\chi ,
\end{equation}
also absolutely convergent.
\end{enumerate}
\end{theorem}

\begin{proof}
(a)  By \eqref{eq:slices}, $\sgn(p^{-k})=+1$ and
\[
\chi\bigl(j_p(p^{k}w)\bigr)=p^{-ik\tau}\kappa_p(w)\prod_{q\ne p}\kappa_q(p)^{-k}
=\eta_p(\chi)^{k}\kappa_p(w).
\]
Insert this into \eqref{eq:Wp} with $a=\chi\circ j_p$, which has $a(1)=1$.

\emph{Shell part.}  By Lemma~\ref{lem:tau}(i) and the shell volume $\log p$,
\[
\int_{p^{k}\Z_p^{\times}}a\,\tau_p\dstar x
=(\log p)\,p^{-|k|/2}\,\eta_p^{k}\int_{\Z_p^{\times}}\kappa_p(w)\,dw ,
\]
and the unit integral is $1$ if $\kappa_p=1$ and $0$ otherwise, by
orthogonality.  For $\kappa_p=1$ the sum over
$k\ne0$ is $\sum_{k\ge1}p^{-k/2}(\eta_p^{k}+\eta_p^{-k})$; with
$\zeta:=p^{-1/2}\overline{\eta_p}$ and $\eta_p^{-1}=\overline{\eta_p}$ this is
$\sum_{k\ge1}(\bar\zeta^{k}+\zeta^{k})=2\Rel(\zeta/(1-\zeta))$, absolutely
convergent since $|\zeta|=p^{-1/2}<1$, and bounded in absolute value by
$2p^{-1/2}(1-p^{-1/2})^{-1}$.

\emph{Unit integral.}  For $k=0$, $j_p(u)=(1,u,1,\dots)$, so the second term
of \eqref{eq:Wp} is
$(\log p)\int_{\Z_p^{\times}}(\kappa_p(u)-1)|1-u|_p^{-1}du$ against the
probability Haar measure.  We claim
\begin{equation}\label{eq:conductor}
\int_{\Z_p^{\times}}\frac{\kappa_p(u)-1}{|1-u|_p}\,du=-\cond(\kappa_p).
\end{equation}
We may assume $e:=\cond(\kappa_p)\ge1$, the case $\kappa_p=1$ being
trivial.  Decompose $\Z_p^{\times}$ into the sets
$V_n:=\{u\in\Z_p^{\times}:v_p(1-u)=n\}$, $n\ge0$, on which $|1-u|_p^{-1}=p^{n}$;
one has $V_0=\Z_p^{\times}\setminus(1+p\Z_p)$ and
$V_n=(1+p^{n}\Z_p)\setminus(1+p^{n+1}\Z_p)$ for $n\ge1$.  Write
$\mu_n:=\vol(1+p^{n}\Z_p)=\bigl((p-1)p^{n-1}\bigr)^{-1}$ for $n\ge1$ and
$\mu_0:=\vol(\Z_p^{\times})=1$, and put
\[
A_n:=\int_{1+p^{n}\Z_p}\bigl(\kappa_p(u)-1\bigr)du\quad(n\ge1),\qquad
A_0:=\int_{\Z_p^{\times}}\bigl(\kappa_p(u)-1\bigr)du,
\]
so that $\int_{V_n}(\kappa_p-1)\,du=A_n-A_{n+1}$ for every $n\ge0$.  Since
$\kappa_p$ is trivial on $1+p^{n}\Z_p$ if and only if $n\ge e$, orthogonality
of characters on the group $1+p^{n}\Z_p$ (respectively on $\Z_p^{\times}$ when
$n=0$, using $e\ge1$) gives $\int_{1+p^n\Z_p}\kappa_p\,du=\mu_n$ for $n\ge e$
and $0$ for $n<e$.  Hence $A_n=0$ for $n\ge e$ and $A_n=-\mu_n$ for $n<e$.
Therefore
\[
\int_{\Z_p^{\times}}\frac{\kappa_p(u)-1}{|1-u|_p}\,du
=\sum_{n\ge0}p^{n}\bigl(A_n-A_{n+1}\bigr)
=-\sum_{n=0}^{e-1}p^{n}\mu_n+\sum_{n=0}^{e-2}p^{n}\mu_{n+1}
=-\mu_0+\sum_{m=1}^{e-1}\bigl(p^{m-1}-p^{m}\bigr)\mu_m .
\]
For $m\ge1$ one has $(p^{m-1}-p^{m})\mu_m=p^{m-1}(1-p)/\bigl((p-1)p^{m-1}\bigr)
=-1$, so the right-hand side is $-1-(e-1)=-e$, which is \eqref{eq:conductor}.
The integrand of \eqref{eq:conductor} vanishes on $1+p^{e}\Z_p$ and is bounded
by $2p^{e}$ elsewhere.

For $\kappa_p=1$ the unit integral vanishes, since the integrand is zero, and
for $\kappa_p\ne1$ the shell part vanishes.  Combining the two parts gives
\eqref{eq:Wphat}.  Finally
$\zeta_p=p^{-1/2}e^{i(\tau\log p+\theta_p)}$ with
$\theta_p:=\arg\prod_{q\ne p}\kappa_q(p)$ a constant, so
$2\Rel(\zeta/(1-\zeta))$ is a real-analytic $2\pi/\log p$-periodic function of
$\tau$; the ramified case is constant.

(b)  By \eqref{eq:slices}, $\chi(j_{\infty}(x))=\sgn(x)^{\varepsilon}
|x|^{i\tau}$, so the two profiles of Definition~\ref{def:trace} are
$b_{+}(t)=e^{i\tau t}$ and $b_{-}(t)=(-1)^{\varepsilon}e^{i\tau t}$.  Put
these into \eqref{eq:Winf} (Convention~\ref{conv:Winf}).  The $K_{c}$-integral
is $(-1)^{\varepsilon}\int_0^{\infty}K_c(t)\,2\cos(\tau t)\,dt$, and the
remaining terms form the regularized $K_{s}$-integral, which is the same for
both parities; since $2K_{s}(|t|)=1/(2|\sinh(t/2)|)$ and
$2K_{c}(|t|)=1/(2\cosh(t/2))$, this reads
\begin{equation}\label{eq:Winfhat2}
\Wh_{\infty}(\varepsilon,\tau)=\frac12\int'_{\R}
\frac{e^{i\tau t}\,dt}{2|\sinh(t/2)|}
+(-1)^{\varepsilon}\cdot\frac12\int_{\R}\frac{e^{i\tau t}\,dt}{2\cosh(t/2)} ,
\end{equation}
where, by definition,
\[
\frac12\int'_{\R}\frac{e^{i\tau t}\,dt}{2|\sinh(t/2)|}
:=(\gamma+\log4\pi)+\int_0^{\infty}\bigl[2K_{s}(t)\cos\tau t-2K(t)e^{-t/2}\bigr]dt,
\]
\[
\frac12\int_{\R}\frac{e^{i\tau t}\,dt}{2\cosh(t/2)}=\int_0^{\infty}2K_{c}(t)\cos\tau t\,dt .
\]
Both right-hand sides are real, and both integrals converge absolutely, as in
Convention~\ref{conv:Winf}.  One has
\begin{equation}\label{eq:sech}
\int_{\R}\frac{e^{i\tau t}\,dt}{2\cosh(t/2)}=\frac{\pi}{\cosh(\pi\tau)}
\end{equation}
\cite[\S3.981]{GR}, so the second term of \eqref{eq:Winfhat2} is
$(-1)^{\varepsilon}\pi/(2\cosh\pi\tau)$; this is the only place where the
parity enters.

\emph{The closed form.}  For $\varepsilon=0$ the two profiles coincide,
$b_{+}=b_{-}=b$ with $b(t)=e^{i\tau t}$, and $b(0)=1$; so the even form
\eqref{eq:Winfeven} gives
\[
\Wh_{\infty}(0,\tau)=(\gamma+\log4\pi)+\int_0^{\infty}K(t)
\bigl[2\cos\tau t-2e^{-t/2}\bigr]dt,\qquad
K(t)=\frac{e^{-t/2}}{1-e^{-2t}} .
\]
Substituting $u=2t$ turns the integral into
\[
\mathcal J(\tau):=\int_0^{\infty}
\frac{e^{-u/4}\cos(\tau u/2)-e^{-u/2}}{1-e^{-u}}\,du
=\int_0^{\infty}\frac{\tfrac12\bigl(e^{-\alpha u}+e^{-\bar\alpha u}\bigr)
-e^{-u/2}}{1-e^{-u}}\,du,\qquad \alpha:=\tfrac14+\tfrac{i\tau}{2},
\]
the integrand being continuous at $u=0$ and exponentially small at infinity,
so that the integral converges absolutely.  Gauss's integral representation
of the digamma function \cite[\S5.9]{DLMF},
\[
\psi(s)=-\gamma+\int_0^{\infty}\frac{e^{-u}-e^{-su}}{1-e^{-u}}\,du
\qquad(\Rel s>0),
\]
gives, for $\Rel s>0$,
$\int_0^{\infty}\bigl(e^{-su}-e^{-u/2}\bigr)(1-e^{-u})^{-1}du
=\psi(\tfrac12)-\psi(s)$, by subtracting the two representations.  Applying
this with $s=\alpha$ and with $s=\bar\alpha$, and halving,
\[
\mathcal J(\tau)=\psi(\tfrac12)-\Rel\,\psi\bigl(\tfrac14+\tfrac{i\tau}{2}\bigr),
\qquad \psi(\tfrac12)=-\gamma-2\log2
\]
\cite[\S5.4]{DLMF}.  Hence $\Wh_{\infty}(0,\tau)=\gamma+\log4\pi-\gamma-2\log2-
\Rel\,\psi(\tfrac14+\tfrac{i\tau}{2})=\log\pi-
\Rel\,\psi(\tfrac14+\tfrac{i\tau}{2})$, which is \eqref{eq:Winfclosed}.
Compare the archimedean term of Weil's explicit formula
\cite{Weil1952,CCWeilArch}.

\emph{The odd component.}  Subtracting the two cases of \eqref{eq:Winfhat2}
and using \eqref{eq:sech},
$\Wh_{\infty}(1,\tau)=\Wh_{\infty}(0,\tau)-\pi/\cosh(\pi\tau)$.  The reflection
formula $\psi(1-z)-\psi(z)=\pi\cot\pi z$ \cite[\S5.5]{DLMF} at
$z=\tfrac14+\tfrac{i\tau}2$, together with
$\Rel\cot(\tfrac\pi4+\tfrac{i\pi\tau}2)=\sin(\tfrac\pi2)/(\cosh(\pi\tau)-\cos\tfrac\pi2)
=1/\cosh(\pi\tau)$ and $\Rel\psi(\bar w)=\Rel\psi(w)$, gives
$\Rel\psi(\tfrac34+\tfrac{i\tau}2)-\Rel\psi(\tfrac14+\tfrac{i\tau}2)=\pi/\cosh(\pi\tau)$.
Hence
\begin{equation}\label{eq:Winfodd}
\Wh_{\infty}(1,\tau)=\log\pi-\Rel\,\psi\bigl(\tfrac34+\tfrac{i\tau}{2}\bigr).
\end{equation}
Both closed forms are real-analytic in $\tau$.

\emph{The asymptotic.}  By \cite[\S5.11]{DLMF},
$\psi(z)=\log z-\tfrac1{2z}+O(|z|^{-2})$ as $z\to\infty$ in
$|\arg z|\le\pi/2$.  For $a\in\{\tfrac14,\tfrac34\}$ and $z=a+\tfrac{i\tau}2$
one has $\Rel\log z=\tfrac12\log(a^{2}+\tau^{2}/4)=\log|\tau/2|+O(\tau^{-2})$
and $\Rel(1/z)=a/|z|^{2}=O(\tau^{-2})$, so
$\Wh_{\infty}(\varepsilon,\tau)=-\log|\tau|+\log2\pi+O(\tau^{-2})$, which
gives \eqref{eq:Winfhat}.

\emph{Monotonicity and the maximum.}  By the partial fraction expansion
$\psi(z)=-\gamma+\sum_{n\ge0}\bigl[\tfrac1{n+1}-\tfrac1{n+z}\bigr]$
\cite[\S5.7]{DLMF}, locally uniformly convergent off $-\N_0$,
\[
\Rel\,\psi\bigl(\tfrac14+\tfrac{i\tau}{2}\bigr)
=-\gamma+\sum_{n\ge0}\Bigl[\frac1{n+1}
-\frac{n+\tfrac14}{(n+\tfrac14)^{2}+\tau^{2}/4}\Bigr].
\]
Each subtracted term is a positive decreasing function of $|\tau|$, so the sum
is increasing in $|\tau|$; hence $\Rel\,\psi(\tfrac14+\tfrac{i\tau}{2})$ is
increasing in $|\tau|$ and $\Wh_{\infty}(0,\cdot)$ attains its maximum at
$\tau=0$:
\begin{equation}\label{eq:maxarch}
\max_{\tau}\Wh_{\infty}(0,\tau)=\log\pi-\psi(\tfrac14)
=\log\pi+\gamma+3\log2+\tfrac{\pi}{2}\approx5.372 ,
\end{equation}
with the value $\psi(\tfrac14)=-\gamma-3\log2-\tfrac{\pi}{2}$ from
\cite[\S5.4]{DLMF}.  Since
$\Wh_{\infty}(1,\tau)=\Wh_{\infty}(0,\tau)-\pi/\cosh(\pi\tau)\le
\Wh_{\infty}(0,\tau)$, the same bound holds for $\varepsilon=1$, and
$\Wh_{\infty}\le\log\pi-\psi(\tfrac14)$ on all of $\Ch$.

(c)  By (a), $\Wh_p\le2(\log p)p^{-1/2}(1-p^{-1/2})^{-1}$ on all of $\Ch$, the
ramified values being $\le0$, and by (b) $\Wh_{\infty}\le\log\pi-\psi(\tfrac14)$
on all of $\Ch$; since $S_f$ is finite, $m<\infty$.  On each component every
$\Wh_p$ is real, real-analytic and bounded, and $\Wh_{\infty}$ is real,
real-analytic and tends to $-\infty$; so $\Th$ is real, real-analytic, and
tends to $-\infty$ along every component.  Hence $\Th+c$ is not the zero
function on any component, and a real-analytic function on $\R$ which is not
identically zero has isolated zeros.

(d)  Let $H\in\DA$.  By Fourier inversion on the locally compact abelian group
$C_S$, $H(w)=\int_{\Ch}\widehat H(\chi)\chi(w)\,d\chi$, where $\widehat H$ is
supported on finitely many components and is a Schwartz function of $\tau$ on
each.  In particular $H(1)=\int\widehat H\,d\chi$, which accounts for the term
$cH(1)$.  Insert the inversion formula into \eqref{eq:Wp} and
\eqref{eq:Winf} and exchange the $\chi$-integral with the integrals over
$\Q_v^{\times}$; this gives \eqref{eq:plancherelv}, and summing over $v$ gives
\eqref{eq:plancherel}.  The exchange is justified by Fubini's theorem as
follows.  At a finite place the shell integrands are bounded by
$|\widehat H(\chi)|\,p^{-|k|/2}$ on the shell $p^{k}\Z_p^{\times}$, summable in $k$ and integrable in $\chi$, and
the integrand of the unit integral is bounded by
$2p^{\cond(\kappa_p)}|\widehat H(\chi)|$, as noted after \eqref{eq:conductor},
where $\cond(\kappa_p)$ is bounded on the support of $\widehat H$.  At the
archimedean place the integrand in $(t,\tau)$ is $\widehat H(\chi)$ times
\[
2K_{s}(t)(\cos\tau t-1)+\bigl(2K_{s}(t)-2K(t)e^{-t/2}\bigr)
\pm2K_{c}(t)\cos\tau t .
\]
Here $|\cos\tau t-1|\le\min(2,\tau^{2}t^{2}/2)$, $2K_{s}(t)\le t^{-1}$ for
$t>0$ and $2K_{s}(t)\le2e^{-t/2}$ for $t\ge1$, so
$\int_0^{\infty}2K_{s}(t)|\cos\tau t-1|\,dt\le C(1+\log(2+|\tau|))$; the
functions $2K_{s}(t)-2K(t)e^{-t/2}=(2\sinh(t/2))^{-1}-(\sinh t)^{-1}$ and
$2K_{c}$ are integrable on $(0,\infty)$.  Thus the absolute value of the
integrand has $t$-integral at most $C'(1+\log(2+|\tau|))|\widehat H(\chi)|$,
which is integrable over $\Ch$.  The same bounds, together with the
boundedness of $\Wh_p$ on each component and \eqref{eq:Winfhat}, show that the
right-hand sides of \eqref{eq:plancherel} and \eqref{eq:plancherelv} converge
absolutely.
\end{proof}

\begin{remark}\label{rem:symbolshape}
Expanding \eqref{eq:Wphat} at an unramified place gives
\begin{equation}\label{eq:symbolexp}
\Wh_p(\chi)=\sum_{k\ge1}2(\log p)\,p^{-k/2}
\cos\bigl(k\tau\log p+k\theta_p\bigr),
\end{equation}
with $\theta_p=\arg\prod_{q\ne p}\kappa_q(p)$.  The weights $p^{-k/2}$ are the
values $\tau_p(p^{\pm k})$ of Lemma~\ref{lem:tau}, and $\log p$ is the shell
volume of Definition~\ref{def:trace}.
\end{remark}

\subsection{Mittag--Leffler form and the Tate poles}

\begin{theorem}\label{thm:ML}
Fix a component $\mathfrak c=(\varepsilon,(\kappa_q))$ of $\Ch$ and write
$\tau$ for its real coordinate.
\begin{enumerate}
\item[(a)] Let $p\in S_f$ with $\kappa_p=1$, and put $L:=\log p$ and
$\theta_p:=\arg\prod_{q\ne p}\kappa_q(p)$.  Then $\Wh_p(\mathfrak c,\cdot)$
extends from $\R$ to a meromorphic function on $\C$, $2\pi/L$-periodic in
$\tau$, given by
\begin{equation}\label{eq:MLp}
\Wh_p(\tau)=-L-i\sum_{k\in\Z}
\Bigl[\frac{1}{\tau-\tau_k^{+}}-\frac{1}{\tau-\tau_k^{-}}\Bigr],
\qquad \tau_k^{\pm}:=\pm\frac{i}{2}+\frac{2\pi k-\theta_p}{L},
\end{equation}
the series converging absolutely and locally uniformly off the poles.
All poles are simple; the residue at $\tau_k^{+}$ is $-i$ and at $\tau_k^{-}$
is $+i$, for every $k$ and every $p$.  For $\kappa_p\ne1$ the function
$\Wh_p(\mathfrak c,\cdot)$ is the constant $-\cond(\kappa_p)\log p$.
\item[(b)] With $\omega_n:=i(2n+\tfrac12)$, $n\ge0$,
\begin{equation}\label{eq:MLinf}
\Wh_{\infty}(0,\tau)=\log\pi+\gamma+\sum_{n\ge0}
\Bigl[\frac{-i}{\tau-\omega_n}+\frac{i}{\tau+\omega_n}-\frac{1}{n+1}\Bigr],
\end{equation}
locally uniformly convergent off the poles; $\Wh_{\infty}(0,\cdot)$ is
meromorphic on $\C$ with simple poles at $\pm\omega_n$ and residues $\mp i$.
\item[(c)] On the principal component $\mathfrak c_0$ one has $\theta_p=0$ for
every $p$, so $\tau_0^{\pm}=\pm i/2=\pm\omega_0$: every place of $S$
contributes a simple pole at $\tau=\pm i/2$ with residue $\mp i$, and
\begin{equation}\label{eq:residues}
\Res_{\tau=+i/2}\Th(\mathfrak c_0,\cdot)=-i\bigl(|S_f|+1\bigr),\qquad
\Res_{\tau=-i/2}\Th(\mathfrak c_0,\cdot)=+i\bigl(|S_f|+1\bigr).
\end{equation}
\item[(d)] Let $f\in C_c^{\infty}(\R)$ and $F(w):=f(\log|w|_S)$, so that
$F\in\DA$, $\widehat F$ is supported on $\mathfrak c_0$, and
$\widehat F(\tau):=\widehat F(\mathfrak c_0,\tau)=\int f(t)e^{-i\tau t}dt$.
Then for every $h>\tfrac12$
\begin{equation}\label{eq:contourp}
W_p(F)=\sum_{k\in\Z}\widehat F\Bigl(\frac{i}{2}+\frac{2\pi k}{\log p}\Bigr)
+\frac{1}{2\pi}\int_{\Im\tau=h}\widehat F(\tau)\Wh_p(\tau)\,d\tau ,
\end{equation}
both terms absolutely convergent, and the sum is the Poisson periodization
\begin{equation}\label{eq:poisson}
\sum_{k\in\Z}\widehat F\Bigl(\frac i2+\frac{2\pi k}{\log p}\Bigr)
=(\log p)\sum_{m\in\Z}f(m\log p)\,p^{m/2}.
\end{equation}
\item[(e)] For the same $F$ and any $\tfrac12<h<\tfrac52$,
\begin{equation}\label{eq:contourinf}
W_{\infty}(F)=\widehat F(i/2)
+\frac{1}{2\pi}\int_{\Im\tau=h}\widehat F(\tau)\Wh_{\infty}(0,\tau)\,d\tau
=\lambda_+(F)+\frac{1}{2\pi}\int_{\Im\tau=h}\widehat F(\tau)
\Wh_{\infty}(0,\tau)\,d\tau ,
\end{equation}
and likewise
\begin{equation}\label{eq:contourinfminus}
W_{\infty}(F)=\lambda_-(F)+\frac{1}{2\pi}\int_{\Im\tau=-h}\widehat F(\tau)
\Wh_{\infty}(0,\tau)\,d\tau .
\end{equation}
\end{enumerate}
In (d) and (e) the functions $\Wh_p$ and $\Wh_{\infty}(0,\cdot)$ on the lines
$\Im\tau=\pm h$ are the meromorphic continuations of (a) and (b).
\end{theorem}

\begin{proof}
(a)  By \eqref{eq:Wphat} with
$\zeta=\zeta_p(\chi)=e^{-L/2+i(\tau L+\theta_p)}$ one has, for real $\tau$,
\begin{equation}\label{eq:twoexp}
\Wh_p=2L\,\Rel\frac{\zeta}{1-\zeta}
=L\Bigl[\frac{1}{1-e^{-w}}+\frac{1}{1-e^{-\tilde w}}-2\Bigr],\qquad
w:=\tfrac L2-i(\tau L+\theta_p),\quad \tilde w:=\tfrac L2+i(\tau L+\theta_p),
\end{equation}
since $\zeta/(1-\zeta)=1/(1-\zeta)-1$ and $\bar\zeta=e^{-\tilde w}$ for real
$\tau$.  The right-hand side is meromorphic in $\tau\in\C$ and
$2\pi/L$-periodic, since $\tau\mapsto\tau+2\pi/L$ sends $w\mapsto w-2\pi i$ and
$\tilde w\mapsto\tilde w+2\pi i$.  Apply the Mittag--Leffler expansion
\begin{equation}\label{eq:coth}
\frac{1}{1-e^{-z}}=\frac12+\frac12\coth\frac z2
=\frac12+\lim_{K\to\infty}\sum_{|k|\le K}\frac{1}{z-2\pi ik}
\end{equation}
\cite[\S4.36]{DLMF}, valid for $z\notin2\pi i\Z$ with locally uniform
convergence, the terms $\pm k$ combining to $2z/(z^{2}+4\pi^{2}k^{2})$.
Substituting $z=w$ and $z=\tilde w$,
\[
\Wh_p=-L+L\lim_{K\to\infty}\sum_{|k|\le K}\Bigl(\frac1{w-2\pi ik}
+\frac1{\tilde w-2\pi ik}\Bigr).
\]
Now $\tilde w-2\pi ik=iL(\tau-\tau_k^{+})$, because
$iL\bigl(\tau-\tfrac i2-\tfrac{2\pi k-\theta_p}{L}\bigr)
=iL\tau+\tfrac L2+i\theta_p-2\pi ik=\tilde w-2\pi ik$; and
$w-2\pi ik=-iL(\tau-\tau_{-k}^{-})$, because
$-iL\bigl(\tau+\tfrac i2+\tfrac{2\pi k+\theta_p}{L}\bigr)
=-iL\tau+\tfrac L2-i\theta_p-2\pi ik=w-2\pi ik$.  Hence
$L/(\tilde w-2\pi ik)=-i/(\tau-\tau_k^{+})$ and
$L/(w-2\pi ik)=+i/(\tau-\tau_{-k}^{-})$.  Re-indexing $k\mapsto-k$ in the
second family, which does not change the sum over $|k|\le K$, gives
\[
\Wh_p=-L-i\lim_{K\to\infty}\sum_{|k|\le K}
\Bigl[\frac{1}{\tau-\tau_k^{+}}-\frac{1}{\tau-\tau_k^{-}}\Bigr].
\]
Since $\tau_k^{+}-\tau_k^{-}=i$, the $k$-th bracket equals
$i\bigl((\tau-\tau_k^{+})(\tau-\tau_k^{-})\bigr)^{-1}=O(k^{-2})$ locally
uniformly in $\tau$, so the series converges absolutely and \eqref{eq:MLp}
holds.  The points $\tau_k^{\pm}$ are distinct, so the poles are simple with
the stated residues.

(b)  By \eqref{eq:Winfclosed}, $\Wh_{\infty}(0,\tau)=\log\pi-\Rel\,
\psi(\tfrac14+\tfrac{i\tau}2)$ for real $\tau$.  Use the partial fraction
expansion
$\psi(z)=-\gamma+\sum_{n\ge0}\bigl[\tfrac1{n+1}-\tfrac1{n+z}\bigr]$
\cite[\S5.7]{DLMF}, locally uniformly convergent off $-\N_0$.  For real $\tau$,
$\Rel\,\psi(\tfrac14+\tfrac{i\tau}2)=\tfrac12\bigl[\psi(\tfrac14+\tfrac{i\tau}2)
+\psi(\tfrac14-\tfrac{i\tau}2)\bigr]$, and both terms are meromorphic in
$\tau$, so the right-hand side of \eqref{eq:MLinf} is the analytic
continuation.  Since $n+\tfrac14+\tfrac{i\tau}2=\tfrac i2(\tau-\omega_n)$ and
$n+\tfrac14-\tfrac{i\tau}2=-\tfrac i2(\tau+\omega_n)$, one has
$\bigl(n+\tfrac14+\tfrac{i\tau}2\bigr)^{-1}=-2i(\tau-\omega_n)^{-1}$ and
$\bigl(n+\tfrac14-\tfrac{i\tau}2\bigr)^{-1}=2i(\tau+\omega_n)^{-1}$, whence
\[
\tfrac12\bigl[\psi(\tfrac14+\tfrac{i\tau}2)+\psi(\tfrac14-\tfrac{i\tau}2)\bigr]
=-\gamma+\sum_{n\ge0}\Bigl[\frac1{n+1}+\frac{i}{\tau-\omega_n}
-\frac{i}{\tau+\omega_n}\Bigr],
\]
and subtracting from $\log\pi$ gives \eqref{eq:MLinf}.  The bracketed terms are
$O(n^{-2})$ for fixed $\tau$, so the series converges absolutely and locally
uniformly off the poles.

(c)  $\tau_0^{\pm}=\pm i/2-\theta_p/L$ equals $\pm i/2$ if and only if
$\theta_p\equiv0\pmod{2\pi}$, and on $\mathfrak c_0$ every $\kappa_q$ is
trivial, so $\theta_p=0$; and $\omega_0=i/2$.  Summing the residues of the
$|S_f|$ finite terms and the archimedean one gives \eqref{eq:residues}.

(d)  Since $F$ is constant on the compact factor of $C_S$, $\widehat F$
vanishes off $\mathfrak c_0$, and \eqref{eq:plancherelv} gives
$W_p(F)=\tfrac1{2\pi}\int_{\R}\widehat F\Wh_p\,d\tau$.  The function
$\widehat F$ is entire and satisfies \eqref{eq:PW}.
Fix $h>\tfrac12$.  In the closed strip $0\le\Im\tau\le h$ the poles of
$\Wh_p$ are the points $i/2+2\pi k/L$, $k\in\Z$, all at height
$\tfrac12$ by \eqref{eq:MLp} with $\theta_p=0$, the family $\tau_k^{-}$ lying
at height $-\tfrac12$.  Put $R_j:=2\pi(j+\tfrac12)/L$ and integrate
$\widehat F\Wh_p$ over the positively oriented boundary of the rectangle
$[-R_j,R_j]\times[0,h]$.  By $2\pi/L$-periodicity in $\Re\tau$, the modulus of
$\Wh_p$ is bounded on the two vertical sides by a constant independent of $j$,
each side being a translate by an integer multiple of $2\pi/L$ of the fixed
compact segment $\{\Re\tau=\pi/L,\ 0\le\Im\tau\le h\}$, which avoids all poles;
with the rapid decay of $\widehat F$ in $\Re\tau$, uniform for
$0\le\Im\tau\le h$, the two side contributions tend to $0$.  On the two
horizontal sides $\Wh_p$ is bounded: for $\Im\tau=y$ one has
$|e^{-w}|=e^{-L/2-yL}$ and $|e^{-\tilde w}|=e^{-L/2+yL}$ by \eqref{eq:twoexp},
so at $y=0$ both moduli are $e^{-L/2}<1$, and at $y=h>\tfrac12$ the first is
$<e^{-L}<1$ while the second is $e^{L(h-1/2)}>1$; in both cases $|1-e^{-w}|$
and $|1-e^{-\tilde w}|$ are bounded below along the whole line.  Letting
$j\to\infty$, the residue theorem gives
\[
\int_{\Im\tau=0}\widehat F\Wh_p\,d\tau-\int_{\Im\tau=h}\widehat F\Wh_p\,d\tau
=2\pi i\sum_{k\in\Z}\Res_{\tau=i/2+2\pi k/L}\bigl(\widehat F\Wh_p\bigr)
=2\pi i\sum_{k\in\Z}(-i)\widehat F\Bigl(\frac i2+\frac{2\pi k}{L}\Bigr),
\]
the sum converging absolutely by \eqref{eq:PW} at height $\tfrac12$.
Dividing by $2\pi$ gives \eqref{eq:contourp}.  For \eqref{eq:poisson} set
$g(t):=f(t)e^{t/2}\in C_c^{\infty}(\R)$; then
$\widehat F(i/2+2\pi k/L)=\int f(t)e^{t/2}e^{-2\pi ikt/L}dt=\widehat g(2\pi k/L)$,
and the Poisson summation formula \cite[\S7.2]{Hormander} for
$g\in C_c^{\infty}(\R)$ gives $\sum_k\widehat g(2\pi k/L)=L\sum_mg(mL)$.

(e)  By \eqref{eq:plancherelv}, as in (d),
$W_{\infty}(F)=\tfrac1{2\pi}\int_{\R}\widehat F(\tau)\Wh_{\infty}(0,\tau)\,d\tau$.
By (b), in the strip $0\le\Im\tau\le h<\tfrac52$ the only pole of
$\Wh_{\infty}(0,\cdot)$ is $\omega_0=i/2$, of residue $-i$, and there is none
on the line $\Im\tau=h$.  For $\tau=\sigma+iy$ with $0\le y\le h$ the points
$\tfrac14\pm\tfrac{i\tau}2$ have real parts in
$[\tfrac14-\tfrac h2,\tfrac14+\tfrac h2]$ and imaginary parts $\pm\sigma/2$,
so for $|\sigma|\ge1$ they lie in a sector $|\arg z|\le\pi-\delta$, where
$\psi(z)=\log z+O(|z|^{-1})$ \cite[\S5.11]{DLMF}.  Hence
$|\Wh_{\infty}(0,\tau)|\le C\log(2+|\sigma|)$ in the strip for $|\sigma|\ge1$.
With \eqref{eq:PW}, the integrals of $\widehat F\,\Wh_{\infty}(0,\cdot)$ over
the vertical sides of the rectangle $[-\rho,\rho]\times[0,h]$ tend to $0$ as
$\rho\to\infty$ and the two horizontal integrals converge absolutely.  The
residue theorem gives
$\int_{\Im\tau=0}-\int_{\Im\tau=h}=2\pi i\cdot(-i)\widehat F(i/2)$, and
$\widehat F(i/2)=\lambda_+(F)$ by \eqref{eq:lambdaFT}; this is
\eqref{eq:contourinf}.  In the strip $-h\le\Im\tau\le0$ the only pole is
$-\omega_0=-i/2$, of residue $+i$, and the same estimates apply; the residue
theorem for the rectangle $[-\rho,\rho]\times[-h,0]$ gives
$\int_{\Im\tau=-h}-\int_{\Im\tau=0}=2\pi i\cdot(+i)\widehat F(-i/2)$, which is
\eqref{eq:contourinfminus} since $\widehat F(-i/2)=\lambda_-(F)$.
\end{proof}

\begin{remark}\label{rem:MLconseq}
By \eqref{eq:contourinf} and \eqref{eq:contourinfminus}, the functionals
$\lambda_{\pm}(F)=\widehat F(\pm i/2)$ are the residue contributions of the
archimedean symbol at its poles $\pm i/2$.  At a finite place the residue term
in \eqref{eq:contourp} is the sum of $\widehat F$ over
$i/2+(2\pi/\log p)\Z$, which by \eqref{eq:poisson} is $\log p$ times the sum of
$f(t)e^{t/2}$ over $t\in(\log p)\Z$.  By \eqref{eq:residues} the residue of
$\Th(\mathfrak c_0,\cdot)$ at $\pm i/2$ is $\mp i(|S_f|+1)$, each place of $S$
contributing $\mp i$.  On the components with $\varepsilon=1$ the archimedean
symbol is \eqref{eq:Winfodd}, whose continuation has its poles at
$\pm i(2n+\tfrac32)$, $n\ge0$, so it has no pole at $\pm i/2$.
\end{remark}

\section{The numerical theory}\label{sec:numerical}

\subsection{The pairing and its radical}

\begin{definition}\label{def:pairing}
For $F,G\in\DA$ put
\begin{equation}\label{eq:IA}
I_A(F,G):=T\bigl(F\star G^{\dagger}\bigr),
\end{equation}
and let
\begin{equation}\label{eq:radical}
\Nrad:=\bigl\{F\in\DA:\ I_A(F,G)=0\ \text{for every }G\in\DA\bigr\},
\qquad N_A:=\DA/\Nrad .
\end{equation}
\end{definition}

A \emph{component} of $\Ch$ is given by a parity $\varepsilon$ and a system
$(\kappa_q)_{q\in S_f}$; for a component $\mathfrak c$ and a function $\phi$ on
$\Ch$ we write $\phi(\mathfrak c,\tau)$ for its value at the point of
$\mathfrak c$ with real coordinate $\tau$.

\begin{lemma}\label{lem:Nseminorm}
For $G\in\DA$ put
\begin{equation}\label{eq:Nsemi}
N(G):=\int_{\Ch}|\widehat G(\chi)|\bigl(|\Th(\chi)|+|c|\bigr)d\chi .
\end{equation}
Then $N(G)<\infty$ for every $G\in\DA$, and $N$ is bounded on every bounded
subset of $\DA$.  Consequently $N$ is a continuous seminorm on $\DA$, and for
all $F,G\in\DA$
\begin{equation}\label{eq:Nineq}
\int_{\Ch}|\widehat F(\chi)|\,|\widehat G(\chi)|\,\bigl(|\Th(\chi)|+|c|\bigr)d\chi
\le\bigl(\sup_{\chi}|\widehat F(\chi)|\bigr)N(G)
\le\lVert F\rVert_{L^{1}}\,N(G).
\end{equation}
\end{lemma}

\begin{proof}
A bounded set $B\subseteq\DA$ lies in one Fr\'echet step: there are a compact
$K\subseteq C_S$ and an integer $N_0\ge0$ such that every $G\in B$ is supported
in $K$, is invariant under $\Theta_{N_0}=\prod_q(1+q^{N_0}\Z_q)$, and has all step
seminorms bounded by constants depending only on $B$.  Invariance forces
$\widehat G(\varepsilon,\tau;(\kappa_q))=0$ unless every $\kappa_q$ is trivial
on $1+q^{N_0}\Z_q$, and there are only finitely many such systems; with the two
parities this leaves finitely many components, say $d$ of them.  On each of
them, with $\chi_{\mathfrak c}$ the restriction of the characters of
$\mathfrak c$ to the compact part of $C_S$ and
$G_{\mathfrak c}(t):=\int G(e^{t}k)\overline{\chi_{\mathfrak c}(k)}\,
d\mu_{\mathrm{cpt}}(k)$, one has
$\widehat G(\mathfrak c,\tau)=\int_{\R}G_{\mathfrak c}(t)e^{-i\tau t}dt$.  The
functions $G_{\mathfrak c}$ are supported in the fixed compact logarithmic
window of $K$, and $|G_{\mathfrak c}|$, $|G_{\mathfrak c}''|$ are bounded by a
constant depending only on $B$; integrating by parts twice gives
$|\widehat G(\mathfrak c,\tau)|\le C'_B\min(1,\tau^{-2})$, uniformly in
$\mathfrak c$ and in $G\in B$.  By Theorem~\ref{thm:symbol}(a),(b),
\begin{equation}\label{eq:logbound}
|\Th(\mathfrak c,\tau)|\le A_S+\sum_{p\in S_f}\cond(\kappa_p)\log p
+B_S\log(2+|\tau|)=:A_{\mathfrak c}+B_S\log(2+|\tau|),
\end{equation}
where $A_S$ and $B_S$ depend only on $S$: the unramified terms of
\eqref{eq:Wphat} are bounded, $\Wh_{\infty}$ is continuous and satisfies
\eqref{eq:Winfhat} uniformly in the parity, and the ramified terms
$-\cond(\kappa_p)\log p$ of \eqref{eq:Wphat} are constant on each component.  The
constant $A_{\mathfrak c}$ depends on the component through its conductors and
is unbounded over $\Ch$.  Let $A_B$ be the maximum of $A_{\mathfrak c}$ over
the $d$ components met by $B$.  Then
\[
N(G)\le d\,C'_B\int_{\R}\min(1,\tau^{-2})
\bigl(A_B+|c|+B_S\log(2+|\tau|)\bigr)\frac{d\tau}{2\pi}<\infty
\]
for every $G\in B$.  The case $B=\{G\}$ gives finiteness for a single $G$.  A
seminorm on a Fr\'echet space which is bounded on bounded sets is continuous,
and a seminorm on an inductive limit which is continuous on every step is
continuous; so $N$ is continuous on $\DA$.  The first inequality of
\eqref{eq:Nineq} is the definition of $N$, and the second is
$|\widehat F(\chi)|\le\lVert F\rVert_{L^{1}}$.
\end{proof}

\begin{proposition}\label{prop:bilin}
For $F,G,H\in\DA$ one has $(G^{\dagger})^{\wedge}=\overline{\widehat G}$,
\begin{equation}\label{eq:daggerreal}
T(H^{\dagger})=\overline{T(H)},
\end{equation}
\begin{equation}\label{eq:IAfourier}
I_A(F,G)=\int_{\Ch}\widehat F(\chi)\,\overline{\widehat G(\chi)}\,
\bigl(\Th(\chi)+c\bigr)\,d\chi ,
\qquad |I_A(F,G)|\le\lVert F\rVert_{L^{1}}\,N(G).
\end{equation}
The form $I_A$ is linear in the first variable, conjugate-linear in the
second, Hermitian, i.e.\ $\overline{I_A(F,G)}=I_A(G,F)$, and jointly continuous
on $\DA\times\DA$.
\end{proposition}

\begin{proof}
Substituting $v=w^{-1}$ and using the inversion invariance of $\dstar w$,
\[
(G^{\dagger})^{\wedge}(\chi)=\int\overline{G(w^{-1})}\,\overline{\chi(w)}\dstar w
=\overline{\int G(v)\overline{\chi(v)}\dstar v}=\overline{\widehat G(\chi)} .
\]
By \eqref{eq:plancherel} and the reality of $\Th$ and $c$
(Theorem~\ref{thm:symbol}(c)),
$T(H^{\dagger})=\int\overline{\widehat H}(\Th+c)\,d\chi=\overline{T(H)}$.
Since the Fourier transform carries convolution to the pointwise product,
$(F\star G^{\dagger})^{\wedge}=\widehat F\,\overline{\widehat G}$, and
\eqref{eq:plancherel} gives the formula for $I_A$; the bound is
\eqref{eq:Nineq}.  Sesquilinearity is clear from \eqref{eq:IAfourier}, and
Hermitian symmetry follows from it because $\Th+c$ is real.  The seminorms
$\lVert\cdot\rVert_{L^{1}}$ and $N$ are continuous on $\DA$, the first because
it is bounded on bounded sets of each step and the second by
Lemma~\ref{lem:Nseminorm}; so the bound in \eqref{eq:IAfourier} gives joint
continuity.
\end{proof}

\begin{proposition}\label{prop:ideal}
For $F,G,H\in\DA$,
\[
I_A(F\star H,G)=I_A(F,G\star H^{\dagger}),\qquad
I_A(F^{\dagger},G)=\overline{I_A(F,G^{\dagger})} .
\]
Consequently $\Nrad$ is a closed ideal of $\DA$, stable under $\dagger$.
\end{proposition}

\begin{proof}
By commutativity and $(G\star H^{\dagger})^{\dagger}=G^{\dagger}\star H$,
\[
I_A(F\star H,G)=T\bigl(F\star H\star G^{\dagger}\bigr)
=T\bigl(F\star(G\star H^{\dagger})^{\dagger}\bigr)=I_A(F,G\star H^{\dagger}),
\]
and by \eqref{eq:daggerreal},
$I_A(F^{\dagger},G)=T\bigl((F\star G)^{\dagger}\bigr)=\overline{T(F\star G)}
=\overline{I_A(F,G^{\dagger})}$.  If $F\in\Nrad$, the right-hand sides vanish
for all $G,H$, so $F\star H\in\Nrad$ and $F^{\dagger}\in\Nrad$.  Finally
$\Nrad=\bigcap_G\ker I_A(\cdot,G)$ is an intersection of kernels of continuous
functionals.
\end{proof}

\begin{theorem}\label{thm:radical}
$\Nrad=\{0\}$ for every value of the constant $c$.  Hence $N_A=\DA$, and the
form $I_A$ is nondegenerate on $\DA$.
\end{theorem}

\begin{proof}
Let $F\in\Nrad$.  Fix a component $\mathfrak c$, with parity $\varepsilon$ and
system $(\kappa_q)$, and let $g\in C_c^{\infty}(\R^{\times})$ satisfy
$g(-x)=(-1)^{\varepsilon}g(x)$.  In the coordinates of Lemma~\ref{thm:split}
put
\[
G\bigl(x_{\infty},(w_q)\bigr):=g(x_{\infty})\prod_{q\in S_f}\kappa_q(w_q).
\]
Then $G\in\DA$.  By orthogonality of characters on the compact part,
$\widehat G$ vanishes off $\mathfrak c$, and by \eqref{eq:archmeasure}
\[
\widehat G(\mathfrak c,\tau)=\int_{\R^{\times}}g(x)\sgn(x)^{\varepsilon}
|x|^{-i\tau}\frac{dx}{2|x|}=\int_{\R}g(e^{t})e^{-i\tau t}\,dt ,
\]
the Fourier transform of the function $t\mapsto g(e^{t})$, which is an
arbitrary element of $C_c^{\infty}(\R)$.  Thus \eqref{eq:IAfourier} gives
\[
\int_{\R}m(\tau)\,\overline{\widehat h(\tau)}\,d\tau=0
\quad\text{for every }h\in C_c^{\infty}(\R),\qquad
m(\tau):=\widehat F(\mathfrak c,\tau)\bigl(\Th(\mathfrak c,\tau)+c\bigr).
\]
The function $m$ is continuous and, by \eqref{eq:logbound}, is the product of
a Schwartz function and a function of at most logarithmic growth; in
particular it is integrable.  The functions $\overline{\widehat h}$,
$h\in C_c^{\infty}(\R)$, are the Fourier transforms of the elements of
$C_c^{\infty}(\R)$ and are therefore dense in the Schwartz space, hence dense
in $C_0(\R)$ for the uniform norm; so $m=0$.  By Theorem~\ref{thm:symbol}(c)
the zeros of $\Th(\mathfrak c,\cdot)+c$ are isolated, so
$\widehat F(\mathfrak c,\tau)=0$ off a discrete set, and by continuity
$\widehat F$ vanishes on $\mathfrak c$.  As $\mathfrak c$ was arbitrary,
$\widehat F=0$.  Since $\widehat F$ is integrable, Fourier inversion gives
$F=0$.
\end{proof}

The proof uses only that $\Th+c$ is real-analytic and not identically zero on
each component, and of at most logarithmic growth there.

\subsection{The Tate boundary sequence}

\begin{definition}\label{def:prim}
The \emph{primitive} subspace is
$\DA^{\mathrm{prim}}:=\ker\lambda_-\cap\ker\lambda_+=\ker M$.  Put
\[
I_{\mathrm{Tate}}(F,G):=\beta_T(MF,MG)=\lambda_-(F)\overline{\lambda_+(G)}
+\lambda_+(F)\overline{\lambda_-(G)},\qquad
I_{\mathrm{vec}}:=(I_A,I_{\mathrm{Tate}}).
\]
\end{definition}

For $w\in C_S$ and $G\in\DA$ we write $G_w(v):=G(w^{-1}v)$ for the translate.
Then
\begin{equation}\label{eq:translate}
\widehat{G_w}(\chi)=\overline{\chi(w)}\,\widehat G(\chi),\qquad
\lambda_{\mp}(G_w)=|w|_S^{\mp1/2}\lambda_{\mp}(G),\qquad N(G_w)=N(G).
\end{equation}

\begin{lemma}\label{lem:orthprim}
Let $F\in\DA$ satisfy $I_A(F,G)=0$ for every $G\in\DA^{\mathrm{prim}}$.  Then
$F=0$.
\end{lemma}

\begin{proof}
The linear functional $G\mapsto\overline{I_A(F,G)}$ vanishes on
$\ker\lambda_-\cap\ker\lambda_+$, hence is a linear combination of $\lambda_-$
and $\lambda_+$: there are scalars $a,b$ with
$I_A(F,G)=a\,\overline{\lambda_-(G)}+b\,\overline{\lambda_+(G)}$ for every
$G\in\DA$.  Choose $G\in\DA$ with $G\ge0$ and $G\ne0$, so that
$\lambda_{\pm}(G)>0$.  By \eqref{eq:IAfourier} and \eqref{eq:translate},
\[
\bigl|a\,|w|_S^{-1/2}\lambda_-(G)+b\,|w|_S^{1/2}\lambda_+(G)\bigr|
=|I_A(F,G_w)|\le\lVert F\rVert_{L^{1}}N(G)
\qquad\text{for every }w\in C_S .
\]
Letting $|w|_S\to0$ gives $a=0$, and letting $|w|_S\to\infty$ gives $b=0$.
Hence $I_A(F,\cdot)=0$ on $\DA$, and $F=0$ by Theorem~\ref{thm:radical}.
\end{proof}

\begin{theorem}\label{thm:tateseq}
The sequence
\begin{equation}\label{eq:tateseq}
0\longrightarrow \DA^{\mathrm{prim}}\longrightarrow N_A
\stackrel{M}{\longrightarrow}\BT\longrightarrow 0,\qquad N_A=\DA,\quad
\BT=\C^{2}\ \text{with}\
\beta_T=\begin{pmatrix}0&1\\1&0\end{pmatrix},
\end{equation}
is an exact sequence of $\DA$-modules, where $\DA$ acts on $\BT$ through the
continuous algebra homomorphism $M$.  One has $M(F^{\dagger})=M(F)^{\sharp}$,
and translation by $w\in C_S$ acts on $\BT$ by
$\rho_T(w)=(|w|_S^{-1/2},|w|_S^{1/2})$.  The kernel $\DA^{\mathrm{prim}}$ is a
closed ideal of codimension two, stable under $\dagger$ and under
translations, and the restriction of $I_A$ to it is nondegenerate.
\end{theorem}

\begin{proof}
Surjectivity, continuity and multiplicativity of $M$ are
Theorem~\ref{thm:tate}, and $N_A=\DA$ is Theorem~\ref{thm:radical}; exactness
in the middle is the definition of $\DA^{\mathrm{prim}}$.  The kernel of a
continuous algebra homomorphism onto $\C^{2}$ is a closed ideal of codimension
two.  The identity
$\lambda_{\mp}(F^{\dagger})=\overline{\lambda_{\pm}(F)}$ of
Theorem~\ref{thm:tate} is $M(F^{\dagger})=M(F)^{\sharp}$, which also gives the
stability under $\dagger$; the action of translations is
\eqref{eq:translate}.  Nondegeneracy on the kernel is
Lemma~\ref{lem:orthprim} applied to $F\in\DA^{\mathrm{prim}}$.
\end{proof}

\begin{theorem}\label{thm:nosplit}
\begin{enumerate}
\item[(a)] No nonzero finite-dimensional subspace of $\DA$ is stable under
convolution by all elements of $\DA$.  In particular there is no nonzero
$\Phi\in\DA$ with
\begin{equation}\label{eq:eigen}
a\star\Phi=\lambda_{\mp}(a)\,\Phi\qquad\text{for every }a\in\DA ,
\end{equation}
and there is no $\DA$-module map $s:\BT\to N_A$ with $M\circ s=\id$.
\item[(b)] The $I_A$-orthogonal complement of $\DA^{\mathrm{prim}}$ in $\DA$ is
zero, and so is its $I_{\mathrm{vec}}$-orthogonal complement.
\end{enumerate}
\end{theorem}

\begin{proof}
(a)  Let $V\subseteq\DA$ be a nonzero finite-dimensional subspace stable under
convolution.  The operators $\Phi\mapsto a\star\Phi$, $a\in\DA$, commute, so
they have a common eigenvector in $V$: there are $\Phi\in V$, $\Phi\ne0$, and
a function $\mu:\DA\to\C$ with $a\star\Phi=\mu(a)\Phi$ for every $a$.  Taking
Fourier transforms, $\widehat a(\chi)\widehat\Phi(\chi)=
\mu(a)\widehat\Phi(\chi)$ for every $\chi\in\Ch$.  If
$\widehat\Phi(\chi_0)\ne0$ and $\widehat\Phi(\chi_1)\ne0$, then
$\widehat a(\chi_0)=\mu(a)=\widehat a(\chi_1)$ for every $a\in\DA$, that is,
$\int a(w)\bigl(\overline{\chi_0(w)}-\overline{\chi_1(w)}\bigr)\dstar w=0$ for
every $a\in\DA$, which forces $\chi_0=\chi_1$ since both are continuous.  So
$\widehat\Phi$ vanishes off a single point; being continuous in $\tau$ on each
component, it vanishes identically, and $\Phi=0$ by Fourier inversion.  This
contradiction proves the first assertion.  A solution $\Phi\ne0$ of
\eqref{eq:eigen} would span a stable line.  If $s$ existed, then
$\Phi:=s(e_-)$ would satisfy $M(\Phi)=e_-\ne0$ and, since
$M(a)e_-=\lambda_-(a)e_-$ in $\BT$, also
$a\star\Phi=s(M(a)e_-)=\lambda_-(a)\Phi$.

(b)  The first assertion is Lemma~\ref{lem:orthprim}.  For the second,
$I_{\mathrm{Tate}}(F,G)=\beta_T(MF,MG)=0$ whenever
$G\in\DA^{\mathrm{prim}}$, so the two orthogonal complements coincide.
\end{proof}

\begin{remark}\label{rem:weakhyp}
By Theorem~\ref{thm:signature}(a),(b), for $c>-m$, where $m=\sup_{\Ch}\Th$, the Hermitian form $I_A$ is
indefinite.  Any two vectors on which it takes values of opposite sign span a
plane of signature $(1,1)$, so there are $F,G\in\DA$ with
$I_A(F,F)=I_A(G,G)=0$ and $I_A(F,G)=1$.  Such a hyperbolic plane depends on
the choice of $F$ and $G$, and by Theorem~\ref{thm:nosplit}(a) it is not
stable under convolution.  It is therefore not a module complement of
$\DA^{\mathrm{prim}}$ in $\DA$.
\end{remark}

\subsection{The separated completion}

Let $\DA'$ denote the space of continuous conjugate-linear functionals on
$\DA$, with the topology of uniform convergence on bounded subsets of $\DA$.
Since $G\mapsto\overline G$ is a continuous conjugate-linear involution of
$\DA$, the map $\ell\mapsto\ell(\overline{\,\cdot\,})$ is a linear
homeomorphism of $\DA'$ onto the strong dual of $\DA$.

\begin{theorem}\label{thm:completion}
Put
\begin{equation}\label{eq:J}
J:N_A=\DA\longrightarrow\DA'\oplus\BT,\qquad
J(F):=\bigl(I_A(F,\cdot),\ M(F)\bigr).
\end{equation}
\begin{enumerate}
\item[(a)] $\DA'\oplus\BT$ is a complete nuclear space, and $J$ is linear,
continuous and injective.
\item[(b)] $J$ has dense image.  Thus the closure $\NAd:=\overline{J(N_A)}$ is
the whole space:
\begin{equation}\label{eq:NAdall}
\NAd=\DA'\oplus\BT .
\end{equation}
\item[(c)] The pairing
$I_{\mathrm{vec}}\bigl((\ell,m),G\bigr):=\bigl(\ell(G),\beta_T(m,MG)\bigr)$
of $\NAd$ with $\DA$ extends $I_{\mathrm{vec}}$ along $J$ and is separating
in the first variable.
\item[(d)] The formulas
\[
a\cdot(\ell,m):=\bigl(\ell(a^{\dagger}\star\,\cdot\,),\,M(a)m\bigr),\qquad
(\ell,m)^{\dagger}:=\bigl(\ell^{\dagger},m^{\sharp}\bigr),\quad
\ell^{\dagger}(G):=\overline{\ell(G^{\dagger})},
\]
define a $\DA$-module structure on $\NAd$ in which every $a\in\DA$ acts by a
continuous operator, and a continuous conjugate-linear involution, such that
$J(a\star F)=a\cdot J(F)$, $J(F^{\dagger})=J(F)^{\dagger}$ and
$(a\cdot x)^{\dagger}=a^{\dagger}\cdot x^{\dagger}$.
\end{enumerate}
\end{theorem}

\begin{proof}
(a)  The space $\DA$ is a countable strict inductive limit of nuclear
Fr\'echet spaces.  It is bornological, so its strong dual is complete
\cite[Ch.~IV, \S6]{Schaefer}; and the strong dual is nuclear, because every
bounded subset of $\DA$ lies in a step, so that the strong dual is the
projective limit of the strong duals of the steps
\cite[Ch.~IV, \S4]{Schaefer}, which are nuclear
\cite[Ch.~IV, \S9]{Schaefer}, \cite[Ch.~50]{Treves}, and projective limits of
nuclear spaces are nuclear \cite[Ch.~III, \S7]{Schaefer}.  Hence
$\DA'\oplus\BT$ is complete and nuclear.  If $B\subseteq\DA$ is bounded then
$\sup_{G\in B}|I_A(F,G)|\le\lVert F\rVert_{L^{1}}\sup_BN$ by
\eqref{eq:IAfourier}, and $\sup_BN<\infty$ by Lemma~\ref{lem:Nseminorm}; so
$F\mapsto I_A(F,\cdot)$ is continuous from $\DA$ to $\DA'$, and $M$ is
continuous.  If $J(F)=0$ then $F\in\Nrad=\{0\}$.

(b)  The space $\DA$ is complete, barrelled and nuclear, hence a Montel space,
hence reflexive \cite[Ch.~III, \S7 and Ch.~IV, \S5]{Schaefer}.  Therefore
every continuous linear functional on $\DA'\oplus\BT$ has the form
$(\ell,m)\mapsto\ell(G)+\phi(m)$ with $G\in\DA$ and $\phi$ a linear
functional on $\BT$.  Suppose that it vanishes on $J(\DA)$:
$I_A(F,G)+\phi(MF)=0$ for every $F\in\DA$.  For $F\in\DA^{\mathrm{prim}}$ this
gives $I_A(G,F)=\overline{I_A(F,G)}=0$, so $G=0$ by
Lemma~\ref{lem:orthprim}.  Then $\phi(MF)=0$ for every $F$, and $\phi=0$
because $M$ is surjective.  By the Hahn--Banach theorem $J(\DA)$ is dense.

(c)  The extension property is the definition of $J$.  If
$I_{\mathrm{vec}}((\ell,m),G)=0$ for all $G$ then $\ell=0$, and $m=0$ because
$M$ is surjective and $\beta_T$ nondegenerate.

(d)  For fixed $a$ the map $G\mapsto a^{\dagger}\star G$ is linear and
continuous on $\DA$, so its transpose $\ell\mapsto\ell(a^{\dagger}\star\cdot)$
is continuous on $\DA'$; and
$(a\star b)^{\dagger}\star G=b^{\dagger}\star(a^{\dagger}\star G)$ shows that
this is an action.  The map $\ell\mapsto\ell^{\dagger}$ is well defined,
because $G\mapsto\overline{\ell(G^{\dagger})}$ is conjugate-linear and
continuous, and it is a continuous conjugate-linear involution, because
$\dagger$ maps bounded sets of $\DA$ to bounded sets.  By
Proposition~\ref{prop:ideal},
$I_A(a\star F,G)=I_A(F,a^{\dagger}\star G)$ and
$I_A(F^{\dagger},G)=\overline{I_A(F,G^{\dagger})}$; together with
$M(a\star F)=M(a)M(F)$ and $M(F^{\dagger})=M(F)^{\sharp}$ these give
$J(a\star F)=a\cdot J(F)$ and $J(F^{\dagger})=J(F)^{\dagger}$.  The identity
$(a\cdot x)^{\dagger}=a^{\dagger}\cdot x^{\dagger}$ holds on $J(\DA)$, since
$(a\star F)^{\dagger}=a^{\dagger}\star F^{\dagger}$, and extends to $\NAd$ by
continuity and (b).
\end{proof}

\subsection{The rulings}

\begin{theorem}\label{thm:rulings}
Fix a real $\vartheta\in C_c^{\infty}((0,\infty))$ with $\vartheta\ge0$,
$\int_{\R}\vartheta(s)\,ds=1$ and $\supp\vartheta\subseteq[a,b]$,
$0<a<b<\infty$.  For $n\ge1$ define $F_n,G_n\in\DA$ by
\begin{equation}\label{eq:spread}
F_n(w):=n^{-1}|w|_S^{-1/2}\vartheta\Bigl(\frac{\log|w|_S}{n}\Bigr),\qquad
G_n(w):=n^{-1}|w|_S^{+1/2}\vartheta\Bigl(-\frac{\log|w|_S}{n}\Bigr),
\end{equation}
so that $G_n=F_n^{\dagger}$.  Then:
\begin{enumerate}
\item[(a)] $\lambda_+(F_n)=1$ and $0<\lambda_-(F_n)\le e^{-na}$;
$\lambda_-(G_n)=1$ and $0<\lambda_+(G_n)\le e^{-na}$.  Hence $M(F_n)\to e_+$
and $M(G_n)\to e_-$ in $\BT$.
\item[(b)] $\widehat F_n$ and $\widehat G_n=\overline{\widehat F_n}$ are
supported on the principal component $\mathfrak c_0$, and
\begin{equation}\label{eq:supFhat}
\sup_{\chi}|\widehat F_n(\chi)|\le\lVert F_n\rVert_{L^{1}}\le e^{-na/2},\qquad
N(F_n)=N(G_n)\le C_1\,n^{-2}e^{-na/2},
\end{equation}
with a constant $C_1$ depending only on $S$, $c$ and $\vartheta$.
\item[(c)] $I_A(F_n,\cdot)\to0$ and $I_A(G_n,\cdot)\to0$ in $\DA'$, that is,
uniformly on every bounded subset of $\DA$.
\item[(d)] $J(F_n)\to(0,e_+)$ and $J(G_n)\to(0,e_-)$ in $\DA'\oplus\BT$.
\item[(e)] The limits $e_{\pm}^{\mathrm{dist}}:=(0,e_{\pm})$ satisfy
$a\cdot e_{\pm}^{\mathrm{dist}}=\lambda_{\pm}(a)e_{\pm}^{\mathrm{dist}}$ for
every $a\in\DA$, and $(e_{\pm}^{\mathrm{dist}})^{\dagger}=
e_{\mp}^{\mathrm{dist}}$.
\end{enumerate}
\end{theorem}

\begin{proof}
Write $t=\log|w|_S$.  By Lemma~\ref{thm:split} and the normalization
\eqref{eq:archmeasure}, $\dstar w=dt\otimes\mu_{\mathrm{cpt}}$ with
$\mu_{\mathrm{cpt}}$ a probability measure on the compact part; so for a
function $\Phi$ of $|w|_S$ alone,
$\int_{C_S}\Phi(|w|_S)\dstar w=\int_{\R}\Phi(e^{t})dt$.  The functions $F_n$
are smooth, depend on $|w|_S$ alone and are supported in $na\le t\le nb$, so
they lie in $\DA$; and $F_n^{\dagger}(w)=F_n(w^{-1})=G_n(w)$ since $\vartheta$
is real.

(a)  Substituting $t=ns$,
\[
\lambda_-(F_n)=\int_{\R}n^{-1}\vartheta(t/n)e^{-t}dt=\int\vartheta(s)e^{-ns}ds,
\qquad
\lambda_+(F_n)=\int_{\R}n^{-1}\vartheta(t/n)dt=1 .
\]
The first is positive because $\vartheta\ge0$ and $\vartheta\ne0$, and is at
most $e^{-na}$ because $s\ge a$ on $\supp\vartheta$.  The statements for
$G_n=F_n^{\dagger}$ follow from
$\lambda_{\mp}(F^{\dagger})=\overline{\lambda_{\pm}(F)}$.

(b)  $F_n$ is constant on the compact part of $C_S$, so by orthogonality
$\widehat F_n$ vanishes off $\mathfrak c_0$, where
\[
\widehat F_n(\tau)=\int_{\R}n^{-1}e^{-t/2}\vartheta(t/n)e^{-i\tau t}dt=
\int\vartheta(s)e^{-ns(\frac12+i\tau)}ds .
\]
Also $\lVert F_n\rVert_{L^{1}}=\int\vartheta(s)e^{-ns/2}ds\le e^{-na/2}$.
Integrating by parts twice in $s$,
\[
\widehat F_n(\tau)=\frac{1}{n^{2}(\frac12+i\tau)^{2}}
\int\vartheta''(s)e^{-ns(\frac12+i\tau)}ds,\qquad
|\widehat F_n(\tau)|\le
\frac{\lVert\vartheta''\rVert_{L^{1}}\,e^{-na/2}}{n^{2}(\frac14+\tau^{2})} .
\]
With \eqref{eq:logbound} on $\mathfrak c_0$ this gives the bound for $N(F_n)$
with
\[
C_1:=\lVert\vartheta''\rVert_{L^{1}}\int_{\R}
\frac{A_{\mathfrak c_0}+|c|+B_S\log(2+|\tau|)}{\frac14+\tau^{2}}\,
\frac{d\tau}{2\pi}<\infty .
\]
Finally $\widehat G_n=\overline{\widehat F_n}$ by
Proposition~\ref{prop:bilin}, so $N(G_n)=N(F_n)$.

(c)  By \eqref{eq:IAfourier}, for every bounded $B\subseteq\DA$,
\[
\sup_{G\in B}|I_A(F_n,G)|\le\lVert F_n\rVert_{L^{1}}
\sup_{G\in B}N(G)\le e^{-na/2}\sup_BN\longrightarrow0 ,
\]
and $\sup_BN<\infty$ by Lemma~\ref{lem:Nseminorm}.  The same holds for $G_n$,
since $\lVert G_n\rVert_{L^{1}}=\lVert F_n\rVert_{L^{1}}$.

(d)  $J(F_n)=(I_A(F_n,\cdot),M(F_n))$; the first component tends to $0$ in
$\DA'$ by (c) and the second to $e_+$ by (a).  Likewise for $G_n$.

(e)  By Theorem~\ref{thm:completion}(d),
$a\cdot(0,e_{\pm})=(0,M(a)e_{\pm})=\lambda_{\pm}(a)(0,e_{\pm})$, the product
in $\BT$ being componentwise, and $(0,e_{\pm})^{\dagger}=(0,e_{\pm}^{\sharp})
=(0,e_{\mp})$.
\end{proof}

\begin{theorem}[Splitting of the completion]\label{thm:splitting}
Let $\bar M:\NAd\to\BT$ be the second projection, $\iota(m):=(0,m)$, and
$N_A^{\mathrm{dist,prim}}:=\ker\bar M=\DA'\oplus0$.  Then $\iota$ is a
continuous injective $\DA$-module map with $\bar M\circ\iota=\id_{\BT}$,
$\bar M\circ J=M$, and
\begin{equation}\label{eq:split}
\NAd=N_A^{\mathrm{dist,prim}}\ \oplus\ \iota(\BT)
\end{equation}
is a decomposition into closed $\DA$-submodules, each stable under $\dagger$,
which acts on the second summand by $\sharp$.  The closure of
$J(\DA^{\mathrm{prim}})$ is $N_A^{\mathrm{dist,prim}}$.
\end{theorem}

\begin{proof}
By \eqref{eq:NAdall} the decomposition \eqref{eq:split} is that of the direct
sum $\DA'\oplus\BT$, and the formulas of Theorem~\ref{thm:completion}(d) act
on the two coordinates separately; this gives all assertions except the last.
For the last, $J(\DA^{\mathrm{prim}})\subseteq\DA'\oplus0$, and a continuous
linear functional $\ell\mapsto\ell(G)$ on $\DA'$ vanishing on
$J(\DA^{\mathrm{prim}})$ has $I_A(G,F)=0$ for all $F\in\DA^{\mathrm{prim}}$,
so $G=0$ by Lemma~\ref{lem:orthprim}.
\end{proof}

Thus the sequence \eqref{eq:tateseq}, which does not split over $\DA$
(Theorem~\ref{thm:nosplit}(a)), splits after completion: the vectors
$e_{\pm}^{\mathrm{dist}}=\iota(e_{\pm})$ are eigenvectors for the two Tate
characters, and they are limits of the sequences $J(F_n)$, $J(G_n)$ of
Theorem~\ref{thm:rulings}.  We call $\iota(\BT)$ the \emph{ruling plane}.

\begin{theorem}\label{thm:rulnull}
Put
\begin{equation}\label{eq:Rboundary}
R_{\partial}:=\bigl\{x\in\NAd:\ \text{the first component of }
I_{\mathrm{vec}}(x,G)\text{ vanishes for every }G\in\DA\bigr\}.
\end{equation}
Then:
\begin{enumerate}
\item[(a)] $R_{\partial}=\iota(\BT)$, and on it the second component of
$I_{\mathrm{vec}}$ is $\beta_T(m,MG)$, where $\beta_T$ has matrix
$\left(\begin{smallmatrix}0&1\\1&0\end{smallmatrix}\right)$ in the basis
$e_{\mp}$;
\item[(b)] $|I_A(F_n,G_m)|\le C_1\,e^{-na/2}m^{-2}e^{-ma/2}$ for all
$n,m\ge1$, so $I_A(F_n,G_m)\to0$ as $n\to\infty$ uniformly in $m$ and as
$m\to\infty$ uniformly in $n$, whereas
$\beta_T(MF_n,MG_m)\to\beta_T(e_+,e_-)=1$ as $n,m\to\infty$;
\item[(c)] $R_{\partial}\cap J(N_A)=0$.
\end{enumerate}
\end{theorem}

\begin{proof}
(a)  The first component of $I_{\mathrm{vec}}((\ell,m),G)$ is $\ell(G)$, which
vanishes for all $G$ if and only if $\ell=0$; the form is
Theorem~\ref{thm:tate}.

(b)  By \eqref{eq:IAfourier} and \eqref{eq:supFhat},
$|I_A(F_n,G_m)|\le\lVert F_n\rVert_{L^{1}}N(G_m)\le
C_1e^{-na/2}m^{-2}e^{-ma/2}$.  By Theorem~\ref{thm:rulings}(a),
$\beta_T(MF_n,MG_m)=\lambda_-(F_n)\overline{\lambda_+(G_m)}+
\lambda_+(F_n)\overline{\lambda_-(G_m)}=\lambda_-(F_n)\lambda_+(G_m)+1\to1$.

(c)  If $J(F)\in R_{\partial}$ then $I_A(F,G)=0$ for every $G$, so $F=0$ by
Theorem~\ref{thm:radical}.
\end{proof}

\begin{definition}\label{def:admissible}
A net $(F_{\alpha})$ in $\DA$ is an \emph{admissible approximation of $e_+$} if
\begin{equation}\label{eq:admissible}
\text{(i) } M(F_{\alpha})\to e_+\ \text{in }\BT,\qquad
\text{(ii) } \lVert F_{\alpha}\rVert_{L^1(\dstar w)}\to0 ;
\end{equation}
similarly for $e_-$.  A net $(G_{\beta})$ in $\DA$ is \emph{$N$-bounded} if
$\sup_{\beta}N(G_{\beta})<\infty$, with $N$ as in \eqref{eq:Nsemi}.
\end{definition}

For $F\ge0$ one has
$\lVert F\rVert_{L^1}\le\tfrac12(\lambda_+(F)+\lambda_-(F))$, because
$|w|_S^{1/2}+|w|_S^{-1/2}\ge2$; so condition (i) alone bounds
$\lVert F_{\alpha}\rVert_{L^{1}}$ only by a quantity tending to $\tfrac12$.
Under (i) and (ii), for every $R>0$ the part of $\lambda_+(F_{\alpha})$ coming
from $\{|w|_S\le R\}$ is at most $R^{1/2}\lVert F_{\alpha}\rVert_{L^{1}}\to0$:
the weighted mass of $F_{\alpha}$ moves to $|w|_S\to\infty$.

\begin{theorem}\label{thm:rulnullclass}
\begin{enumerate}
\item[(a)] If $(F_{\alpha})$ is an admissible approximation of $e_{\pm}$ then
$J(F_{\alpha})\to \iota(e_{\pm})$ in $\DA'\oplus\BT$; in particular
$I_A(F_{\alpha},\cdot)\to0$ uniformly on bounded subsets of $\DA$.
\item[(b)] If moreover $(G_{\beta})$ is an $N$-bounded net, then
$|I_A(F_{\alpha},G_{\beta})|\le\lVert F_{\alpha}\rVert_{L^1}\sup_{\beta}
N(G_{\beta})\to0$ uniformly in $\beta$.  If $(G_{\beta})$ is an $N$-bounded
admissible approximation of $e_{\mp}$, then
\begin{equation}\label{eq:doublelimit}
\lim_{\alpha}\lim_{\beta}I_A(F_{\alpha},G_{\beta})
=\lim_{\beta}\lim_{\alpha}I_A(F_{\alpha},G_{\beta})=0,\qquad\text{while}\quad
\beta_T(MF_{\alpha},MG_{\beta})\to\beta_T(e_{\pm},e_{\mp})=1 .
\end{equation}
\item[(c)] The sequences \eqref{eq:spread} are admissible and $N$-bounded:
$(F_n)$ for $e_+$ and $(G_m)$ for $e_-$.
\end{enumerate}
\end{theorem}

\begin{proof}
(a)  By \eqref{eq:IAfourier}, for $G$ in a bounded set $B$,
$|I_A(F_{\alpha},G)|\le\lVert F_{\alpha}\rVert_{L^1}\sup_BN\to0$, with
$\sup_BN<\infty$ by Lemma~\ref{lem:Nseminorm}; the second component is
\eqref{eq:admissible}(i).

(b)  The inequality is \eqref{eq:IAfourier}, and it gives
$\lim_{\alpha}I_A(F_{\alpha},G_{\beta})=0$ for every $\beta$.  For fixed
$\alpha$, Hermitian symmetry and \eqref{eq:IAfourier} give
$|I_A(F_{\alpha},G_{\beta})|=|I_A(G_{\beta},F_{\alpha})|\le
\lVert G_{\beta}\rVert_{L^{1}}N(F_{\alpha})\to0$, because $(G_{\beta})$ is
admissible.  So both iterated limits exist and vanish.  The last assertion
holds because $\beta_T$ is continuous.

(c)  This is Theorem~\ref{thm:rulings}(a),(b).
\end{proof}

\begin{remark}\label{rem:noextension}
Theorem~\ref{thm:rulnullclass} does not extend $I_A$ continuously to
$\NAd\times\NAd$.  The space $\NAd$ is a $\DA$-module paired with $\DA$ by
$I_{\mathrm{vec}}$; the product of two of its elements is not defined, and no
form is defined on $\NAd\times\NAd$.  In particular the first summand of
\eqref{eq:split} carries only its pairing with $\DA$.
\end{remark}

\begin{remark}\label{rem:nocontradiction}
Theorem~\ref{thm:nosplit}(a) excludes eigenvectors of the Tate characters in
$\DA$.  The eigenvectors $e_{\pm}^{\mathrm{dist}}$ of
Theorem~\ref{thm:rulings}(e) lie in $\NAd$ and, by
Theorem~\ref{thm:rulnull}(c), outside $J(N_A)$.  Thus $\BT$ is a quotient of
$\DA$ by Theorem~\ref{thm:tateseq} and a direct summand of the completion
$\NAd$ by Theorem~\ref{thm:splitting}.
\end{remark}

\section{Presentations and the determinant}\label{sec:presented}

The trace $T$ is defined on $\DA$.  In this section we extend it additively
to finite graded families of elements of $\DA$, which record the smeared
objects $E_F$ of Definition~\ref{def:smear} together with a degree.

\begin{definition}\label{def:Corrtr}
A \emph{presentation} is a family
$\mathfrak f=(F_{n,i})_{n\in\Z,\,1\le i\le k_n}$ of elements of $\DA$, with
integers $k_n\ge0$ such that $k_n=0$ for all but finitely many $n$.  We write
$\mathfrak f_n:=(F_{n,i})_{1\le i\le k_n}$ for its part of degree $n$, and
$\Corrtr$ for the set of presentations.  Put
\begin{equation}\label{eq:s}
s(\mathfrak f):=\sum_n(-1)^{n}\sum_{i=1}^{k_n}F_{n,i}\in\DA .
\end{equation}
For presentations $\mathfrak f=(F_{n,i})_{i\le k_n}$ and
$\mathfrak g=(G_{n,j})_{j\le l_n}$ define:
the \emph{sum} $\mathfrak f\oplus\mathfrak g$, whose part of degree $n$ is the
concatenation $\mathfrak f_n\sqcup\mathfrak g_n$; the \emph{shift}
$\mathfrak f[1]_n:=\mathfrak f_{n+1}$; the \emph{cone}
$\Cone(\mathfrak f,\mathfrak g)_n:=\mathfrak g_n\sqcup\mathfrak f_{n+1}$; the
\emph{product}
$(\mathfrak f\star\mathfrak g)_n:=(F_{i,a}\star G_{j,b})_{i+j=n,\,a\le k_i,\,b\le l_j}$;
and the \emph{adjoint}
$\mathfrak f^{\dagger}_n:=(F_{n,i}^{\dagger})_{1\le i\le k_n}$.
\end{definition}

\begin{lemma}\label{lem:slaws}
For presentations $\mathfrak f,\mathfrak g$,
\begin{gather*}
s(\mathfrak f\oplus\mathfrak g)=s(\mathfrak f)+s(\mathfrak g),\qquad
s(\mathfrak f[1])=-s(\mathfrak f),\qquad
s(\Cone(\mathfrak f,\mathfrak g))=s(\mathfrak g)-s(\mathfrak f),\\
s(\mathfrak f\star\mathfrak g)=s(\mathfrak f)\star s(\mathfrak g),\qquad
s(\mathfrak f^{\dagger})=s(\mathfrak f)^{\dagger}.
\end{gather*}
\end{lemma}

\begin{proof}
The first and the last are immediate.  For the shift,
$s(\mathfrak f[1])=\sum_n(-1)^{n}\sum_iF_{n+1,i}=-s(\mathfrak f)$, and the
cone is $\mathfrak g\oplus\mathfrak f[1]$.  For the product,
\[
s(\mathfrak f\star\mathfrak g)=\sum_n(-1)^{n}\sum_{i+j=n}\sum_{a,b}
F_{i,a}\star G_{j,b}=
\Bigl(\sum_i(-1)^{i}\sum_aF_{i,a}\Bigr)\star
\Bigl(\sum_j(-1)^{j}\sum_bG_{j,b}\Bigr),
\]
all sums being finite.
\end{proof}

\begin{proposition}\label{prop:faithful}
Let $\varrho$ be the regular representation of $L^{1}(C_S)$ on $L^{2}(C_S)$,
$\varrho(F)f:=F\star f$.  Then $\varrho$ is injective.  In particular
$F\mapsto\varrho(F)$ is injective on $\DA$.
\end{proposition}

\begin{proof}
Let $(e_k)$ be a sequence in $\DA$ with $e_k\ge0$, $\int e_k\dstar w=1$ and
supports shrinking to the identity (Remark~\ref{rem:approxid}).  Then
$e_k\in L^{1}\cap L^{2}$, and $F\star e_k\to F$ in $L^{1}(C_S)$ for every
$F\in L^{1}(C_S)$, by the continuity of translation in $L^{1}$.  If
$\varrho(F)=0$ then $F\star e_k=0$ for every $k$, so $F=0$.
\end{proof}

\begin{definition}\label{def:target}
Let $\Picb(\C)$ be the group $\C^{\times}$, whose element $s$ we write
$(\C,s)$ and regard as the line $\C$ with base point $s$ and with the metric
\begin{equation}\label{eq:metric}
\lVert1\rVert:=|s|^{-1/2},\qquad\text{equivalently}\qquad
\lVert1\rVert^{2}=|s|^{-1} .
\end{equation}
The product of $\Picb(\C)$ is written $\otimes$, its unit $(\C,1)$ is written
$\mathrm{unit}$, and $\overline{(\C,s)}:=(\C,\bar s)$.  Put
\begin{equation}\label{eq:target}
T_A:=\Picb(\C)\times\Picb(\C),
\end{equation}
with the componentwise operations.  For a presentation $\mathfrak f$ put
\begin{equation}\label{eq:TrA}
T(\mathfrak f):=T\bigl(s(\mathfrak f)\bigr)\in\C,\qquad
\Tr_A(\mathfrak f):=\bigl(\bigl(\C,e^{T(\mathfrak f)}\bigr),\
\mathrm{unit}\bigr)\in T_A,
\end{equation}
and for $F\in\DA$ put $\Det^{\mathrm{rel}}_S(F):=e^{T(F)}$.
\end{definition}

With the convention \eqref{eq:metric}, the element
$\bigl(\C,\Det^{\mathrm{rel}}_S(F\star G^{\dagger})\bigr)=(\C,e^{I_A(F,G)})$
has $\lVert1\rVert^{2}=e^{-\Rel I_A(F,G)}$.  The second factor of $T_A$ is
used in Section~\ref{sec:RR}.

\begin{theorem}[Additivity]\label{thm:additivity}
For presentations $\mathfrak f,\mathfrak g$,
\[
T(\mathfrak f\oplus\mathfrak g)=T(\mathfrak f)+T(\mathfrak g),\qquad
T(\mathfrak f[1])=-T(\mathfrak f),\qquad
T(\Cone(\mathfrak f,\mathfrak g))=T(\mathfrak g)-T(\mathfrak f).
\]
Equivalently
$\Tr_A(\mathfrak f\oplus\mathfrak g)=\Tr_A(\mathfrak f)\otimes\Tr_A(\mathfrak g)$,
$\Tr_A(\mathfrak f[1])=\Tr_A(\mathfrak f)^{-1}$ and
$\Tr_A(\Cone(\mathfrak f,\mathfrak g))=\Tr_A(\mathfrak g)\otimes
\Tr_A(\mathfrak f)^{-1}$.  In particular $\Tr_A$ is unchanged when a
presentation consisting of one element $F$ in each of two adjacent degrees is
added.
\end{theorem}

\begin{proof}
Lemma~\ref{lem:slaws} and the linearity of $T$ give the identities for $T$,
and $e^{x+y}=e^{x}e^{y}$ gives those for $\Tr_A$.  The presentation with $F$
in degrees $n$ and $n+1$ has $s=(-1)^{n}(F-F)=0$.
\end{proof}

\begin{theorem}[Symmetry and reality]\label{thm:cyclic}
$\Tr_A(\mathfrak f\star\mathfrak g)=\Tr_A(\mathfrak g\star\mathfrak f)$ and
$\Tr_A(\mathfrak f^{\dagger})=\overline{\Tr_A(\mathfrak f)}$ for all
presentations $\mathfrak f,\mathfrak g$; that is,
$T(\mathfrak f\star\mathfrak g)=T(\mathfrak g\star\mathfrak f)$ and
$T(\mathfrak f^{\dagger})=\overline{T(\mathfrak f)}$.
\end{theorem}

\begin{proof}
By Lemma~\ref{lem:slaws} and the commutativity of $\DA$,
$s(\mathfrak f\star\mathfrak g)=s(\mathfrak f)\star s(\mathfrak g)=
s(\mathfrak g\star\mathfrak f)$.  Reality is \eqref{eq:daggerreal} applied to
$s(\mathfrak f^{\dagger})=s(\mathfrak f)^{\dagger}$.
\end{proof}

\section{The trace Banach algebra}\label{sec:banach}

In this section we construct a complete algebra norm on $\DA$ for which $T$,
and hence $\exp T$, is continuous, and we study the associated determinant
germ.

\begin{definition}\label{def:trnorm}
For $F\in\DA$ let $N(F)$ be the seminorm \eqref{eq:Nsemi}, and put
\begin{equation}\label{eq:trnorm}
\Lambda(F):=\int_{C_S}|F(w)|\bigl(|w|_S^{1/2}+|w|_S^{-1/2}\bigr)\dstar w,
\qquad
\lVert F\rVert_{\mathrm{tr}}:=\Lambda(F)+N(F).
\end{equation}
Here $N(F)<\infty$ by Lemma~\ref{lem:Nseminorm} and $\Lambda(F)<\infty$ because
$F$ has compact support.  Since $|w|_S^{1/2}+|w|_S^{-1/2}\ge2$,
\begin{equation}\label{eq:trnormineq}
\lVert F\rVert_{L^{1}}\le\tfrac12\Lambda(F)\le\tfrac12\lVert
F\rVert_{\mathrm{tr}},\qquad
|\lambda_{\mp}(F)|\le\Lambda(F)\le\lVert F\rVert_{\mathrm{tr}} .
\end{equation}
Let $\Atr$ be the completion of $\DA$ for $\lVert\cdot\rVert_{\mathrm{tr}}$.
The same formulas define $\Lambda(H)$ and $N(H)$ in $[0,\infty]$ for every
$H\in L^{1}(C_S)$.
\end{definition}

\begin{theorem}\label{thm:banach}
\begin{enumerate}
\item[(a)] For all $F,G\in\DA$,
$\lVert F\star G\rVert_{\mathrm{tr}}\le\lVert F\rVert_{\mathrm{tr}}
\lVert G\rVert_{\mathrm{tr}}$ and $\lVert F^{\dagger}\rVert_{\mathrm{tr}}
=\lVert F\rVert_{\mathrm{tr}}$.  Hence $\Atr$ is a commutative Banach
$\dagger$-algebra.  The natural map $\Atr\to L^{1}(C_S)$ is injective and
identifies $\Atr$ with a $\dagger$-subalgebra of $L^{1}(C_S)$ on which
$\lVert H\rVert_{\mathrm{tr}}=\Lambda(H)+N(H)$.  Moreover
$\widehat H\in L^{1}(\Ch,d\chi)$ for every $H\in\Atr$.
\item[(b)] $\lambda_-$ and $\lambda_+$ extend uniquely to continuous linear
functionals on $\Atr$ of norm $\le1$; they remain multiplicative and satisfy
$\lambda_{\mp}(H^{\dagger})=\overline{\lambda_{\pm}(H)}$, so that
$M=(\lambda_-,\lambda_+):\Atr\to\BT$ is a continuous surjective algebra
homomorphism with $M(H^{\dagger})=M(H)^{\sharp}$, extending that of
Theorem~\ref{thm:tate}.
\item[(c)] $T$ extends uniquely to a continuous linear functional on $\Atr$,
given by $T(H)=\int\widehat H(\Th+c)\,d\chi$, with $|T(H)|\le N(H)\le
\lVert H\rVert_{\mathrm{tr}}$.  Hence $\Det^{\mathrm{rel}}_S=\exp T$ extends to
a continuous function on $\Atr$ with
$\Det^{\mathrm{rel}}_S(H+K)=\Det^{\mathrm{rel}}_S(H)\Det^{\mathrm{rel}}_S(K)$.
\item[(d)] For every $H\in\Atr$ and every $r\ge1$,
\begin{equation}\label{eq:powerest}
\bigl|T(H^{\star r})\bigr|\le\lVert H\rVert_{L^{1}}^{\,r-1}N(H)
\le\lVert H\rVert_{L^{1}}^{\,r-1}\lVert H\rVert_{\mathrm{tr}} .
\end{equation}
\item[(e)] For $F\in\Atr$ and $z\in\C$ with $|z|\,\lVert F\rVert_{L^{1}}<1$
the series
\begin{equation}\label{eq:germ}
\log\Delta_F(z):=\sum_{r\ge1}(-1)^{r+1}\frac{z^{r}T(F^{\star r})}{r}
\end{equation}
converges absolutely and
\begin{equation}\label{eq:germint}
\log\Delta_F(z)=\int_{\Ch}\Log\bigl(1+z\widehat F(\chi)\bigr)
\bigl(\Th(\chi)+c\bigr)d\chi ,
\end{equation}
$\Log$ the principal branch, the integral being absolutely convergent.  The
function $\log\Delta_F$ is holomorphic in $z$ on that disc, and for fixed $z$
it is continuous in $F$ on the open set
$\{F\in\Atr:|z|\,\lVert F\rVert_{L^{1}}<1\}$.  We put
$\Delta_F(z):=\exp\log\Delta_F(z)$.
\item[(f)] Let $F,G\in\Atr$ and
$|z|<\min\bigl(1/\lVert F\rVert_{L^{1}},1/\lVert G\rVert_{L^{1}}\bigr)$, and
put $H:=F+G+zF\star G\in\Atr$.  Then
$1+z\widehat H=(1+z\widehat F)(1+z\widehat G)$ takes its values in
$\C\setminus(-\infty,0]$, the integral below converges absolutely, and
\begin{equation}\label{eq:multiplicative}
\Delta_F(z)\,\Delta_G(z)=\exp\int_{\Ch}\Log\bigl(1+z\widehat H(\chi)\bigr)
\bigl(\Th(\chi)+c\bigr)d\chi .
\end{equation}
If moreover
$(1+|z|\lVert F\rVert_{L^{1}})(1+|z|\lVert G\rVert_{L^{1}})<2$, then
$|z|\,\lVert H\rVert_{L^{1}}<1$ and the right-hand side is $\Delta_H(z)$.
Finally $\Delta_{F^{\dagger}}(\bar z)=\overline{\Delta_F(z)}$.
\item[(g)] Let $p\in S_f$ and let
$\Wh^{u}_p(\chi):=-\cond(\kappa_p)\log p$ be the ramified part of the symbol
\eqref{eq:Wphat} at $p$.  For every $H\in\Atr$,
\[
N^{u}_p(H):=\int_{\Ch}|\widehat H(\chi)|\,\cond(\kappa_p)\log p\,d\chi<\infty,
\qquad
W^{u}_p(H):=\int_{\Ch}\widehat H(\chi)\,\Wh^{u}_p(\chi)\,d\chi ,
\]
and for $H\in\DA$ the number $W^{u}_p(H)$ is the second term of \eqref{eq:Wp}.
For $F\in\Atr$ and $r\ge1$,
\[
\bigl|W^{u}_p(F^{\star r})\bigr|\le\lVert F\rVert_{L^{1}}^{\,r-1}N^{u}_p(F),
\]
and for $|z|\,\lVert F\rVert_{L^{1}}<1$ the series
$\sum_{r\ge1}(-1)^{r+1}z^{r}W^{u}_p(F^{\star r})/r$ converges absolutely, is
holomorphic in $z$, and equals
$\int_{\Ch}\Log(1+z\widehat F)\,\Wh^{u}_p\,d\chi$.
\end{enumerate}
\end{theorem}

\begin{proof}
(a)  For $\sigma=\pm\tfrac12$ write
$\lVert F\rVert_{1,\sigma}=\int|F(w)||w|_S^{\sigma}\dstar w$, so
$\Lambda=\lVert\cdot\rVert_{1,1/2}+\lVert\cdot\rVert_{1,-1/2}$.  Since
$w\mapsto|w|_S^{\sigma}$ is multiplicative,
$|F\star G(w)||w|_S^{\sigma}\le\int|F(u)||u|_S^{\sigma}|G(u^{-1}w)|
|u^{-1}w|_S^{\sigma}\dstar u$,
and Fubini's theorem gives
$\lVert F\star G\rVert_{1,\sigma}\le\lVert F\rVert_{1,\sigma}
\lVert G\rVert_{1,\sigma}$.  Adding the two values of $\sigma$,
$\Lambda(F\star G)\le\Lambda(F)\Lambda(G)$.  Since $(F\star G)^{\wedge}=
\widehat F\widehat G$ and $\lVert\widehat G\rVert_{\infty}\le\lVert
G\rVert_{L^{1}}\le\tfrac12\Lambda(G)$ by \eqref{eq:trnormineq},
\[
N(F\star G)=\int|\widehat F||\widehat G|(|\Th|+|c|)d\chi
\le\tfrac12\Lambda(G)N(F).
\]
Adding, $\lVert F\star G\rVert_{\mathrm{tr}}\le\Lambda(G)
(\Lambda(F)+N(F))\le\lVert F\rVert_{\mathrm{tr}}\lVert G\rVert_{\mathrm{tr}}$.
For the involution, $\dstar w$ is inversion invariant and
$|w^{-1}|_S^{\sigma}=|w|_S^{-\sigma}$, so $\Lambda(F^{\dagger})=\Lambda(F)$;
and $(F^{\dagger})^{\wedge}=\overline{\widehat F}$
(Proposition~\ref{prop:bilin}), so $N(F^{\dagger})=N(F)$.

For the realization, let $(F_j)$ be a
$\lVert\cdot\rVert_{\mathrm{tr}}$-Cauchy sequence in $\DA$.  It is Cauchy in
$L^{1}\bigl((|w|_S^{1/2}+|w|_S^{-1/2})\dstar w\bigr)$, hence converges there
and in $L^{1}(C_S)$ to some $H$, and $\widehat F_j\to\widehat H$ uniformly.
The sequence $(\widehat F_j)$ is Cauchy in
$L^{1}\bigl((|\Th|+|c|)d\chi\bigr)$, hence converges there to some $\phi$; a
subsequence converges almost everywhere for that measure, so
$\phi=\widehat H$ almost everywhere for it.  Hence
$\Lambda(F_j-H)+N(F_j-H)\to0$ and
$\lVert F_j\rVert_{\mathrm{tr}}\to\Lambda(H)+N(H)$.  So the map sending the
class of $(F_j)$ to $H$ is well defined and isometric for $\Lambda+N$, hence
injective.  It is multiplicative and compatible with $\dagger$, because
$F_j\star G_j\to H\star K$ and $F_j^{\dagger}\to H^{\dagger}$ in $L^{1}$ when
$F_j\to H$ and $G_j\to K$ in $L^{1}$.

For the last assertion of (a), let $m_{\infty}:=\sup\Wh_{\infty}<\infty$ and
$\beta_q:=2(\log q)q^{-1/2}(1-q^{-1/2})^{-1}$.  By
Theorem~\ref{thm:symbol}(a),(b), for every $\chi=(\varepsilon,\tau;(\kappa_q))$,
\begin{equation}\label{eq:lowerTh}
-\Th(\chi)\ \ge\ \sum_{q:\,\kappa_q\ne1}\cond(\kappa_q)\log q-\Wh_{\infty}(\chi)
-\sum_{q\in S_f}\beta_q .
\end{equation}
Let $E:=\{\chi:|\Th(\chi)|<1\}$.  On $E$ the sum
$\sum_{\kappa_q\ne1}\cond(\kappa_q)\log q$ is at most
$1+\max(m_{\infty},0)+\sum_q\beta_q$, which leaves finitely many systems
$(\kappa_q)$; and $-\Wh_{\infty}(\chi)\le1+\sum_q\beta_q$, which by
\eqref{eq:Winfhat} bounds $|\tau|$.  So $E$ has finite measure $|E|$, and for
$H\in\Atr$
\[
\int_{\Ch}|\widehat H|\,d\chi\le\int_{\Ch\setminus E}|\widehat H|\,
(|\Th|+|c|)\,d\chi+\lVert\widehat H\rVert_{\infty}|E|
\le N(H)+\lVert H\rVert_{L^{1}}|E|<\infty .
\]

(b)  $|\lambda_{\mp}(F)|\le\lVert F\rVert_{1,\mp1/2}\le\Lambda(F)\le
\lVert F\rVert_{\mathrm{tr}}$, whence the unique continuous extension of norm
$\le1$.  Multiplicativity and the involution rule pass to the completion,
both sides of each identity being $\lVert\cdot\rVert_{\mathrm{tr}}$-continuous
and agreeing on the dense subalgebra $\DA$; surjectivity holds already on
$\DA$.

(c)  The functional $H\mapsto\int\widehat H(\Th+c)\,d\chi$ is defined on
$\Atr$ and bounded by $N(H)$, by the definition \eqref{eq:Nsemi} of $N$, and
it agrees with $T$ on $\DA$ by \eqref{eq:plancherel}.  Additivity of $T$
gives the last identity.

(d)  $(H^{\star r})^{\wedge}=\widehat H^{r}$, so by (c)
$|T(H^{\star r})|\le\int|\widehat H|^{r}(|\Th|+|c|)d\chi\le
\lVert\widehat H\rVert_{\infty}^{r-1}N(H)\le\lVert H\rVert_{L^{1}}^{r-1}N(H)$.

(e)  Put $\rho:=|z|\,\lVert F\rVert_{L^{1}}<1$.  By (d),
\[
\sum_{r\ge1}\frac{|z|^{r}|T(F^{\star r})|}{r}\le|z|\,N(F)\sum_{r\ge1}
\rho^{\,r-1}=\frac{|z|\,N(F)}{1-\rho}<\infty ,
\]
and the same bound shows that the power series converges locally uniformly on
the disc, so its sum is holomorphic.  For \eqref{eq:germint}, insert
$T(F^{\star r})=\int\widehat F^{r}(\Th+c)d\chi$ and interchange sum and
integral.  This is legitimate by dominated convergence:
$|z\widehat F(\chi)|\le\rho$, so
$\sum_r|z\widehat F(\chi)|^{r}/r\le|z||\widehat F(\chi)|/(1-\rho)$, and
$|\widehat F|(|\Th|+|c|)$ is integrable; and
$\sum_r(-1)^{r+1}(z\widehat F)^{r}/r=\Log(1+z\widehat F)$ pointwise.  For the
continuity in $F$, fix $z$ and $F_0$ with $|z|\lVert F_0\rVert_{L^{1}}<1$.  By
\eqref{eq:trnormineq} there are a
$\lVert\cdot\rVert_{\mathrm{tr}}$-neighbourhood $U$ of $F_0$ and constants
$\rho_0<1$, $C$ with $|z|\lVert F\rVert_{L^{1}}\le\rho_0$ and $N(F)\le C$ for
$F\in U$.  Each term $T(F^{\star r})$ is
$\lVert\cdot\rVert_{\mathrm{tr}}$-continuous by (a) and (c), and by (d) the
tail $\sum_{r>R}$ of the series is at most $|z|C\rho_0^{R}/(1-\rho_0)$ on $U$;
so the series converges uniformly on $U$.

(f)  Put $\rho:=|z|\max(\lVert F\rVert_{L^{1}},\lVert G\rVert_{L^{1}})<1$.
For each $\chi$ the numbers $1+z\widehat F(\chi)$ and $1+z\widehat G(\chi)$ lie
in the open disc of radius $1$ centred at $1$, so their arguments lie in
$(-\pi/2,\pi/2)$.  Hence their product is not in $(-\infty,0]$, and the sum
of their principal logarithms, having imaginary part in $(-\pi,\pi)$, is the
principal logarithm of the product:
\[
\Log(1+z\widehat F)+\Log(1+z\widehat G)
=\Log\bigl((1+z\widehat F)(1+z\widehat G)\bigr)=\Log(1+z\widehat H),
\]
since $\widehat H=\widehat F+\widehat G+z\widehat F\widehat G$.  The left-hand
side is bounded in modulus by $|z|(|\widehat F|+|\widehat G|)/(1-\rho)$, which
is integrable against $(|\Th|+|c|)d\chi$.  Integrating against
$(\Th+c)d\chi$, using \eqref{eq:germint} for $F$ and for $G$, and
exponentiating gives \eqref{eq:multiplicative}.  Under the additional
hypothesis,
\[
|z|\,\lVert H\rVert_{L^{1}}\le|z|\lVert F\rVert_{L^{1}}+|z|\lVert
G\rVert_{L^{1}}+|z|^{2}\lVert F\rVert_{L^{1}}\lVert G\rVert_{L^{1}}
=(1+|z|\lVert F\rVert_{L^{1}})(1+|z|\lVert G\rVert_{L^{1}})-1<1 ,
\]
so (e) applies to the element $H$ at the point $z$ and identifies the
right-hand side of \eqref{eq:multiplicative} with $\Delta_H(z)$.  For the
conjugation rule, $\lVert F^{\dagger}\rVert_{L^{1}}=\lVert F\rVert_{L^{1}}$,
$(F^{\dagger})^{\wedge}=\overline{\widehat F}$,
$\Log(\bar\zeta)=\overline{\Log(\zeta)}$ for $\zeta\notin(-\infty,0]$, and
$\Th+c$ is real; so
$\int\Log(1+\bar z\,\overline{\widehat F})(\Th+c)d\chi=
\overline{\int\Log(1+z\widehat F)(\Th+c)d\chi}$.

(g)  By \eqref{eq:lowerTh}, $\cond(\kappa_p)\log p\le|\Th(\chi)|+A'$ with
$A':=\max(m_{\infty},0)+\sum_q\beta_q$, so by (a)
\[
N^{u}_p(H)\le N(H)+A'\int_{\Ch}|\widehat H|\,d\chi<\infty\qquad(H\in\Atr),
\]
and $W^{u}_p(H)$ is defined, with $|W^{u}_p(H)|\le N^{u}_p(H)$.  Let
$H\in\DA$ and $a:=H\circ j_p$.  By Fourier inversion, in which only finitely
many components occur, and by $j_p(u)=(1,u,1,\dots)$ for $u\in\Z_p^{\times}$,
\[
a(u)-a(1)=\int_{\Ch}\widehat H(\chi)\bigl(\kappa_p(u)-1\bigr)d\chi .
\]
For each of the finitely many $\kappa_p$ that occur, $\kappa_p(u)-1$ vanishes
on $1+p^{\cond(\kappa_p)}\Z_p$, so $(\kappa_p(u)-1)|1-u|_p^{-1}$ is bounded and
Fubini's theorem applies; with \eqref{eq:conductor},
\[
(\log p)\int_{\Z_p^{\times}}\frac{a(u)-a(1)}{|1-u|_p}\,du
=-(\log p)\int_{\Ch}\widehat H(\chi)\,\cond(\kappa_p)\,d\chi=W^{u}_p(H).
\]
For $F\in\Atr$, estimating as in (d),
$|W^{u}_p(F^{\star r})|\le\int|\widehat F|^{r}\cond(\kappa_p)\log p\,d\chi\le
\lVert F\rVert_{L^{1}}^{\,r-1}N^{u}_p(F)$.  The assertions on the series
follow as in (e), with $\Wh^{u}_p$ in place of $\Th+c$ and $N^{u}_p$ in place
of $N$.
\end{proof}

\begin{theorem}[Presentations over $\Atr$]\label{thm:Atrcat}
Let $\varrho$ be the regular representation of
Proposition~\ref{prop:faithful}.
\begin{enumerate}
\item[(a)] For $H,K\in\Atr$,
$\lVert\varrho(H)\rVert_{\mathrm{op}}\le\lVert H\rVert_{L^{1}}\le
\tfrac12\lVert H\rVert_{\mathrm{tr}}$,
$\varrho(H\star K)=\varrho(H)\varrho(K)$ and
$\varrho(H^{\dagger})=\varrho(H)^{*}$.
\item[(b)] $\varrho$ is injective on $\Atr$.
\item[(c)] Let $\Corrtr(\Atr)$ be the set of presentations with entries in
$\Atr$, with $s$ and the operations of Definition~\ref{def:Corrtr}, and with
$T(\mathfrak f):=T(s(\mathfrak f))$ and $\Tr_A$ defined by \eqref{eq:TrA}.
Then Lemma~\ref{lem:slaws} and Theorems~\ref{thm:additivity} and
\ref{thm:cyclic} hold for $\Corrtr(\Atr)$.  If
$\mathfrak f^{(k)}=(H^{(k)}_{n,i})$ are presentations over $\Atr$ with the
same integers $k_n$, and $H^{(k)}_{n,i}\to H_{n,i}$ in
$\lVert\cdot\rVert_{\mathrm{tr}}$ for all $n,i$, then
$\Tr_A(\mathfrak f^{(k)})\to\Tr_A(\mathfrak f)$ for
$\mathfrak f=(H_{n,i})$.
\end{enumerate}
\end{theorem}

\begin{proof}
(a)  $\lVert H\star f\rVert_{2}\le\lVert H\rVert_{1}\lVert f\rVert_{2}$ is
Young's inequality on $C_S$; the second inequality is \eqref{eq:trnormineq}.
Multiplicativity is the associativity $(H\star K)\star f=H\star(K\star f)$ for
$H,K\in L^{1}$ and $f\in L^{2}$, which follows from Fubini's theorem.  For the
involution, the substitution $u\mapsto u^{-1}$ gives
$(H^{\dagger}\star g)(v)=\int\overline{H(u)}g(uv)\dstar u$, and therefore
$\langle H\star f,g\rangle=\langle f,H^{\dagger}\star g\rangle$ for
$f,g\in L^{2}$.

(b)  This is Proposition~\ref{prop:faithful} composed with the injection
$\Atr\to L^{1}(C_S)$ of Theorem~\ref{thm:banach}(a).

(c)  The proofs of Lemma~\ref{lem:slaws} and of
Theorems~\ref{thm:additivity} and \ref{thm:cyclic} use only that $\DA$ is a
commutative $\dagger$-algebra, that $T$ is linear, and \eqref{eq:daggerreal}.
By Theorem~\ref{thm:banach}(a),(c), $\Atr$ is a commutative
$\dagger$-algebra and $T$ is linear and continuous on it; and
\eqref{eq:daggerreal} extends to $\Atr$ by continuity, $\dagger$ being
isometric.  The last assertion follows from
$s(\mathfrak f^{(k)})\to s(\mathfrak f)$ in $\Atr$, a finite sum, and the
continuity of $T$.
\end{proof}

\begin{remark}\label{rem:whynorm}
The weight $\Lambda$ cannot be replaced by the $L^{1}$ norm.  For the norm
$\lVert F\rVert_{L^{1}}+N(F)$ the functions $F_n$ of \eqref{eq:spread} tend to
zero, by \eqref{eq:supFhat}, while $\lambda_+(F_n)=1$ for every $n$; so
$\lambda_+$, and hence $M$, would not be continuous.  For the norm
\eqref{eq:trnorm} one has
$\lVert F_n\rVert_{\mathrm{tr}}\ge\Lambda(F_n)\ge\lambda_+(F_n)=1$.
\end{remark}

\section{Change of \texorpdfstring{$S$}{S}: inflation, restriction, and the limit trace}\label{sec:charts}

Everything so far was constructed for a fixed finite set $S$ of places.  In
this section we compare the constructions for $S\subseteq S'\subseteq S''$,
all finite, containing $\infty$ and at least one finite place.  We write
$C_S$, $\DA(S)$, $T_S$, $I_A^{S}$ for the objects of the previous sections when
$S$ must be named.

\subsection{Inflation and restriction}

\begin{definition}\label{def:chartsystem}
For $S\subseteq S'$ let
\begin{equation}\label{eq:piTS}
\pi_{S'S}:C_{S'}\longrightarrow C_S
\end{equation}
be the coordinate projection in the coordinates
$C_{S'}\cong\R^{\times}\times\prod_{p\in S'_f}\Z_p^{\times}$ and
$C_S\cong\R^{\times}\times\prod_{p\in S_f}\Z_p^{\times}$ of
Lemma~\ref{thm:split}.  It is a continuous surjective group homomorphism with
compact kernel $K_{S'S}=\prod_{\ell\in S'_f\setminus S_f}\Z_{\ell}^{\times}$,
whose Haar measure $dk$ we normalize to total mass one, and
$\pi_{S'S}\pi_{S''S'}=\pi_{S''S}$.  Since the modulus is the modulus of the
archimedean coordinate (Lemma~\ref{thm:split}),
\begin{equation}\label{eq:chartsystem}
|\pi_{S'S}(w)|_S=|w|_{S'}\qquad(w\in C_{S'}),
\end{equation}
and the Haar measures of Section~\ref{sec:carrier} satisfy
$\dstar w_{S'}=\dstar w_S\otimes dk$.
\end{definition}

\begin{theorem}\label{thm:chartsystem}
Let $S\subseteq S'\subseteq S''$.
\begin{enumerate}
\item[(a)] (Inflation)  $\Inf_{SS'}:=\pi_{S'S}^{*}:\DA(S)\to\DA(S')$ is an
injective $\dagger$-algebra homomorphism, and
$\Inf_{S'S''}\Inf_{SS'}=\Inf_{SS''}$.
\item[(b)] (Restriction)  The map $R_{S'S}:\DA(S')\to\DA(S)$,
\begin{equation}\label{eq:RTS}
(R_{S'S}G)(v):=\int_{K_{S'S}}G(\tilde v\,k)\,dk\qquad
\bigl(\tilde v\in\pi_{S'S}^{-1}(v)\bigr),
\end{equation}
is well defined, is a $\dagger$-algebra homomorphism, and satisfies
$R_{S'S}\circ\Inf_{SS'}=\id_{\DA(S)}$.
\item[(c)] (Presentations)  Applied entrywise to presentations
(Definition~\ref{def:Corrtr}), $\Inf_{SS'}$ and $R_{S'S}$ commute with
$\oplus$, shift, cone, $\star$ and $\dagger$, and
$s(\Inf_{SS'}\mathfrak f)=\Inf_{SS'}s(\mathfrak f)$,
$s(R_{S'S}\mathfrak g)=R_{S'S}s(\mathfrak g)$.
\item[(d)] $M\circ\Inf_{SS'}=M$.
\end{enumerate}
\end{theorem}

\begin{proof}
(a)  Let $F\in\DA(S)$.  The function $F\circ\pi_{S'S}$ is smooth in the
archimedean coordinate and locally constant in the others, and its support
$\pi_{S'S}^{-1}(\supp F)$ is compact because the kernel is compact; so
$\Inf_{SS'}F\in\DA(S')$.  For $F,G\in\DA(S)$ and $w\in C_{S'}$, writing
$u=(x,k)$ with $x\in C_S$, $k\in K_{S'S}$,
\[
(\Inf F\star\Inf G)(w)=\int_{C_S}\int_{K_{S'S}}F(x)\,G\bigl(x^{-1}\pi_{S'S}(w)\bigr)
\,dk\dstar x=(F\star G)\bigl(\pi_{S'S}(w)\bigr),
\]
since $dk$ has total mass one.  Compatibility with $\dagger$ holds because
$\pi_{S'S}$ is a homomorphism.  Injectivity follows from (b), and the
composition rule from $\pi_{S'S}\pi_{S''S'}=\pi_{S''S}$.

(b)  The integrand in \eqref{eq:RTS} is a function on the coset
$\pi_{S'S}^{-1}(v)=\tilde vK_{S'S}$ and $dk$ is translation invariant, so the
integral does not depend on $\tilde v$.  The function $R_{S'S}G$ is supported in the compact set
$\pi_{S'S}(\supp G)$; it is smooth in the archimedean coordinate, on which
$\pi_{S'S}$ is the identity, because differentiation passes under the integral
over the compact group; and it is invariant under the image in
$\prod_{p\in S_f}\Z_p^{\times}$ of any open subgroup under which $G$ is
invariant.  So $R_{S'S}G\in\DA(S)$.  For $G,H\in\DA(S')$, Fubini's theorem and
$\dstar u=\dstar x\otimes dk'$ give
\[
R_{S'S}(G\star H)(v)=\int_{C_{S'}}G(u)\,(R_{S'S}H)\bigl(\pi_{S'S}(u)^{-1}v\bigr)
\dstar u=\int_{C_S}(R_{S'S}G)(x)\,(R_{S'S}H)(x^{-1}v)\dstar x .
\]
Compatibility with $\dagger$ is inversion invariance of $dk$.  Finally
$\Inf_{SS'}F$ is constant on the cosets of $K_{S'S}$, so its average over a
coset is the value of $F$.

(c)  The operations on presentations are defined entrywise from $+$, $\star$
and $\dagger$ in $\DA$, and $s$ is a signed sum; so (c) follows from (a) and
(b).

(d)  By \eqref{eq:chartsystem} and $\dstar w_{S'}=\dstar w_S\otimes dk$,
\[
\lambda_{\mp}(\Inf_{SS'}F)=\int_{C_{S'}}F(\pi_{S'S}w)\,|\pi_{S'S}w|_{S}^{\mp1/2}
\dstar w=\int_{C_S}F(x)|x|_S^{\mp1/2}\dstar x=\lambda_{\mp}(F).\qedhere
\]
\end{proof}

\subsection{The comparison of traces}

For a rational number $t\ne0$ which is a unit at every $p\in S_f$ we write
$\delta_S(t)\in C_S$ for the element with coordinate $t$ at $\infty$ and at
every $p\in S_f$.  By Lemma~\ref{thm:split}, $|\delta_S(t)|_S=|t|$.

\begin{theorem}\label{thm:chartcomp}
Let $S\subseteq S'$.  For every $F\in\DA(S)$,
\begin{equation}\label{eq:chartcomp}
T_{S'}(\Inf_{SS'}F)=T_S(F)+\sum_{\ell\in S'_f\setminus S_f}(\log\ell)
\sum_{k\ne0}\ell^{-|k|/2}F\bigl(\delta_S(\ell^{-k})\bigr),
\end{equation}
and the inner sums are finite.
\end{theorem}

\begin{proof}
We write $j_{S,v}$ for the maps \eqref{eq:slices} of $S$.

Let $p\in S_f$ and $x=p^{k}u\in\Q_p^{\times}$.  By \eqref{eq:slices} for $S'$,
the element $j_{S',p}(x)$ has coordinate $p^{-k}$ at $\infty$, $u$ at $p$, and
$p^{-k}$ at every other place of $S'$.  Applying $\pi_{S'S}$ deletes the
coordinates at $S'_f\setminus S_f$ and leaves $j_{S,p}(x)$.  So
$(\Inf_{SS'}F)\circ j_{S',p}=F\circ j_{S,p}$ on $\Q_p^{\times}$, and the term
$W_p$ of \eqref{eq:Wp} is the same for $(S',\Inf_{SS'}F)$ and for $(S,F)$.
In the same way $\pi_{S'S}\circ j_{S',\infty}=j_{S,\infty}$, so the terms
$W_{\infty}$ agree.

Let $\ell\in S'_f\setminus S_f$ and $x=\ell^{k}u$.  The element $j_{S',\ell}(x)$
has coordinate $\ell^{-k}$ at $\infty$, $u$ at $\ell$, and $\ell^{-k}$ at every
other place of $S'$, so
\[
a(x):=(\Inf_{SS'}F)\bigl(j_{S',\ell}(x)\bigr)=F\bigl(\delta_S(\ell^{-k})\bigr),
\]
which does not depend on $u$.  For $k=0$ this gives $a(u)=F(1)=a(1)$ for every
$u\in\Z_{\ell}^{\times}$, so the integral over the units in \eqref{eq:Wp}
vanishes at $\ell$.  For $k\ne0$, $\tau_{\ell}(\ell^{k}u)=\ell^{-|k|/2}$ by
Lemma~\ref{lem:tau}(i) and the shell $\ell^{k}\Z_{\ell}^{\times}$ has volume
$\log\ell$; hence
$W_{\ell}(\Inf_{SS'}F)=(\log\ell)\sum_{k\ne0}\ell^{-|k|/2}F(\delta_S(\ell^{-k}))$.
Only finitely many $k$ contribute, since $F$ has compact support and
$|\delta_S(\ell^{-k})|_S=\ell^{-k}$.

Finally $(\Inf_{SS'}F)(1)=F(1)$ and the constant $c$ is the same for $S$ and
$S'$ (Convention~\ref{conv:c}).  Summing over $v\in S'$ gives
\eqref{eq:chartcomp}.
\end{proof}

\begin{proposition}\label{prop:resnotrace}
Let $S\subsetneq S'$.  There is $G\in\DA(S')$ with
$T_{S'}(G)\ne T_S(R_{S'S}G)$.
\end{proposition}

\begin{proof}
Choose $\ell\in S'_f\setminus S_f$ and $F\in\DA(S)$ with $F\ge0$ and
$F(\delta_S(\ell^{-1}))>0$, and put $G:=\Inf_{SS'}F$.  Then $R_{S'S}G=F$, and
by \eqref{eq:chartcomp} the difference $T_{S'}(G)-T_S(F)$ is a sum of
nonnegative terms, one of which is
$(\log\ell)\ell^{-1/2}F(\delta_S(\ell^{-1}))>0$.
\end{proof}

\begin{remark}\label{rem:restrictionstatus}
Thus the trace is not compatible with the restriction maps $R_{S'S}$.  It is
compatible with the inflation maps up to the correction \eqref{eq:chartcomp},
which is a finite sum over the new primes; and the limit trace of the next
theorem is invariant under inflation.
\end{remark}

\subsection{The colimit trace}

For $F\in\DA(S)$, $F\ne0$, put
$\operatorname{rad}(F):=\sup\{|\log|w|_S|:w\in\supp F\}<\infty$, and
$\operatorname{rad}(0):=0$.  By \eqref{eq:chartsystem},
$\operatorname{rad}(\Inf_{SS'}F)=\operatorname{rad}(F)$.

\begin{theorem}\label{thm:colimtrace}
Let $F\in\DA(S)$.
\begin{enumerate}
\item[(a)] For a prime $\ell\notin S_f$ with $\log\ell>\operatorname{rad}(F)$,
the term of \eqref{eq:chartcomp} at $\ell$ vanishes.
\item[(b)] For all finite $S'\supseteq
S\cup\{\ell:\log\ell\le\operatorname{rad}(F)\}$ the value
$T_{S'}(\Inf_{SS'}F)$ is the same.  Writing $T_{\infty}(F)$ for it,
\begin{equation}\label{eq:Tinf}
T_{\infty}(F)=T_S(F)+\sum_{\ell\notin S_f,\ \log\ell\le\operatorname{rad}(F)}
(\log\ell)\sum_{k\ne0}\ell^{-|k|/2}F\bigl(\delta_S(\ell^{-k})\bigr),
\end{equation}
a finite sum.
\item[(c)] $T_{\infty}\circ\Inf_{SS'}=T_{\infty}$ for all $S\subseteq S'$.
Hence on the filtered colimit of $\dagger$-algebras
\begin{equation}\label{eq:colim}
\DA^{\infty}:=\operatorname*{colim}_{S}\DA(S)\qquad
\text{(transition maps }\Inf_{SS'})
\end{equation}
the functional $T_{\infty}$ is well defined and linear, it satisfies
$T_{\infty}(H^{\dagger})=\overline{T_{\infty}(H)}$, and
$I_{\infty}(F,G):=T_{\infty}(F\star G^{\dagger})$ is a Hermitian form on
$\DA^{\infty}$.
\item[(d)] $M$ descends to a surjective $\dagger$-algebra homomorphism
$M:\DA^{\infty}\to\BT$, and the sequence
$0\to\operatorname*{colim}_S\DA(S)^{\mathrm{prim}}\to\DA^{\infty}
\xrightarrow{M}\BT\to0$ is exact.
\item[(e)] If $F,G\in\DA(S)$ and
$\operatorname{rad}(F\star G^{\dagger})<\log\ell$ for every prime
$\ell\notin S_f$, then $I_{\infty}(F,G)=I_A^{S}(F,G)$.
\end{enumerate}
\end{theorem}

\begin{proof}
(a)  For $k\ne0$,
$\bigl|\log|\delta_S(\ell^{-k})|_S\bigr|=|k|\log\ell\ge\log\ell>
\operatorname{rad}(F)$, so $F(\delta_S(\ell^{-k}))=0$.

(b)  In \eqref{eq:chartcomp} the terms with
$\log\ell>\operatorname{rad}(F)$ vanish by (a), and the remaining ones are the
terms of \eqref{eq:Tinf}, whatever $S'$ is.

(c)  Let $S''\supseteq S'$ contain all primes $\ell$ with
$\log\ell\le\operatorname{rad}(F)=\operatorname{rad}(\Inf_{SS'}F)$.  By (b),
applied to $F$ and to $\Inf_{SS'}F$, both $T_{\infty}(F)$ and
$T_{\infty}(\Inf_{SS'}F)$ equal
$T_{S''}(\Inf_{S'S''}\Inf_{SS'}F)=T_{S''}(\Inf_{SS''}F)$.  So $T_{\infty}$ is
defined on $\DA^{\infty}$.  Two elements of $\DA^{\infty}$ are represented in
a common $\DA(S)$, and for $S''$ large $T_{\infty}$ is computed on both, and on
their sum, by $T_{S''}\circ\Inf_{SS''}$; so linearity and
$T_{\infty}(H^{\dagger})=\overline{T_{\infty}(H)}$ follow from the linearity of
$T_{S''}$ and \eqref{eq:daggerreal}, the maps $\Inf$ being $\dagger$-algebra
homomorphisms.  The form $I_{\infty}$ is then Hermitian as in
Proposition~\ref{prop:bilin}.

(d)  By Theorem~\ref{thm:chartsystem}(d), $M$ is compatible with the
transition maps, which therefore carry $\DA(S)^{\mathrm{prim}}$ into
$\DA(S')^{\mathrm{prim}}$.  At each $S$ the sequence
$0\to\DA(S)^{\mathrm{prim}}\to\DA(S)\to\BT\to0$ is exact
(Theorem~\ref{thm:tateseq}), and filtered colimits are exact.

(e)  By (a) applied to $F\star G^{\dagger}$, the sum in \eqref{eq:Tinf} is
empty.
\end{proof}

In terms of \eqref{eq:Tinf}, $T_S$ consists of the terms of $T_{\infty}$ at the
places of $S$.

\section{Effectivity}\label{sec:eff}

\begin{definition}\label{def:eff}
The \emph{effective cone} is
\[
\Eff_A^{\mathrm{tr}}:=\{F\in\DA:\ F(u)\ge0\ \text{for every }u\in C_S\}.
\]
A presentation $\mathfrak f=(F_{n,i})$ is \emph{effective} if every $F_{n,i}$
lies in $\Eff_A^{\mathrm{tr}}$.
\end{definition}

\begin{proposition}\label{prop:effsmear}
$\Eff_A^{\mathrm{tr}}$ is a convex cone in $\DA$, closed in the topology of
$\DA$, stable under $\star$ and $\dagger$, and
$\Eff_A^{\mathrm{tr}}\cap(-\Eff_A^{\mathrm{tr}})=\{0\}$.  Effective
presentations are stable under $\oplus$, shift, cone, $\star$ and $\dagger$.
\end{proposition}

\begin{proof}
Sums and nonnegative multiples of nonnegative functions are nonnegative.
Evaluation at a point is continuous on $\DA$, so the cone is an intersection of
closed half-spaces of the real subspace of real-valued functions, which is
closed.  For $F,G\ge0$ one has
$(F\star G)(w)=\int F(u)G(u^{-1}w)\dstar u\ge0$ and
$F^{\dagger}(u)=\overline{F(u^{-1})}\ge0$.  If $F\ge0$ and $-F\ge0$ then
$F=0$.  The operations on presentations only rearrange entries or form
$\star$-products and $\dagger$ of entries.
\end{proof}

The next statement compares $\Eff_A^{\mathrm{tr}}$ with the cone of elements
acting as positive operators.  Let $\varrho$ be the regular representation of
Proposition~\ref{prop:faithful}, so that $\varrho(F)$ is the convolution
operator $f\mapsto F\star f$ on $L^{2}(C_S,\dstar w)$.

\begin{proposition}\label{thm:cone}
Let $F\in\DA$.
\begin{enumerate}
\item[(a)] $\varrho(F)$ is positive semidefinite if and only if
$\widehat F(\chi)\ge0$ for every $\chi\in\Ch$.  In particular
$\lambda(H\star H^{\dagger})$ is positive semidefinite for every $H\in\DA$.
\item[(b)] There is $F\in\Eff_A^{\mathrm{tr}}$ such that $\varrho(F)$ is
self-adjoint and not positive semidefinite.
\item[(c)] There is $H\in\DA$ with
$H\star H^{\dagger}\notin\Eff_A^{\mathrm{tr}}$.
\end{enumerate}
\end{proposition}

\begin{proof}
(a)  By the Plancherel theorem for the locally compact abelian group $C_S$,
the Fourier transform is a unitary map $L^{2}(C_S)\to L^{2}(\Ch)$ which
carries $\varrho(F)$ to multiplication by the bounded continuous function
$\widehat F$.  A multiplication operator by a continuous function is positive
semidefinite if and only if the function is nonnegative.  The transform of
$H\star H^{\dagger}$ is $|\widehat H|^{2}$.

(b)  Let $F\in\DA$ satisfy $F\ge0$, $F\ne0$, $F^{\dagger}=F$ and $F\equiv0$
near $1$; for instance $F=F_0+F_0^{\dagger}$ with $F_0\ge0$, $F_0\ne0$
supported in $\{|w|_S>2\}$.  Then $\widehat F$ is real, integrable and
continuous, and by Fourier inversion $\int_{\Ch}\widehat F\,d\chi=F(1)=0$.
Since $\widehat F\ne0$, it takes negative values, and (a) applies.

(c)  Let $\varphi\in\DA$ with $\varphi\ge0$, $\varphi\ne0$, supported in
$V:=\{w:|\log|w|_S|<\tfrac14\}$, and put $\psi:=\varphi\star\varphi^{\dagger}$.
Then $\psi$ is supported in $\{|\log|w|_S|<\tfrac12\}$ and
$\psi(1)=\int\varphi^{2}\dstar u>0$.  Fix $w\in C_S$ with $|w|_S=e$, and put
$\varphi_w(v):=\varphi(w^{-1}v)$ and $H:=\varphi-\varphi_w$.  Since
$(\varphi_w)^{\dagger}(v)=\varphi^{\dagger}(wv)$,
\[
H\star H^{\dagger}(v)=2\psi(v)-\psi(w^{-1}v)-\psi(wv),
\]
and at $v=w$ this equals $2\psi(w)-\psi(1)-\psi(w^{2})=-\psi(1)<0$, because
$w$ and $w^{2}$ lie outside the support of $\psi$.
\end{proof}

\begin{theorem}\label{thm:coneproper}
Let $\mathcal E:=\overline{J(\Eff_A^{\mathrm{tr}})}\subseteq\NAd$ be the
closure of the image of the effective cone under the map $J$ of
\eqref{eq:J}.  Then $\mathcal E$ is a closed convex cone and
\begin{equation}\label{eq:properdist}
\mathcal E\cap(-\mathcal E)=\{0\}.
\end{equation}
\end{theorem}

\begin{proof}
$J$ is linear, and the closure of a convex cone is a convex cone.  Let
$x=(\ell,m)\in\mathcal E\cap(-\mathcal E)$ and choose nets $(F_{\alpha})$,
$(G_{\beta})$ in $\Eff_A^{\mathrm{tr}}$ with $J(F_{\alpha})\to x$ and
$J(G_{\beta})\to-x$ in $\DA'\oplus\BT$.  We use nets because the strong dual
$\DA'$ is not metrizable.

\emph{Step 1.}  The projection $\bar M$ is continuous, so
$m=\lim_{\alpha}M(F_{\alpha})$ and $-m=\lim_{\beta}M(G_{\beta})$.  For $F\ge0$
one has $\lambda_{\mp}(F)=\int F|w|_S^{\mp1/2}\dstar w\ge0$, so $M(F)$ lies in
the closed quadrant $\R_{\ge0}^{2}$, and so do both limits.  Hence $m=0$, that
is, $\lambda_-(F_{\alpha})\to0$ and $\lambda_+(F_{\alpha})\to0$.

\emph{Step 2.}  For every $w$, $|w|_S^{1/2}+|w|_S^{-1/2}\ge2$; hence for
$F\ge0$
\begin{equation}\label{eq:amgm}
\lVert F\rVert_{L^{1}}=\int F\dstar w\le\tfrac12\int
F(w)\bigl(|w|_S^{1/2}+|w|_S^{-1/2}\bigr)\dstar w
=\tfrac12\bigl(\lambda_+(F)+\lambda_-(F)\bigr).
\end{equation}
By Step 1, $\lVert F_{\alpha}\rVert_{L^{1}}\to0$, so
$\sup_{\chi}|\widehat F_{\alpha}(\chi)|\to0$.

\emph{Step 3.}  For every $G\in\DA$, Lemma~\ref{lem:Nseminorm} gives
$|I_A(F_{\alpha},G)|\le\bigl(\sup|\widehat F_{\alpha}|\bigr)N(G)\to0$.  Since
$I_A(F_{\alpha},G)\to\ell(G)$, we get $\ell=0$.  Hence $x=0$.
\end{proof}

\begin{theorem}\label{thm:coneRR}
Let $F,G\in\Eff_A^{\mathrm{tr}}$.
\begin{enumerate}
\item[(a)] $M(F)=(\lambda_-(F),\lambda_+(F))\in\R_{\ge0}^{2}$, and both
entries are positive unless $F=0$.
\item[(b)] $M(F+G)=M(F)+M(G)$, $M(F\star G)=M(F)M(G)$ with the componentwise
product, and $M(F^{\dagger})=(\lambda_+(F),\lambda_-(F))$.
\item[(c)] One has
\begin{equation}\label{eq:contractive}
\beta_T(MF,MG)=\lambda_-(F)\lambda_+(G)+\lambda_+(F)\lambda_-(G)\ge0,
\end{equation}
with equality if and only if $F=0$ or $G=0$.
\end{enumerate}
\end{theorem}

\begin{proof}
(a)  $\lambda_{\mp}(F)\ge0$ for $F\ge0$, and it is positive unless the
continuous function $F$ vanishes almost everywhere, hence everywhere.  (b) is
Theorem~\ref{thm:tate}, the entries being real by (a).  (c)  The two products
are nonnegative by (a), and if $F\ne0\ne G$ they are positive.
\end{proof}

Thus $M$ maps the effective cone into the closed positive quadrant of $\BT$,
as the bidegree does for effective correspondences on a curve.  By
Proposition~\ref{thm:cone}, effectivity of $F$ and nonnegativity of
$\widehat F$ are independent conditions.

\section{The Riemann--Roch pairing}\label{sec:RR}

In this section $\Picb(\C)$ denotes the group $\C^{\times}$ with the metric
\eqref{eq:metric}, $\lVert1\rVert_s=|s|^{-1/2}$, and $\sharp$ is the
conjugate-linear swap $(z_-,z_+)^{\sharp}=(\overline{z_+},\overline{z_-})$ of
$\BT$.

\subsection{The pairing}

\begin{definition}\label{def:lper}
For $b,b'\in\BT$ put
\begin{equation}\label{eq:lper}
\lambda_{\mathrm{per}}(b,b'):=\exp\beta_T(b,b')\in\Picb(\C),\qquad\text{so}\quad
\lVert1\rVert^{2}_{\lambda_{\mathrm{per}}(b,b')}=
\exp\bigl(-\Rel\,\beta_T(b,b')\bigr).
\end{equation}
We call $\lambda_{\mathrm{per}}$ the \emph{ruling determinant}.  For real
$b=(d_1,d_2)$ one has
$\lVert1\rVert_{\lambda_{\mathrm{per}}(b,b)}=\exp(-d_1d_2)$.
\end{definition}

\begin{lemma}\label{lem:lperbiext}
For $b,b',b''\in\BT$,
\[
\lambda_{\mathrm{per}}(b+b'',b')=\lambda_{\mathrm{per}}(b,b')
\lambda_{\mathrm{per}}(b'',b'),\qquad
\lambda_{\mathrm{per}}(b,b'+b'')=\lambda_{\mathrm{per}}(b,b')
\lambda_{\mathrm{per}}(b,b''),
\]
\[
\lambda_{\mathrm{per}}(b',b)=\overline{\lambda_{\mathrm{per}}(b,b')},\qquad
\lambda_{\mathrm{per}}(b^{\sharp},b'^{\sharp})=
\overline{\lambda_{\mathrm{per}}(b,b')}.
\]
\end{lemma}

\begin{proof}
The form $\beta_T(b,b')=b_-\overline{b'_+}+b_+\overline{b'_-}$ is additive in
each variable and Hermitian, and
$\beta_T(b^{\sharp},b'^{\sharp})=\overline{b_+}\,b'_-+\overline{b_-}\,b'_+
=\overline{\beta_T(b,b')}$.  Apply $\exp$.
\end{proof}

\begin{definition}\label{def:RR}
For $F,G\in\Atr$ put
\begin{equation}\label{eq:LS}
L_S(F,G):=\exp I_A(F,G)=\exp T(F\star G^{\dagger})\in\Picb(\C),
\end{equation}
\begin{equation}\label{eq:lRR}
\lambda_{RR}(F,G):=\bigl(L_S(F,G),\ \lambda_{\mathrm{per}}(MF,MG)\bigr)
\in\Picb(\C)\times\Picb(\C).
\end{equation}
For $x=(\ell,m)\in\NAd$ and $G\in\DA$ put
$\lambda_{RR}(x,G):=\bigl(e^{\ell(G)},\lambda_{\mathrm{per}}(m,MG)\bigr)$, and
for $b,b'\in\BT$ put
$\lambda_{RR}(\iota(b),\iota(b')):=(1,\lambda_{\mathrm{per}}(b,b'))$.  For
$F\in\DA$ one has $\lambda_{RR}(J(F),G)=\lambda_{RR}(F,G)$, since
$J(F)=(I_A(F,\cdot),MF)$.  The componentwise logarithm of $\lambda_{RR}$ is
\begin{equation}\label{eq:logRR}
\bigl(I_A(F,G),\ \beta_T(MF,MG)\bigr)=I_{\mathrm{vec}}(F,G).
\end{equation}
\end{definition}

\begin{theorem}\label{thm:scalartrivial}
Let $(F_{\alpha})$ be an admissible approximation of $e_+$ and $(G_{\beta})$
an admissible $N$-bounded approximation of $e_-$
(Definition~\ref{def:admissible}).  Then
\begin{equation}\label{eq:trivial}
L_S(F_{\alpha},G_{\beta})\longrightarrow1
\quad(\alpha\to\infty,\ \text{uniformly in }\beta),\qquad
\lambda_{\mathrm{per}}(MF_{\alpha},MG_{\beta})\longrightarrow
\lambda_{\mathrm{per}}(e_+,e_-)=e ,
\end{equation}
the second limit being taken in both indices.  The element $e\in\Picb(\C)$
has $\lVert1\rVert^{2}=e^{-1}\ne1$.  Such pairs of approximations exist.
\end{theorem}

\begin{proof}
By Theorem~\ref{thm:rulnullclass}(b),
$|I_A(F_{\alpha},G_{\beta})|\le\lVert F_{\alpha}\rVert_{L^{1}}
\sup_{\beta}N(G_{\beta})\to0$, and
$\beta_T(MF_{\alpha},MG_{\beta})\to\beta_T(e_+,e_-)=1$.  Apply $\exp$.
Existence is Theorem~\ref{thm:rulnullclass}(c).
\end{proof}

So along admissible $N$-bounded approximations the first component of
$\lambda_{RR}$ tends to the value $1$ which Definition~\ref{def:RR} assigns on
$\iota(\BT)\times\iota(\BT)$, while the second component tends to
$\lambda_{\mathrm{per}}(e_+,e_-)=e$.

\begin{remark}\label{rem:whyblind}
By
Lemma~\ref{lem:Nseminorm}, $I_A(\cdot,G)$ is continuous for the norm
$\sup_{\chi}|\widehat F(\chi)|$.  The functionals $\lambda_{\mp}$ are not: for
$F\ge0$, $F\ne0$, the translates $F_w(v)=F(w^{-1}v)$ satisfy
$\sup|\widehat{F_w}|=\sup|\widehat F|$, while
$\lambda_{\mp}(F_w)=|w|_S^{\mp1/2}\lambda_{\mp}(F)$ is unbounded in $w$.
\end{remark}

\begin{remark}\label{rem:notafunctor}
The trace is not a function of the bidegree $M$: by
Theorem~\ref{thm:separate} there are $F,F'\in\DA$ with $M(F)=M(F')$ and
$T(F)\ne T(F')$.
\end{remark}

\begin{theorem}\label{thm:RR}
\begin{enumerate}
\item[(a)] For $F,F',G,G'\in\Atr$,
\[
\lambda_{RR}(F+F',G)=\lambda_{RR}(F,G)\lambda_{RR}(F',G),\qquad
\lambda_{RR}(F,G+G')=\lambda_{RR}(F,G)\lambda_{RR}(F,G'),
\]
\[
\lambda_{RR}(G,F)=\overline{\lambda_{RR}(F,G)},\qquad
\lambda_{RR}(F^{\dagger},G^{\dagger})=\overline{\lambda_{RR}(F,G)},
\]
products and conjugation being componentwise.
\item[(b)] $\lambda_{RR}$ is continuous on $\Atr\times\Atr$ for
$\lVert\cdot\rVert_{\mathrm{tr}}$, and on $\NAd\times\DA$ it is separately
continuous.
\item[(c)] For $S\subseteq S'$ and $F,G\in\DA(S)$,
$\lambda_{\mathrm{per}}(M\Inf_{SS'}F,M\Inf_{SS'}G)=\lambda_{\mathrm{per}}(MF,MG)$,
and $L_{S'}(\Inf_{SS'}F,\Inf_{SS'}G)$ is obtained from $L_S(F,G)$ by applying
\eqref{eq:chartcomp} to $F\star G^{\dagger}$.  On $\DA^{\infty}$ the pair
$\bigl(\exp I_{\infty}(F,G),\lambda_{\mathrm{per}}(MF,MG)\bigr)$ is well
defined.
\item[(d)] In the splitting \eqref{eq:split}, for $F,G\in\DA$,
\begin{equation}\label{eq:weilformula}
\lambda_{RR}(F,G)=\lambda_{RR}\bigl((I_A(F,\cdot),0),G\bigr)\cdot
\lambda_{RR}\bigl(\iota(MF),G\bigr),
\end{equation}
where the first factor is $(L_S(F,G),1)$ and the second is
$(1,\lambda_{\mathrm{per}}(MF,MG))$.  For
$x\in N_A^{\mathrm{dist,prim}}$ and $G\in\DA^{\mathrm{prim}}$ the second
component of $\lambda_{RR}(x,G)$ is $1$.
\item[(e)] For $v\in S$ and $F,G\in\DA$ put
$L^{(v)}(F,G):=\exp W_v(F\star G^{\dagger})$.  Then
\[
L_S(F,G)=e^{c\,(F\star G^{\dagger})(1)}\prod_{v\in S}L^{(v)}(F,G),\qquad
\frac{L^{(v)}(F+G,F+G)}{L^{(v)}(F,F)\,L^{(v)}(G,G)}
=\exp\bigl(2\,\Rel\,W_v(F\star G^{\dagger})\bigr).
\]
\item[(f)] For $H\in\Atr$ and an integer $N>\lVert H\rVert_{L^{1}}$,
\begin{equation}\label{eq:allamp}
\Bigl|N\sum_{r\ge1}(-1)^{r+1}\frac{T\bigl((H/N)^{\star r}\bigr)}{r}-T(H)\Bigr|
\le\frac{\lVert H\rVert_{L^{1}}\lVert H\rVert_{\mathrm{tr}}}
{2N\bigl(1-\lVert H\rVert_{L^{1}}/N\bigr)} .
\end{equation}
\end{enumerate}
\end{theorem}

\begin{proof}
(a)  $I_A$ is additive in each variable and Hermitian
(Proposition~\ref{prop:bilin}, extended to $\Atr$ by continuity,
Theorem~\ref{thm:banach}), and so is $\beta_T(M\cdot,M\cdot)$, $M$ being
linear.  For the involution, $(F^{\dagger})^{\wedge}=\overline{\widehat F}$
and $\Th+c$ is real, so
$I_A(F^{\dagger},G^{\dagger})=\overline{I_A(F,G)}$ by \eqref{eq:IAfourier};
and $\beta_T(MF^{\dagger},MG^{\dagger})=\beta_T((MF)^{\sharp},(MG)^{\sharp})=
\overline{\beta_T(MF,MG)}$ by Lemma~\ref{lem:lperbiext}.

(b)  By Theorem~\ref{thm:banach}(a),(c),
$|I_A(F,G)|\le\lVert F\star G^{\dagger}\rVert_{\mathrm{tr}}\le
\lVert F\rVert_{\mathrm{tr}}\lVert G\rVert_{\mathrm{tr}}$, so the
sesquilinear form $I_A$ is continuous on $\Atr\times\Atr$; by
Theorem~\ref{thm:banach}(b), $M$ is continuous on $\Atr$, and $\beta_T$ is
continuous on $\BT\times\BT$.  On $\NAd\times\DA$ the map
$((\ell,m),G)\mapsto(\ell(G),\beta_T(m,MG))$ is separately continuous by
Theorem~\ref{thm:completion}.  Compose with $\exp$.

(c)  The first statement is Theorem~\ref{thm:chartsystem}(d).  The second
holds because $\Inf_{SS'}$ is a $\dagger$-algebra homomorphism, so
$\Inf_{SS'}(F\star G^{\dagger})=(\Inf_{SS'}F)\star(\Inf_{SS'}G)^{\dagger}$.
The third is Theorem~\ref{thm:colimtrace}(c),(d).

(d)  By Theorem~\ref{thm:splitting},
$J(F)=(I_A(F,\cdot),0)+\iota(MF)$ in $\NAd$.  By Definition~\ref{def:RR},
$\lambda_{RR}((I_A(F,\cdot),0),G)=\bigl(e^{I_A(F,G)},
\lambda_{\mathrm{per}}(0,MG)\bigr)=(L_S(F,G),1)$ and
$\lambda_{RR}(\iota(MF),G)=(e^{0},\lambda_{\mathrm{per}}(MF,MG))$.  Their
product is $\lambda_{RR}(F,G)$.  If $MG=0$ then
$\lambda_{\mathrm{per}}(m,MG)=1$ for every $m$.

(e)  The first identity is \eqref{eq:tracedef} applied to
$F\star G^{\dagger}$.  For the second, by \eqref{eq:plancherelv} and the
reality of $\Wh_v$ (Theorem~\ref{thm:symbol}) one has
$W_v(H^{\dagger})=\overline{W_v(H)}$, and
$G\star F^{\dagger}=(F\star G^{\dagger})^{\dagger}$; hence
\[
W_v\bigl((F+G)\star(F+G)^{\dagger}\bigr)-W_v(F\star F^{\dagger})
-W_v(G\star G^{\dagger})=W_v(F\star G^{\dagger})+W_v(G\star F^{\dagger})
=2\,\Rel\,W_v(F\star G^{\dagger}).
\]

(f)  The term $r=1$ of the sum is $T(H)/N$.  By \eqref{eq:powerest},
$|T((H/N)^{\star r})|\le(\lVert H\rVert_{L^1}/N)^{r-1}
\lVert H\rVert_{\mathrm{tr}}/N$, so with $x:=\lVert H\rVert_{L^{1}}/N<1$,
\[
N\Bigl|\sum_{r\ge2}(-1)^{r+1}\frac{T\bigl((H/N)^{\star r}\bigr)}{r}\Bigr|
\le\lVert H\rVert_{\mathrm{tr}}\sum_{r\ge2}\frac{x^{r-1}}{r}
\le\lVert H\rVert_{\mathrm{tr}}\,\frac{x}{2(1-x)} .\qedhere
\]
\end{proof}

Formula \eqref{eq:weilformula} reads, in logarithms,
$I_{\mathrm{vec}}(F,G)=(I_A(F,G),0)+(0,\beta_T(MF,MG))$.  On $C\times C$, for
a curve $C$, a divisor class $D$ of bidegree $(d_1,d_2)$ is the sum of a
combination of the two rulings and a class $P$ orthogonal to both, and
$D\cdot D'=P\cdot P'+d_1d_2'+d_2d_1'$ \cite{Weil1948}.  The pair
$(I_A,\beta_T\circ(M\times M))$ has the second term of this formula as its
second component.  Its first component is $I_A(F,G)$ for the whole of $F$ and
$G$ and not a pairing of primitive parts, and the two components are not
added.  By Theorem~\ref{thm:signature}, the first component is indefinite for
$c>-m$.

\subsection{Base change along the Tate boundary}

\begin{theorem}\label{thm:basechange}
Let $\DA^{+}=\C1\oplus\DA$ be the unitization of $\DA$, with
$(a1+F)^{\dagger}=\bar a1+F^{\dagger}$, and extend $M$ to $\DA^{+}$ by
$M(a1+F):=(a,a)+M(F)$.  Then:
\begin{enumerate}
\item[(a)] $M:\DA^{+}\to\BT$ is a surjective unital algebra homomorphism with
$M(X^{\dagger})=M(X)^{\sharp}$.
\item[(b)] For presentations $\mathfrak f,\mathfrak g$, the element
$M(s(\mathfrak f))\in\BT$ satisfies
$M(s(\mathfrak f\oplus\mathfrak g))=M(s(\mathfrak f))+M(s(\mathfrak g))$,
$M(s(\mathfrak f[1]))=-M(s(\mathfrak f))$,
$M(s(\mathfrak f\star\mathfrak g))=M(s(\mathfrak f))M(s(\mathfrak g))$ and
$M(s(\mathfrak f^{\dagger}))=M(s(\mathfrak f))^{\sharp}$.
\item[(c)] Let $I:=\ker M\subseteq\DA^{+}$.  With $\operatorname{Tor}$ and
$I^{2}$ taken over the ring underlying $\DA^{+}$, without topology,
\begin{equation}\label{eq:tor}
\operatorname{Tor}_1^{\DA^{+}}(\BT,\BT)\cong I/I^{2},
\end{equation}
and $I/I^{2}$ surjects onto $\C^{2}$.  In particular $\BT$ is not flat over
$\DA^{+}$.
\end{enumerate}
\end{theorem}

\begin{proof}
(a)  Multiplicativity and surjectivity of $M$ on $\DA$ are
Theorem~\ref{thm:tate}; the extension is multiplicative because $(1,1)$ is the
unit of $\BT$; and
$M(F^{\dagger})=(\overline{\lambda_+(F)},\overline{\lambda_-(F)})=M(F)^{\sharp}$,
$(a,a)^{\sharp}=(\bar a,\bar a)$.

(b)  Apply $M$ to the identities of Lemma~\ref{lem:slaws} and to
$s(\mathfrak f^{\dagger})=s(\mathfrak f)^{\dagger}$, and use (a).

(c)  For a commutative ring $R$ and an ideal $I$, tensoring
$0\to I\to R\to R/I\to0$ with $R/I$ gives the exact sequence
$0\to\operatorname{Tor}_1^{R}(R/I,R/I)\to I\otimes_RR/I\to R/I$, in which the
last map is zero since $I\cdot(R/I)=0$, and $I\otimes_RR/I=I/I^{2}$.  Apply
this with $R=\DA^{+}$ and $R/I=\BT$.

For $X=a1+F\in\DA^{+}$ put $\widehat X:=a+\widehat F$, so that
$(X\star Y)^{\wedge}=\widehat X\widehat Y$ and, by \eqref{eq:lambdaFT},
$M(X)=\bigl(\widehat X(-i/2),\widehat X(i/2)\bigr)$, where $\widehat X$ is
continued to an entire function of $\tau$ on the principal component
$\mathfrak c_0$.  Define $\partial_{\pm}:I\to\C$ by
$\partial_{\pm}(X):=\partial_{\tau}\widehat X(\mathfrak c_0,\pm i/2)$.  These
maps are linear.  If $X,Y\in I$ then
$\widehat X(\pm i/2)=\widehat Y(\pm i/2)=0$, so
$\partial_{\pm}(X\star Y)=\widehat X(\pm i/2)\partial_{\pm}Y+
\widehat Y(\pm i/2)\partial_{\pm}X=0$; hence $\partial_{\pm}$ vanish on
$I^{2}$.

It remains to show that $(\partial_+,\partial_-):I\to\C^{2}$ is surjective.
For $f\in C_c^{\infty}(\R)$ let $F\in\DA$ be the function
$F(w)=f(\log|w|_S)$; then $\widehat F$ is supported on $\mathfrak c_0$,
$\widehat F(\tau)=\int f(t)e^{-i\tau t}dt$, so
$\widehat F(\mp i/2)=\int f(t)e^{\mp t/2}dt$ and
$\partial_{\pm}(F)=-i\int tf(t)e^{\pm t/2}dt$.  The four functions
$e^{t/2},e^{-t/2},te^{t/2},te^{-t/2}$ are linearly independent on $\R$: if
$(a_1+a_2t)e^{t/2}+(a_3+a_4t)e^{-t/2}=0$ for all $t$, then letting
$t\to+\infty$ gives $a_1=a_2=0$ and letting $t\to-\infty$ gives $a_3=a_4=0$.
Hence the four functionals $f\mapsto\int fe^{t/2}$, $\int fe^{-t/2}$,
$\int tfe^{t/2}$, $\int tfe^{-t/2}$ on $C_c^{\infty}(\R)$ are linearly
independent, and the map $C_c^{\infty}(\R)\to\C^{4}$ which they define is
surjective.  The $f$ on which the first two functionals vanish give elements
$F\in I$, and on them the last two functionals take all values in $\C^{2}$.
So $I/I^{2}$ surjects onto $\C^{2}$.  If $\BT$ were flat over $\DA^{+}$, the
group $\operatorname{Tor}_1^{\DA^{+}}(\BT,\BT)$ would vanish.
\end{proof}

\section{The classical model}\label{sec:classical}

The composition law \eqref{eq:starmult} and the involution of
Definition~\ref{def:dagger} have an analogue for graphs of automorphisms of a
curve.

Let $C$ be a smooth projective geometrically connected curve of genus $g$ over
a field $k$, and put $Y=C\times C$ with projections $p_1,p_2$.  For
$K,L\in\Perf(Y)$ the composition is
$K\star L:=Rp_{13*}\bigl(p_{12}^{*}K\otimes^{L}p_{23}^{*}L\bigr)$, with
$p_{ij}$ the projections of $C^{3}$; it is associative and $\mathcal O_{\Delta}$
is a two-sided unit \cite[Ch.~5]{Huybrechts}.  For an automorphism $f$ of $C$
let $i_f:C\to Y$, $x\mapsto(x,f(x))$, with image the graph $\Gamma_f$.  We
assume that $C$ has a theta characteristic $h$ defined over $k$, that is, a
line bundle with $h^{2}\cong\omega_C$; this holds after a finite extension of
$k$.  Let $\sigma$ be the swap of the two factors of $Y$,
$K^{\vee}=R\!\mathcal{H}om(K,\mathcal O_Y)$, and
\[
D_h(K):=\sigma^{*}(K^{\vee})\otimes p_1^{*}h\otimes p_2^{*}h\,[1].
\]

\begin{theorem}\label{thm:classical}
Let $\Lambda$ be an abelian group and $\phi:\Lambda\to\Aut(C)$ a homomorphism.
For $\lambda\in\Lambda$ put $K_{\lambda}:=\mathcal O_{\Gamma_{\phi(\lambda)}}
\in\Perf(Y)$.
\begin{enumerate}
\item[(a)] $K_{\lambda}\star K_{\mu}\cong K_{\lambda\mu}$ for all
$\lambda,\mu\in\Lambda$.
\item[(b)] $K_1=\mathcal O_{\Delta}$ is the unit for $\star$.
\item[(c)] For every $\lambda$, writing $f=\phi(\lambda)$,
\begin{equation}\label{eq:dualitygraph}
D_h(K_{\lambda})\ \cong\
K_{\lambda^{-1}}\otimes p_1^{*}h\otimes p_2^{*}h^{-1}
\ \cong\ i_{f^{-1}*}\bigl(h\otimes(f^{-1})^{*}h^{-1}\bigr).
\end{equation}
The line bundle $h\otimes(f^{-1})^{*}h^{-1}$ has trivial square, and it is
trivial if and only if $f^{*}h\cong h$.  In particular
$D_h(\mathcal O_{\Delta})\cong\mathcal O_{\Delta}$.
\item[(d)] $\chi(L\Delta^{*}\mathcal O_{\Delta})=2-2g$, and
$L\Delta^{*}K_{\lambda}=0$ whenever $\phi(\lambda)$ has no fixed point over an
algebraic closure of $k$.
\end{enumerate}
\end{theorem}

\begin{proof}
(a)  Let $f,f'$ be automorphisms of $C$.  The divisors $\Gamma_f\times C$ and
$C\times\Gamma_{f'}$ of $C^{3}$ are integral and distinct, so a local equation
of the first is a nonzerodivisor on the second, and
$p_{12}^{*}\mathcal O_{\Gamma_f}\otimes^{L}p_{23}^{*}\mathcal O_{\Gamma_{f'}}$
is the structure sheaf of their intersection
$Z=\{(x,f(x),f'(f(x)))\}$.  The projection $p_{13}$ maps $Z$ isomorphically
onto $\Gamma_{f'\circ f}$, so
$\mathcal O_{\Gamma_f}\star\mathcal O_{\Gamma_{f'}}\cong
\mathcal O_{\Gamma_{f'\circ f}}$.  For $f=\phi(\lambda)$, $f'=\phi(\mu)$ one
has $f'\circ f=\phi(\mu\lambda)=\phi(\lambda\mu)$, because $\Lambda$ is
abelian.

(b)  $\phi(1)=\id$ and $\Gamma_{\id}=\Delta$.

(c)  The map $i_f$ is a closed immersion of a smooth curve in a smooth
surface.  Its normal bundle is $i_f^{*}T_Y/T_C=(T_C\oplus f^{*}T_C)/T_C\cong
f^{*}T_C$, and $f^{*}T_C\cong T_C$ by the differential of $f$.  By duality for
the closed immersion $i_f$ and the fundamental local isomorphism
\cite[Ch.~III, \S\S6--7]{HartshorneRD}, for a line bundle $N$ on $C$,
\begin{equation}\label{eq:fli}
\bigl(i_{f*}N\bigr)^{\vee}\ \cong\ i_{f*}\bigl(N^{-1}\otimes T_C\bigr)[-1].
\end{equation}
Next, $\sigma\circ i_f=i_{f^{-1}}\circ f$, because
$\sigma(x,f(x))=(f(x),x)=i_{f^{-1}}(f(x))$; since $\sigma$ is an involution,
$\sigma^{*}=\sigma_{*}$, and so
\[
\sigma^{*}i_{f*}(N)\cong i_{f^{-1}*}\bigl(f_{*}N\bigr)
=i_{f^{-1}*}\bigl((f^{-1})^{*}N\bigr).
\]
Finally $i_{f^{-1}}^{*}(p_1^{*}h\otimes p_2^{*}h)=h\otimes(f^{-1})^{*}h$.
Taking $N=\mathcal O_C$ and using the projection formula,
\[
D_h(\mathcal O_{\Gamma_f})
\cong i_{f^{-1}*}\Bigl((f^{-1})^{*}T_C\otimes h\otimes(f^{-1})^{*}h\Bigr)
\cong i_{f^{-1}*}\Bigl(h\otimes(f^{-1})^{*}\bigl(T_C\otimes h\bigr)\Bigr).
\]
Since $h^{2}\cong\omega_C$ we have $T_C\otimes h\cong h^{-1}$, which gives the
second expression in \eqref{eq:dualitygraph}; the first is equal to it by the
projection formula, since $i_{f^{-1}}^{*}(p_1^{*}h\otimes p_2^{*}h^{-1})
=h\otimes(f^{-1})^{*}h^{-1}$.  The square of this bundle is
$\omega_C\otimes(f^{-1})^{*}\omega_C^{-1}\cong\mathcal O_C$, and it is trivial
if and only if $(f^{-1})^{*}h\cong h$, that is, $f^{*}h\cong h$.

(d)  The diagonal is a Cartier divisor on $Y$ with $\mathcal O_Y(-\Delta)
|_{\Delta}\cong\Omega^{1}_C$, so $L\Delta^{*}\mathcal O_{\Delta}$ is the
complex $\Omega^{1}_C\xrightarrow{\,0\,}\mathcal O_C$, that is,
$\mathcal O_C\oplus\Omega^{1}_C[1]$, and
$\chi=(1-g)-(g-1)=2-2g$.  If $\phi(\lambda)$ has no fixed point then
$\Gamma_{\phi(\lambda)}\cap\Delta=\emptyset$ and the restriction of
$K_{\lambda}$ to $\Delta$ is zero.
\end{proof}

\begin{remark}\label{rem:classicalscope}
By (a) and (b), $\lambda\mapsto K_{\lambda}$ has the properties
\eqref{eq:starmult} of $a\mapsto E_a$.  By (c), $D_h(K_{\lambda})$ differs
from $K_{\lambda^{-1}}$, the analogue of $E_a^{\dagger}=E_{a^{-1}}$, by the
twist $L_{\lambda}:=h\otimes(\phi(\lambda)^{-1})^{*}h^{-1}$.  For $g\ne1$ a
nontrivial twist cannot be removed by replacing $K_{\lambda}$ with
$K_{\lambda}^{a,b}:=i_{f*}(h^{a}\otimes f^{*}h^{b})$, $f=\phi(\lambda)$,
$a,b\in\Z$.  Indeed, with $f'=f^{-1}$ and
$L:=(f')^{*}h\otimes h^{-1}\cong L_{\lambda}^{-1}\cong L_{\lambda}$, formula
\eqref{eq:fli} and the computation in the proof of (c) give
\[
D_h\bigl(K_{\lambda}^{a,b}\bigr)\cong
i_{f'*}\bigl(h^{1-b}\otimes(f')^{*}h^{-a-1}\bigr)
=i_{f'*}\bigl(h^{-a-b}\otimes L^{-a-1}\bigr),\qquad
K_{\lambda^{-1}}^{a,b}=i_{f'*}\bigl(h^{a+b}\otimes L^{b}\bigr).
\]
Since $L^{2}\cong\mathcal O_C$, the two are isomorphic if and only if
$\omega_C^{a+b}\cong L^{a+b+1}$.  For $g\ne1$, comparing degrees gives
$a+b=0$, and the condition becomes $L\cong\mathcal O_C$, that is,
$L_{\lambda}\cong\mathcal O_C$.  The duality $D_h$
is $k$-linear, whereas $\dagger$ is antilinear on morphisms; the comparison
concerns objects and the composition law.  Graphs of maps of degree greater
than one have no analogue among the $E_a$.
\end{remark}

\begin{remark}\label{rem:invertible}
For $a\in C_S$ the bimodule $E_a$ has the bidegree
\begin{equation}\label{eq:bidegree}
\rho_T(a)=\bigl(|a|_S^{-1/2},|a|_S^{+1/2}\bigr)\in\BT
\end{equation}
of \eqref{eq:rhoT}, in the sense that $M(F_w)=\rho_T(w)M(F)$ for the translate
$F_w(v)=F(w^{-1}v)$ of $F\in\DA$.  One has
$E_a\star E_a^{\dagger}=E_a\star E_{a^{-1}}=E_1=\Delta$ by \eqref{eq:starmult}
and Definition~\ref{def:dagger}, and $\beta_T(\rho_T(a),\rho_T(a))=2$ for
every $a$.
\end{remark}

\section{The signature of the form}\label{sec:signature}

A Hermitian form has \emph{infinite positive} (respectively \emph{negative})
\emph{index} if it is positive (respectively negative) definite on some
infinite-dimensional subspace.

\begin{theorem}\label{thm:signature}
For $F\in\DA$,
\begin{equation}\label{eq:IAF}
I_A(F,F)=T\bigl(F\star F^{\dagger}\bigr)
=\int_{\Ch}|\widehat F(\chi)|^{2}\bigl(\Th(\chi)+c\bigr)d\chi ,
\end{equation}
a real number.  Let $m:=\sup_{\Ch}\Th<\infty$
(Theorem~\ref{thm:symbol}(c)).  Then:
\begin{enumerate}
\item[(a)] $I_A$ has infinite negative index, for every $c$.
\item[(b)] If $c>-m$ then $I_A$ has infinite positive index.  If $c\le-m$
then $I_A$ is negative definite on $\DA$.
\item[(c)] One has
\begin{equation}\label{eq:mbound}
m\ \ge\ \log\pi-\psi(\tfrac14)+\sum_{p\in S_f}(\log p)
\frac{2p^{-1/2}}{1-p^{-1/2}}\ >\ \log\pi+\gamma+3\log2+\frac{\pi}{2}
\ >\ 5.37 .
\end{equation}
In particular $m>0$, and for $c=0$ the form has infinite index of both signs.
\item[(d)] The statements (a) and (b) hold for the restriction of $I_A$ to
$\DA^{\mathrm{prim}}$.
\end{enumerate}
\end{theorem}

\begin{proof}
Formula \eqref{eq:IAF} is \eqref{eq:plancherel} applied to $H=F\star
F^{\dagger}$, whose transform is $|\widehat F|^{2}$; the value is real because
$\Th+c$ is.

Fix a component $\mathfrak c$ of $\Ch$.  For $\varphi\in C_c^{\infty}(\R)$ let
$F[\varphi]\in\DA$ be the function, constructed as in the proof of
Theorem~\ref{thm:radical}, whose transform is supported on $\mathfrak c$ and
equals $\widehat\varphi(\tau)=\int\varphi(t)e^{-i\tau t}dt$ there.  Then
\[
I_A(F[\varphi],F[\psi])=\int_{\R}\widehat\varphi(\tau)\,
\overline{\widehat\psi(\tau)}\,(\Th(\mathfrak c,\tau)+c)\,\frac{d\tau}{2\pi}.
\]
We write $\lVert g\rVert^{2}:=\int_{\R}|g|^{2}\frac{d\tau}{2\pi}$ and
$\omega(\tau):=|\Th(\mathfrak c,\tau)|+|c|$.  By \eqref{eq:logbound},
$\omega(\tau)\le A_{\mathfrak c}+|c|+B_S\log(2+|\tau|)$.

\emph{Concentration.}  Fix $\varphi_0\in C_c^{\infty}(\R)$ with
$\lVert\widehat\varphi_0\rVert=1$; $\widehat\varphi_0$ is a Schwartz function.
For $a\in\R$ and $\delta\ge1$ put
$\varphi_{a,\delta}(t):=\delta^{-1/2}e^{iat}\varphi_0(t/\delta)$, so that
$\widehat\varphi_{a,\delta}(\tau)=\delta^{1/2}\widehat\varphi_0(\delta(\tau-a))$
and $\lVert\widehat\varphi_{a,\delta}\rVert=1$.  For every $\rho>0$, the
substitution $s=\delta(\tau-a)$ gives
\begin{equation}\label{eq:weightedtail}
\int_{|\tau-a|>\rho}\omega(\tau)
|\widehat\varphi_{a,\delta}(\tau)|^{2}\frac{d\tau}{2\pi}
\le\int_{|s|>\delta\rho}\Bigl(A_{\mathfrak c}+|c|+B_S\log\bigl(2+|a|+|s|\bigr)\Bigr)
|\widehat\varphi_0(s)|^{2}\frac{ds}{2\pi}\ \longrightarrow\ 0
\end{equation}
as $\delta\to\infty$, the integrand on the right being integrable over $\R$.
In particular $\int_{|\tau-a|\le\rho}|\widehat\varphi_{a,\delta}|^{2}
\frac{d\tau}{2\pi}\to1$.

\emph{A Gram estimate.}  Let $\psi_1,\psi_2,\dots\in C_c^{\infty}(\R)$ with
$\lVert\widehat\psi_j\rVert=1$, and let $I_1,I_2,\dots$ be pairwise disjoint
intervals such that $\sup_{I_j}\omega\le L_j$,
$\int_{\R\setminus I_j}\omega|\widehat\psi_j|^{2}\frac{d\tau}{2\pi}\le
\varepsilon_j$, and, for $j<k$,
$\int_{I_k}\omega|\widehat\psi_j|^{2}\frac{d\tau}{2\pi}\le\eta_k$.  Then for
$j<k$, splitting the integral over $I_j$, $I_k$ and the complement of
$I_j\cup I_k$ and applying the Cauchy--Schwarz inequality on each piece,
\begin{equation}\label{eq:gram}
|I_A(F[\psi_j],F[\psi_k])|\le\int\omega|\widehat\psi_j||\widehat\psi_k|
\frac{d\tau}{2\pi}\le(L_j\varepsilon_k)^{1/2}+(\eta_kL_k)^{1/2}
+(\varepsilon_j\varepsilon_k)^{1/2}.
\end{equation}
Such sequences are constructed inductively.  At step $k$, with
$\psi_1,\dots,\psi_{k-1}$ and $I_1,\dots,I_{k-1}$ given, one first chooses an
interval $I_k$ disjoint from the earlier ones in which each of the functions
$\omega|\widehat\psi_j|^{2}$, $j<k$, has integral at most $\eta_k$.  This is
possible far from the origin, because these functions are integrable, and it
is possible with $I_k$ short inside a bounded set, because they are bounded on
bounded sets.  Then, with $a_k$ the midpoint of $I_k$, one takes
$\psi_k=\varphi_{a_k,\delta_k}$ with $\delta_k$ so large, by
\eqref{eq:weightedtail}, that
$\int_{\R\setminus I_k}\omega|\widehat\psi_k|^{2}\frac{d\tau}{2\pi}\le
\varepsilon_k$ and
$\int_{I_k}|\widehat\psi_k|^{2}\frac{d\tau}{2\pi}\ge\tfrac12$.

\emph{Proof of} (a).  By Theorem~\ref{thm:symbol}(c), $\Th(\mathfrak c,\cdot)+c$
is continuous and tends to $-\infty$.  Put $K_0:=\max(2|c|,1)$ and
$\varepsilon_k=\eta_k:=K_0\,2^{-4k-12}$.  At step $k$ choose a point $a_k$ so
far from the origin that $\Th(\mathfrak c,a_k)+c\le-\tfrac32K_0$, that
$[a_k-1,a_k+1]$ is disjoint from the earlier intervals, and that
$\int_{a_k-1}^{a_k+1}\omega|\widehat\psi_j|^{2}\frac{d\tau}{2\pi}\le\eta_k$ for
$j<k$.  Put $K_k:=-\tfrac23(\Th(\mathfrak c,a_k)+c)\ge K_0$, and by continuity
choose $\rho_k\in(0,1]$ with
$-2K_k\le\Th(\mathfrak c,\cdot)+c\le-K_k$ on
$I_k:=(a_k-\rho_k,a_k+\rho_k)$.  Then
$\omega\le|\Th+c|+2|c|\le2K_k+K_0\le3K_k=:L_k$ on $I_k$.  With
$u_k:=F[\psi_k]/\sqrt{K_k}$ and $G_{jk}:=I_A(u_j,u_k)$,
\[
G_{kk}\le K_k^{-1}\bigl(-\tfrac12K_k+\varepsilon_k\bigr)\le-\tfrac12+2^{-16},
\]
and by \eqref{eq:gram}, for $j<k$, using $K_j,K_k\ge K_0$,
\[
|G_{jk}|\le\sqrt3\Bigl(\frac{\varepsilon_k}{K_k}\Bigr)^{1/2}
+\sqrt3\Bigl(\frac{\eta_k}{K_j}\Bigr)^{1/2}
+\Bigl(\frac{\varepsilon_j\varepsilon_k}{K_jK_k}\Bigr)^{1/2}
\le(2\sqrt3+1)\,2^{-2k-6}\le2^{-j-k-2}.
\]
For a finite combination $\sum_jc_ju_j$,
\[
I_A\Bigl(\sum c_ju_j,\sum c_ju_j\Bigr)\le\bigl(-\tfrac12+2^{-16}\bigr)\sum|c_j|^{2}
+\sum_{j\ne k}2^{-j-k-2}|c_j||c_k|
\le\Bigl(-\tfrac12+2^{-16}+\tfrac1{12}\Bigr)\sum|c_j|^{2},
\]
using $\sum_{j\ne k}2^{-j-k}|c_j||c_k|\le\bigl(\sum_j2^{-j}|c_j|\bigr)^{2}\le
\bigl(\sum_{j\ge1}4^{-j}\bigr)\sum|c_j|^{2}=\tfrac13\sum|c_j|^{2}$.  So $I_A$ is
negative definite on the span of the $u_j$, and the same inequality shows that
the $u_j$ are linearly independent.

\emph{Proof of} (b).  If $c\le-m$ then $\Th+c\le0$ on $\Ch$, and by
Theorem~\ref{thm:symbol}(c) its zero set is discrete in each component.  For
$F\ne0$ the continuous function $|\widehat F|^{2}$ is positive on a nonempty
open set, so \eqref{eq:IAF} gives $I_A(F,F)<0$.  If $c>-m$, choose a component
$\mathfrak c$ and an interval $I=(a-\rho,a+\rho)$ with
$\Th(\mathfrak c,\cdot)+c\ge\kappa>0$ on $I$, and put $L:=\sup_I\omega$.
Run the inductive construction with $I_k$ a subinterval of
$(a+2^{-k}\rho,a+2^{-k+1}\rho)$, $L_k:=L$ and
$\varepsilon_k=\eta_k:=\kappa^{2}2^{-4k-12}/(L+\kappa)$.  Then
$I_A(F[\psi_k],F[\psi_k])\ge\tfrac12\kappa-\varepsilon_k\ge\tfrac14\kappa$,
and by \eqref{eq:gram}, $|I_A(F[\psi_j],F[\psi_k])|\le3\kappa\,2^{-2k-6}\le
\kappa\,2^{-j-k-4}$ for $j<k$.  As in (a),
\[
I_A\Bigl(\sum c_jF[\psi_j],\sum c_jF[\psi_j]\Bigr)\ge
\bigl(\tfrac14-\tfrac1{48}\bigr)\kappa\sum|c_j|^{2},
\]
so $I_A$ is positive definite on the span of the $F[\psi_j]$, which is
infinite dimensional.

(c)  Evaluate $\Th$ at $\tau=0$ on the principal component.  There
$\zeta_p=p^{-1/2}$, so by \eqref{eq:Wphat} the term at $p$ is
$(\log p)\cdot2p^{-1/2}/(1-p^{-1/2})>0$; and the archimedean term is
\eqref{eq:maxarch}, with $\psi(\tfrac14)=-\gamma-3\log2-\pi/2$.  Numerically
$\log\pi+\gamma+3\log2+\pi/2=5.372\ldots$

(d)  $\DA^{\mathrm{prim}}$ has codimension two in $\DA$, so its intersection
with an infinite-dimensional subspace on which $I_A$ is definite is
infinite dimensional, and $I_A$ is definite on it.  The second statement of
(b) restricts to any subspace.
\end{proof}

The argument of (d) applies to every subspace of finite codimension.  Hence
for $c>-m$ the form $I_A$ is semidefinite on no subspace of finite codimension
of $\DA$; in particular $\DA$ is not the orthogonal sum of a plane and a
subspace on which $I_A$ is definite, which is the decomposition that the Hodge
index theorem provides for divisors on a surface.  Moreover the
$I_A$-orthogonal complement of $\DA^{\mathrm{prim}}$ is zero
(Theorem~\ref{thm:nosplit}(b)).

\appendix

\section{The rulings in the limit over \texorpdfstring{$S$}{S}}\label{app:mertens}

This appendix uses Mertens' theorem.  Nothing in
Sections~\ref{sec:carrier}--\ref{sec:signature} depends on it.

Theorem~\ref{thm:scalartrivial} concerns a fixed $S$.  We now compute the
form $I_{\infty}$ of Theorem~\ref{thm:colimtrace} on approximations of the
rulings.  Call $F\in\DA(S)$ \emph{radial with profile} $f\in C_c^{\infty}(\R)$
if $F(w)=f(\log|w|_S)$.  By \eqref{eq:chartsystem}, $\Inf_{SS'}F$ is then
radial with the same profile, so radial elements of $\DA^{\infty}$ and their
profiles are well defined.  Since $\dstar w=dt\otimes\mu_{\mathrm{cpt}}$ with
$\mu_{\mathrm{cpt}}$ a probability measure, the convolution of radial elements
with profiles $f,g$ is radial with profile $f*g$, the additive convolution on
$\R$, and $F^{\dagger}$ is radial with profile $\overline{f(-t)}$.

\begin{lemma}\label{lem:radialtrace}
Let $H\in\DA^{\infty}$ be radial with profile $h$.  Then
\begin{equation}\label{eq:Wpprofile}
T_{\infty}(H)=W_{\infty}(H)+c\,h(0)+\sum_{p}(\log p)\sum_{k\ne0}p^{-|k|/2}
h(-k\log p),
\end{equation}
the sum over all primes $p$ being finite, and
\[
W_{\infty}(H)=(\gamma+\log4\pi)\,h(0)
+\int_0^{\infty}K(t)\bigl[h(t)+h(-t)-2h(0)e^{-t/2}\bigr]dt .
\]
\end{lemma}

\begin{proof}
Represent $H$ in $\DA(S)$.  For $p\in S_f$, by \eqref{eq:slices} the element
$j_p(p^{k}u)$ has archimedean coordinate $p^{-k}$, so
$H(j_p(p^{k}u))=h(-k\log p)$ for every unit $u$.  The integral over the units
in \eqref{eq:Wp} vanishes, and Lemma~\ref{lem:tau}(i) with the shell volume
$\log p$ gives $W_p(H)=(\log p)\sum_{k\ne0}p^{-|k|/2}h(-k\log p)$.  For
$p\notin S_f$ the term at $p$ in \eqref{eq:Tinf} is the same expression,
because $|\delta_S(p^{-k})|_S=p^{-k}$; and it vanishes when $\log p$ exceeds
the radius of the support of $h$.  The two profiles $b_{\pm}$ of
Definition~\ref{def:trace} are both equal to $h$, so $W_{\infty}(H)$ is given
by \eqref{eq:Winfeven}.  Finally $H(1)=h(0)$.
\end{proof}

Let $\Lambda(N)$ be the von Mangoldt function.  By Mertens' theorem
\cite{Mertens}, \cite[Ch.~2]{MontgomeryVaughan}, there is a constant $C_0$
such that
\begin{equation}\label{eq:mertenssum}
\sum_{N\le x}\frac{\Lambda(N)}{N}=\log x+E(x),\qquad|E(x)|\le C_0\quad(x\ge1).
\end{equation}

\begin{theorem}\label{thm:mertens}
Let $c$ be arbitrary.
\begin{enumerate}
\item[(a)] Let $\varphi_1,\varphi_2\in C_c^{\infty}(\R)$ with
$\supp\varphi_i\subseteq[A_i,\infty)$ and $A:=A_1+A_2>0$, and let
$F_i\in\DA^{\infty}$ be radial with profile $e^{-t/2}\varphi_i(t)$.  Then
\begin{equation}\label{eq:step2}
\Bigl|T_{\infty}(F_1\star F_2)-\int\varphi_1\int\varphi_2\Bigr|
\le C_0\min\bigl(\lVert\varphi_1'\rVert_{L^1}\lVert\varphi_2\rVert_{L^1},
\lVert\varphi_1\rVert_{L^1}\lVert\varphi_2'\rVert_{L^1}\bigr)
+\frac{e^{-A}}{1-e^{-2A}}\lVert\varphi_1\rVert_{L^1}\lVert\varphi_2\rVert_{L^1}.
\end{equation}
\item[(b)] Let $F_n,G_m$ be the functions \eqref{eq:spread}, for a real
$\vartheta\in C_c^{\infty}((0,\infty))$ with $\vartheta\ge0$,
$\int\vartheta=1$ and $\supp\vartheta\subseteq[a,b]$, regarded as elements of
$\DA^{\infty}$.  Then, with constants depending only on $\vartheta$ and $c$,
\begin{equation}\label{eq:mertenslimit}
I_{\infty}(F_n,G_m)=1+O\bigl(\max(n,m)^{-1}\bigr),\qquad
I_{\infty}(F_n,F_n)=O(e^{-na}),\qquad I_{\infty}(G_m,G_m)=O(e^{-ma}).
\end{equation}
So the Gram matrix of $I_{\infty}$ on $(G_n,F_n)$ tends to
$\left(\begin{smallmatrix}0&1\\1&0\end{smallmatrix}\right)$, the matrix of
$\beta_T$ in the basis $(e_-,e_+)$.
\item[(c)] There are $\Phi_n\in\DA(S)$, $n\ge1$, such that $(\Phi_n)$ is an
admissible approximation of $e_+$ and $(\Phi_n^{\dagger})$ one of $e_-$
(Definition~\ref{def:admissible}), and $I_{\infty}(\Phi_n,\Phi_n^{\dagger})\to0$.
\end{enumerate}
\end{theorem}

Thus the value $1$ in (b) is a property of approximations whose profiles vary
slowly, in the sense that the right-hand side of \eqref{eq:step2} is small; by
(c) it is not shared by all admissible approximations.

\begin{proof}
(a)  The element $F_1\star F_2$ is radial with profile
\begin{equation}\label{eq:profileconv}
h(t)=\bigl(e^{-t/2}\varphi_1\bigr)*\bigl(e^{-t/2}\varphi_2\bigr)(t)
=e^{-t/2}\varphi(t),\qquad\varphi:=\varphi_1*\varphi_2 ,
\end{equation}
and $\varphi\in C_c^{\infty}(\R)$ is supported in $[A,\infty)$, with
$\int\varphi=\int\varphi_1\int\varphi_2$ and
\[
\lVert\varphi'\rVert_{L^1}\le\min\bigl(\lVert\varphi_1'\rVert_{L^1}
\lVert\varphi_2\rVert_{L^1},\lVert\varphi_1\rVert_{L^1}
\lVert\varphi_2'\rVert_{L^1}\bigr).
\]
In particular $h(0)=0$ and $h(-t)=0$ for
$t>0$.  We apply Lemma~\ref{lem:radialtrace}.

\emph{Finite places.}  Only $k=-j$ with $j\ge1$ contributes to
\eqref{eq:Wpprofile}, and
$p^{-j/2}h(j\log p)=p^{-j}\varphi(j\log p)$.  With $N=p^{j}$,
\[
\sum_p(\log p)\sum_{k\ne0}p^{-|k|/2}h(-k\log p)
=\sum_{N\ge2}\frac{\Lambda(N)}{N}\,\varphi(\log N).
\]
By partial summation and \eqref{eq:mertenssum}, $\varphi$ being smooth with
compact support in $(0,\infty)$,
\begin{align*}
\sum_{N\ge2}\frac{\Lambda(N)}{N}\,\varphi(\log N)
&=-\int_1^{\infty}\bigl(\log x+E(x)\bigr)\varphi'(\log x)\,\frac{dx}{x}
=-\int_0^{\infty}\bigl(t+E(e^{t})\bigr)\varphi'(t)\,dt\\
&=\int\varphi-\int_0^{\infty}E(e^{t})\varphi'(t)\,dt ,
\end{align*}
and the last integral is at most $C_0\lVert\varphi'\rVert_{L^1}$ in absolute
value.

\emph{The real place.}  Here $W_{\infty}(F_1\star F_2)=\int_A^{\infty}K(t)h(t)
\,dt$.  For $t\ge A$ one has
$K(t)=e^{-t/2}(1-e^{-2t})^{-1}\le(1-e^{-2A})^{-1}e^{-t/2}$, so
\[
|W_{\infty}(F_1\star F_2)|\le(1-e^{-2A})^{-1}\int_A^{\infty}
e^{-t}|\varphi(t)|\,dt\le\frac{e^{-A}}{1-e^{-2A}}
\lVert\varphi_1\rVert_{L^1}\lVert\varphi_2\rVert_{L^1}.
\]
Adding the two estimates gives \eqref{eq:step2}.

(b)  By \eqref{eq:chartsystem} the functions \eqref{eq:spread} for different
$S$ correspond under inflation.  $F_n$ is radial with profile
$f_n(t)=e^{-t/2}\vartheta_n(t)$, where
$\vartheta_n(t):=n^{-1}\vartheta(t/n)$ is a probability density supported in
$[na,nb]$ with $\lVert\vartheta_n'\rVert_{L^1}=\lVert\vartheta'\rVert_{L^1}/n$.
The profile of $G_m$ is $m^{-1}e^{t/2}\vartheta(-t/m)=\overline{f_m(-t)}$, so
$G_m=F_m^{\dagger}$ and $I_{\infty}(F_n,G_m)=T_{\infty}(F_n\star F_m)$.  By (a)
with $A=(n+m)a\ge2a$,
\[
\bigl|I_{\infty}(F_n,G_m)-1\bigr|\le
\frac{C_0\lVert\vartheta'\rVert_{L^1}}{\max(n,m)}
+\frac{e^{-(n+m)a}}{1-e^{-4a}} .
\]

For the second estimate put $H:=F_n\star F_n^{\dagger}$, radial with profile
\[
\eta(t)=\int f_n(u)\overline{f_n(u-t)}\,du=e^{t/2}B(t),\qquad
B(t):=n^{-1}\int_{\R}e^{-ns}\vartheta(s)\vartheta(s-t/n)\,ds ,
\]
by the substitution $u=ns$.  The integrand vanishes unless $s\ge a$ and
$s-t/n\ge a$, and $\eta$ is supported in $|t|\le n(b-a)$.  For $t>0$ this gives
$e^{-ns}\le e^{-na-t}$, and for $t\le0$ it gives $e^{-ns}\le e^{-na}$; since
$\int\vartheta(s)\vartheta(s-t/n)\,ds\le\lVert\vartheta\rVert_{\infty}$,
\begin{equation}\label{eq:hprime}
|B(t)|\le\lVert\vartheta\rVert_{\infty}n^{-1}e^{-na}\min(1,e^{-t}),\qquad
|\eta(t)|\le\lVert\vartheta\rVert_{\infty}n^{-1}e^{-na}e^{-|t|/2}.
\end{equation}
The same argument applied to
$B'(t)=-n^{-2}\int e^{-ns}\vartheta(s)\vartheta'(s-t/n)\,ds$ and
$\eta'=\tfrac12\eta+e^{t/2}B'$ gives
$|\eta'(t)|\le C_{\vartheta}n^{-1}e^{-na}e^{-|t|/2}$ with
$C_{\vartheta}:=\lVert\vartheta\rVert_{\infty}+\lVert\vartheta'\rVert_{\infty}$.
We estimate the three parts of \eqref{eq:Wpprofile} for $H$.  First,
$|c\,\eta(0)|\le|c|\lVert\vartheta\rVert_{\infty}n^{-1}e^{-na}$.  Second, for
$k=-j<0$, $p^{-j/2}|\eta(j\log p)|=|B(j\log p)|\le
\lVert\vartheta\rVert_{\infty}n^{-1}e^{-na}p^{-j}$, and for $k=j>0$,
$p^{-j/2}|\eta(-j\log p)|=p^{-j}|B(-j\log p)|\le
\lVert\vartheta\rVert_{\infty}n^{-1}e^{-na}p^{-j}$; only $p^{j}\le e^{n(b-a)}$
occurs, so by \eqref{eq:mertenssum}
\[
\sum_p(\log p)\sum_{k\ne0}p^{-|k|/2}|\eta(-k\log p)|\le
2\lVert\vartheta\rVert_{\infty}n^{-1}e^{-na}
\sum_{N\le e^{n(b-a)}}\frac{\Lambda(N)}{N}
\le2\lVert\vartheta\rVert_{\infty}e^{-na}\bigl(b-a+C_0\bigr).
\]
Third, in the formula for $W_{\infty}(H)$ of Lemma~\ref{lem:radialtrace}, for
$0<t\le1$ one has $K(t)\le1/t$ and
\[
\bigl|\eta(t)+\eta(-t)-2\eta(0)e^{-t/2}\bigr|\le
2t\sup|\eta'|+t\,|\eta(0)| ,
\]
while for $t\ge1$ one has $K(t)\le(1-e^{-2})^{-1}e^{-t/2}$ and the bracket is
at most $4\sup_{s}\bigl(|\eta(s)|e^{|s|/2}\bigr)e^{-t/2}$.  By
\eqref{eq:hprime} and the bound for $\eta'$, $W_{\infty}(H)=O(n^{-1}e^{-na})$.
Hence $I_{\infty}(F_n,F_n)=T_{\infty}(H)=O(e^{-na})$.  Finally
$G_m\star G_m^{\dagger}=F_m^{\dagger}\star F_m$, so
$I_{\infty}(G_m,G_m)=I_{\infty}(F_m,F_m)$; and
$I_{\infty}(G_m,F_n)=\overline{I_{\infty}(F_n,G_m)}$.

(c)  Let $\rho\in C_c^{\infty}((-1,1))$ with $\rho\ge0$ and $\int\rho=1$, and
put $t_n:=n\log6$, $\varepsilon_n:=36^{-n}/4$,
$\rho_n(t):=\varepsilon_n^{-1}\rho((t-t_n)/\varepsilon_n)$.  Let $\Phi_n$ be
radial with profile $e^{-t/2}\rho_n(t)$.  Then $\Phi_n\ge0$,
$\lambda_+(\Phi_n)=\int\rho_n=1$,
$\lambda_-(\Phi_n)=\int e^{-t}\rho_n\le e^{-t_n+1}$ and
$\lVert \Phi_n\rVert_{L^1}=\int e^{-t/2}\rho_n\le e^{-(t_n-1)/2}$; so $(\Phi_n)$ is
an admissible approximation of $e_+$, and $(\Phi_n^{\dagger})$ one of $e_-$,
because $M(F^{\dagger})=M(F)^{\sharp}$ and
$\lVert F^{\dagger}\rVert_{L^1}=\lVert F\rVert_{L^1}$.  Now
$I_{\infty}(\Phi_n,\Phi_n^{\dagger})=T_{\infty}(\Phi_n\star \Phi_n)$, and
$\rho_n*\rho_n$ is supported in the open interval
$|t-\log36^{n}|<36^{-n}/2$.  For an integer $N\ne36^{n}$ one has
$|\log N-\log36^{n}|\ge\log(1+36^{-n})\ge36^{-n}/2$, and $36^{n}$ is not a
prime power.  So $(\rho_n*\rho_n)(\log N)=0$ for every prime power $N$, the
sum over the primes in \eqref{eq:Wpprofile} vanishes for $\Phi_n\star \Phi_n$, and
the value at $0$ of its profile is $0$.  By the estimate of the real place in
(a), with $A=2(t_n-\varepsilon_n)$,
$|T_{\infty}(\Phi_n\star \Phi_n)|\le e^{-A}(1-e^{-2A})^{-1}\to0$.
\end{proof}

\begin{remark}\label{rem:mertensmeaning}
The value $1$ in \eqref{eq:mertenslimit} comes from the main term $\log x$ of
\eqref{eq:mertenssum}.  For a fixed $S$ only the primes of $S$ contribute to
$T_S$, the corresponding sum is bounded, and one obtains the limit $0$ of
Theorem~\ref{thm:rulnull}(b).
\end{remark}


\begin{thebibliography}{99}

\bibitem{BondalKapranov}
A.~I.~Bondal and M.~M.~Kapranov, \emph{Enhanced triangulated categories},
Mat. Sb. \textbf{181} (1990), no.~5, 669--683; English translation in
Math. USSR-Sb. \textbf{70} (1991), no.~1, 93--107.

\bibitem{BrodzkiLykova}
J.~Brodzki and Z.~A.~Lykova, \emph{Excision in cyclic type homology of
Fr\'echet algebras}, Bull. London Math. Soc. \textbf{33} (2001), 283--291.

\bibitem{Brylinski}
J.-L.~Brylinski, \emph{Algebras associated with group actions and their
homology}, Brown University preprint, 1987.

\bibitem{Connes99}
A.~Connes, \emph{Trace formula in noncommutative geometry and the zeros of
the Riemann zeta function}, Selecta Math. (N.S.) \textbf{5} (1999), 29--106.

\bibitem{CCWeilArch}
A.~Connes and C.~Consani, \emph{Weil positivity and trace formula: the
archimedean place}, Selecta Math. (N.S.) \textbf{27} (2021), no.~4, Paper
No.~77.

\bibitem{DLMF}
\emph{NIST Digital Library of Mathematical Functions},
\texttt{https://dlmf.nist.gov/}, Chapters 4 and 5.

\bibitem{GetzlerJones}
E.~Getzler and J.~D.~S.~Jones, \emph{The cyclic homology of crossed product
algebras}, J. Reine Angew. Math. \textbf{445} (1993), 161--174.

\bibitem{GR}
I.~S.~Gradshteyn and I.~M.~Ryzhik, \emph{Table of Integrals, Series, and
Products}, 7th ed., Academic Press, 2007.

\bibitem{HartshorneRD}
R.~Hartshorne, \emph{Residues and Duality}, Lecture Notes in Math.
\textbf{20}, Springer, 1966.

\bibitem{Helemskii}
A.~Ya.~Helemskii, \emph{The Homology of Banach and Topological Algebras},
Mathematics and its Applications (Soviet Series) \textbf{41}, Kluwer, 1989.

\bibitem{Hormander}
L.~H\"ormander, \emph{The Analysis of Linear Partial Differential Operators
I}, 2nd ed., Springer, 1990.

\bibitem{Huybrechts}
D.~Huybrechts, \emph{Fourier--Mukai Transforms in Algebraic Geometry}, Oxford
Mathematical Monographs, Oxford University Press, 2006.

\bibitem{MacLane}
S.~Mac~Lane, \emph{Homology}, Grundlehren der mathematischen Wissenschaften
\textbf{114}, Springer, 1963.

\bibitem{Mertens}
F.~Mertens, \emph{Ein Beitrag zur analytischen Zahlentheorie},
J. Reine Angew. Math. \textbf{78} (1874), 46--62.

\bibitem{MeyerIdele}
R.~Meyer, \emph{On a representation of the idele class group related to
primes and zeros of $L$-functions}, Duke Math. J. \textbf{127} (2005),
519--595.

\bibitem{MontgomeryVaughan}
H.~L.~Montgomery and R.~C.~Vaughan, \emph{Multiplicative Number Theory I.
Classical Theory}, Cambridge Studies in Advanced Mathematics \textbf{97},
Cambridge University Press, 2007.

\bibitem{Nistor}
V.~Nistor, \emph{Group cohomology and the cyclic cohomology of crossed
products}, Invent. Math. \textbf{99} (1990), 411--424.

\bibitem{Schaefer}
H.~H.~Schaefer, \emph{Topological Vector Spaces}, Graduate Texts in Math.
\textbf{3}, Springer, 1971.

\bibitem{Tate}
J.~Tate, \emph{Fourier analysis in number fields and Hecke's
zeta-functions}, thesis, Princeton, 1950; reprinted in: J.~W.~S.~Cassels and
A.~Fr\"ohlich (eds.), \emph{Algebraic Number Theory}, Academic Press, 1967,
305--347.

\bibitem{Treves}
F.~Tr\`eves, \emph{Topological Vector Spaces, Distributions and Kernels},
Academic Press, 1967.

\bibitem{Weil1948}
A.~Weil, \emph{Sur les courbes alg\'ebriques et les vari\'et\'es qui s'en
d\'eduisent}, Hermann, Paris, 1948.

\bibitem{Weil1952}
A.~Weil, \emph{Sur les ``formules explicites'' de la th\'eorie des nombres
premiers}, Comm. S\'em. Math. Univ. Lund (1952), Tome suppl\'ementaire,
252--265.

\bibitem{Wodzicki}
M.~Wodzicki, \emph{Excision in cyclic homology and in rational algebraic
$K$-theory}, Ann. of Math. (2) \textbf{129} (1989), 591--639.

\end{thebibliography}

\end{document}
